
\chapter{Free field archetypes}\label{chap:free-fields}
% Regime I: quadratic (Convention~\ref{conv:regime-tags}).

\noindent
\emph{Type signature for this chapter:} every free-field algebra
$\mathcal A\in\{\mathcal F, \cH_k, \beta\gamma_\lambda, bc_\lambda,
V_\Lambda\}$ is a chart algebra
$A_b=\mathrm{End}_{\mathcal{C}^{\mathrm{op}}_X}(b)$ on the
factorisation dg-cat $\mathcal{C}^{\mathrm{op}}$ on $(X, D, \tau)$
at the standard vacuum $b$. Theorem~A
(Theorem~\ref{thm:A-infinity-2} in
Chapter~\ref{chap:theorem-A-infinity-2}) is the $1\!\leftrightarrow\!2$
reconstruction at the chart-to-twisting-coalgebra step;
chart-class enumeration is at level~$1$, the $5\times 5$
$\kappa$-stratification is at level~$5$. The free-field locus
populates rows $\mathsf{G}$ (free fermion, Heisenberg, lattice)
and $\mathsf{C}$ ($\beta\gamma$, $bc$) of the five-archetype matrix,
with shadow depths $r_{\max}\in\{2,4\}$.

The Heisenberg algebra has one generator. Its shadow tower
terminates at degree~$2$, its trace-form collision residue is
\[
r_{\cH_k}^{\mathrm{coll}}(z)=\frac{k}{z},
\]
and every higher bracket on the convolution algebra is zero: the
archetype is Gaussian (class~$\mathsf{G}$).
This is the canonical $\mathbf{G}$-class anchor of the
quadrichotomy
$\{\mathbf{G},\mathbf{L},\mathbf{C},\mathbf{M}\}$ of
Chapter~\ref{chap:shadow-quadrichotomy-platonic}; the
$\mathbf{C}$-class anchor is $\beta\gamma$
(Chapter~\ref{chap:beta-gamma}), reached by adjoining a second
boson with simple-pole symplectic OPE, while the dual $bc$-ghost
pairing realises the universal Koszul conductor
$K_{\beta\gamma} = c_{\beta\gamma} + c_{bc} = 0$ of
Chapter~\ref{chap:universal-conductor} in opposite statistical
sectors. What changes when the Gaussian one-generator input is
replaced by a first-order pair?

The answer separates statistics from arity. Replacing the Heisenberg
current by a single fermionic generator with antisymmetric OPE
$\psi(z)\psi(w) \sim 1/(z{-}w)$ leaves the tower Gaussian: the
antisymmetry kills the would-be contact shadow before it can
propagate, so no shadow beyond degree~$2$ survives. A second boson
coupled through a symplectic pairing
produces the $\beta\gamma$ system (Chapter~\ref{chap:beta-gamma});
here the simple-pole OPE $\beta(z)\gamma(w) \sim 1/(z{-}w)$ has zero
pole-valued collision residue after $d\log$ absorption, but it remains
as a regular ordered contact operator. The standard conformal-weight
family then carries the quartic contact coefficient
$S_4=-5/12$, and the shadow tower extends to degree~$4$
(class~$\mathsf{C}$). Reversing the
statistics of both generators while keeping the same first-order
residue gives the $bc$ ghost system, Koszul dual to
$\beta\gamma$ with matching shadow depth and opposite
$\kappa$-coefficient. The dichotomy is sharp: one generator forces
class~$\mathsf{G}$; two first-order generators coupled by a
nondegenerate simple-pole pairing force class~$\mathsf{C}$, with
Weyl and Clifford statistics giving the $\beta\gamma$ and $bc$
realizations. No finer gradation exists within the free-field regime.
The binary pole data are not uniform: Heisenberg carries the scalar
residue $k/z$, while $\beta\gamma$ and $bc$ carry regular ordered
contact data, Weyl for $\beta\gamma$ and graded Clifford for $bc$.
These contact tensors are not pole-valued collision residues
(Theorem~\ref{thm:bg-ordered-bar}). The symmetric scalar
rational $r$-matrix of the one-generator free fermion is zero.
The class-$\mathsf{C}$ distinction is the degree~$4$ contact
stratum together with these regular ordered contact data.

The modular Koszul triples
(Definition~\ref{def:modular-koszul-triple}) of the free-field
archetypes are:
\begin{equation}\label{eq:fermion-triple}
\mathfrak{T}_{\cF}
\;=\;
\bigl(\,\cF,\;\;
\operatorname{Sym}^{\mathrm{ch}}(\gamma),\;\;
r(z) = 0\,\bigr),
\end{equation}
\begin{equation}\label{eq:betagamma-triple}
\mathfrak{T}_{\beta\gamma}
\;=\;
\bigl(\,\beta\gamma_\lambda,\;\;
bc_\lambda,\;\;
r_{\beta\gamma}^{\mathrm{coll}}(z)=0\,\bigr),
\qquad
\mathfrak{T}_{bc}
\;=\;
\bigl(\,bc_\lambda,\;\;
\beta\gamma_\lambda,\;\;
r_{bc}^{\mathrm{coll}}(z)=0\,\bigr),
\end{equation}
The free fermion $r$-matrix vanishes: the simple pole in the
OPE $\psi(z)\psi(w) \sim 1/(z{-}w)$ is absorbed by $d\log$
extraction (Remark~\ref{rem:fermion-r-matrix-vanishing}).
The $\beta\gamma$--$bc$ pair is different: its simple pole gives no
pole-valued scalar collision residue, but it survives on the
ordered bar side as the regular contact kernel
\[
\Theta_{\beta\gamma}^{\mathrm{ord}}
=\beta\otimes\gamma-\gamma\otimes\beta,
\qquad
\Theta_{bc}^{\mathrm{ord}}
=b\otimes c+c\otimes b
\quad\text{(graded Clifford sign).}
\]
Thus $R=\id+2\pi i\,\Theta$ before symmetric averaging, while
the scalar shadow records the class-$\mathsf{C}$ quartic contact
term.
The stress tensor inside $\beta\gamma_\lambda$ or $bc_\lambda$ is an
autonomous Virasoro subalgebra and has its own one-dimensional
class-$\mathsf{M}$ shadow tower. This ambient stress-tensor lane is not
the free-field archetype: the standard two-generator family is the
finite-depth class-$\mathsf{C}$ object.

\begin{remark}[Shadow obstruction tower for free fields: the Gaussian base case]
\label{rem:free-fields-master-mc}
\index{shadow obstruction tower!free-field archetypes}
Each free-field archetype specializes the universal MC element $\Theta_\cA \in \MC(\gAmod)$
(Theorem~\ref{thm:mc2-bar-intrinsic}). The Gaussian archetype is
characterized by maximal truncation: the shadow
obstruction tower terminates at degree~$2$, so that
$\Theta_\cA^{\leq 2} = \kappa \cdot \eta \otimes \Lambda$ and all
higher projections $\Theta_\cA^{\leq r} = \Theta_\cA^{\leq 2}$ for
$r \geq 3$. Equivalently, the convolution algebra
$\gAmod$ has trivial higher brackets: the modular deformation complex
$\Definfmod(\cA)$ is formal.

The free fermion is Gaussian with
$\kappa(\mathcal F)=c/2=1/4\neq0$ on the Virasoro scalar lane. The
Heisenberg Gaussian has instead
$\kappa(\mathcal H_k)=k$, the level coefficient in
$r_{\cH_k}^{\mathrm{coll}}(z)=k/z$, not $c(\mathcal H_k)/2$.
The $\beta\gamma$ and $bc$ systems break out of the
Gaussian class because the standard conformal-weight family
$\beta\gamma_\lambda$ satisfies
$S_2 = 6\lambda^2 - 6\lambda + 1$, $S_3 = 0$, $S_4 = -5/12$,
and $S_r = 0$ for $r \ge 5$, while its $bc_\lambda$ Verdier-dual
has the opposite scalar coefficient
$\kappa(bc_\lambda)=-(6\lambda^2-6\lambda+1)$ and the same
finite class-C depth. This remains true even though the internal
weight-changing line is shadow-trivial:
$\Theta_{\beta\gamma}^{\le r}\big|_{\mathrm{w.c.}} = 0$ for
$r \ge 2$. Both two-generator free-field systems
($\beta\gamma$ and $bc$) are class~C ($r_{\max} = 4$); the
single-generator free fermion is class~G ($r_{\max} = 2$) by
fermionic antisymmetry. The free-field periodic table therefore
exhibits the transition from the Gaussian to the contact/quartic
archetype within the shadow depth classification
(Table~\ref{tab:shadow-tower-census}).
\end{remark}

\begin{remark}[Three-pillar interpretation: free-field archetypes]
\label{rem:free-fields-three-pillar}
\index{free fields!three-pillar interpretation}
In the three-pillar architecture
(\S\ref{sec:concordance-three-pillars}):
(i)~free bosons produce a \emph{strict} homotopy chiral algebra
(Gaussian archetype, all secondary Borcherds operations
$j'_n = 0$ for $n \geq 3$;
Definition~\ref{def:secondary-borcherds},
cf.\ Example~\ref{ex:cech-hca-heisenberg}). The $\beta\gamma$--$bc$
pairs are non-strict on the bar side: the simple-pole contact
operation produces the first non-linear transfer, and the standard
conformal-weight family has the class-C quartic shadow
$S_4=-5/12$. The auxiliary invariant
$\mu_{\beta\gamma}$ vanishes
(Corollary~\ref{cor:nms-betagamma-mu-vanishing}), but this is not the
vanishing of the class-C quartic shadow. The shadow tower terminates
at degree~$4$;
(ii)~the convolution $sL_\infty$-algebra
(Theorem~\ref{thm:operadic-homotopy-convolution}) is \emph{formal}
for all free-field archetypes: the transferred higher brackets
$\ell_k^{\mathrm{tr}} = 0$ for $k \geq 3$ on free bosons and
fermions, and $\ell_4^{\mathrm{tr}}$ is a coboundary for
$\beta\gamma$--$bc$. The shadow algebra is
$\mathbb{Q}[\kappa]$ (bosonic) or
$\mathbb{Q}[\kappa,S_4,\mu]/(\mu)$ ($\beta\gamma$--$bc$), with
$S_4=-5/12$ on the standard conformal-weight family;
(iii)~the log-FM tropicalization
(Theorem~\ref{thm:planted-forest-tropicalization}) reduces to the
classical FM compactification for free fields: no nodal divisor
insertions are needed because the only free-field singularities are the
Heisenberg double pole and the simple first-order pairings, all handled
by ordinary FM collision residues or regular contact transport. Thus the
planted-forest coefficient algebra
$\mathcal{G}_{\mathrm{pf}}$ contributes no corrections beyond
depth~$1$.
Free-field archetypes are the unique family where Pillar~B is formal
and Pillar~C is classical simultaneously.
\end{remark}

\begin{remark}[Chart-class enumeration: free fermion]
\label{rem:fermion-chart-class-enumeration}
\index{chart-class!free fermion}
\emph{Type signature:} (Open quadrant; chart presentation;
Beilinson level~$1$; HP = factorisation dg-cat
$\mathcal{C}^{\mathrm{op}}$ on $(X, D, \tau)$, vacuum $b$,
$A_b = \mathrm{End}_\mathcal{C}(b)$; archetype label invariant
under Morita equivalence on the level-$1\to 0$ fibre).
The chart algebra
\[
A_b^{\mathcal F} =
\mathrm{End}_{\mathcal{C}^{\mathrm{op}}_X}(b^{\mathcal F})
\;\simeq\; \mathcal F\otimes_{\mathcal{D}_X}(\,\cdot\,)
\]
on the standard NS-vacuum $b^{\mathcal F}$ realises the
single-generator real free fermion. The five-archetype label is
$\mathsf{G}$ (Gaussian, $r_{\max}=2$), not because the OPE is
quadratic (it is in fact a single-pole product) but because
fermionic antisymmetry kills every cubic and higher shadow.
Two Morita-distinct chart presentations exist:
\begin{enumerate}[label=\textup{(G\arabic*)}]
\item \emph{Real (Majorana) chart}: one fermionic generator
$\psi(z)$ with OPE $\psi(z)\psi(w)\sim 1/(z{-}w)$. The single-spin
Pfaffian sector picks $2^{2g}$ spin structures per genus.
\item \emph{Bosonised lattice chart}: under the boson-fermion
correspondence $\mathcal F\otimes\mathcal F \cong
V_{\mathbb Z}=\cH_1$, two real fermions assemble into the
rank-one Heisenberg on the integer lattice. This is a Morita
equivalence between two chart presentations of the same
class-$\mathsf{G}$ archetype: the rank-two Pfaffian sector and the
rank-one Heisenberg determinant sector.
\end{enumerate}
The Koszul-dual chart $\mathcal F^! = \mathrm{Sym}^{\mathrm{ch}}(\gamma)$
exchanges Bose and Fermi statistics on a single generator
(Theorem~\ref{thm:single-fermion-boson-duality}); its chart class is
also $\mathsf{G}$ but on a distinct Morita orbit (no spin-structure
data because $\gamma$ is bosonic). The class-$\mathsf{G}$ row of
the five-archetype matrix therefore admits at least two Morita
orbits in the free-field locus: $(\mathcal F, \mathrm{Sym}^{\mathrm{ch}}(\gamma))$
and the Heisenberg pair $(\cH_k, \cH_k^!)$.
\end{remark}

\begin{remark}[Theorem A in four lanes: free fermion]
\label{rem:fermion-theorem-A-four-lanes}
\index{Theorem A!four lanes!free fermion}
On the chart algebra $A_b^{\mathcal F}=\mathcal F$, Theorem~A
(Theorem~\ref{thm:A-infinity-2} in
Chapter~\ref{chap:theorem-A-infinity-2}) appears in the four lanes
$\mathrm{L}1,\mathrm{L}2,\mathrm{L}3,\mathrm{L}4$:
\begin{itemize}
\item \emph{$\mathrm{L}1$ (chain level, $\mathrm{Ch}(\mathrm{Vect})$).}
The bar--cobar adjunction is realised by the maximally collapsed
genus-0 bar complex $\bar{B}^0_{\mathrm{geom}}(\mathcal F)=\mathbb C$,
$\bar{B}^n_{\mathrm{geom}}(\mathcal F)=0$ for $n\ge 1$
(Theorem~\ref{thm:fermion-bar-complex-genus-0}); the Gaussian
witness chain homotopy $h_G=0$.
\item \emph{$\mathrm{L}2$ (Quillen-equivalence).}
The Vallette bar--cobar Quillen equivalence applies trivially in
the Francis--Gaitsgory ambient because the cofree coalgebra and the
free algebra coincide on a single fermionic generator.
\item \emph{$\mathrm{L}3$ ($(\infty,1)$-categorical).}
The $(\infty,1)$-localisation gives the adjoint equivalence between
the Pfaffian sector $\mathrm{Pf}(\bar\partial_{K^{1/2}})$ and the
symmetric chiral algebra $\mathrm{Sym}^{\mathrm{ch}}(\gamma)$ on
the Verdier-dual side.
\item \emph{$\mathrm{L}4$ ($(\infty,2)$-properad).}
Theorem~$A^{\infty,2}$ inscribes the $(\infty,2)$-categorical adjoint
equivalence on the properadic envelope; the Hackney--Robertson lift
is strict because all higher properadic operations vanish on the
single-generator Gaussian sector.
\end{itemize}
The four lanes coincide on the entire chart locus of $\mathcal F$.
\end{remark}

\begin{remark}[Five-$\kappa$ stratification: free fermion]
\label{rem:fermion-five-kappa-stratification}
\index{kappa stratification@$\kappa$-stratification!free fermion row G}
The five $\kappa$-measurements
$\{\kappa_{\mathrm{cat}}, \kappa^{\mathrm{Hodge}}_{\mathrm{ch}},
\kappa^{\mathrm{Heis}}_{\mathrm{ch}}, \kappa_{\mathrm{BKM}},
\kappa_{\mathrm{fiber}}\}$ on the row $\mathsf{G}$ of the
$5\times 5$ stratification matrix at the free fermion $\mathcal F$
evaluate as
\[
\kappa_{\mathrm{cat}}^{\mathcal F}=\frac14,\qquad
\kappa^{\mathrm{Hodge}}_{\mathrm{ch}}(\mathcal F)=\frac12\cdot\frac14
=\frac18,\qquad
\kappa^{\mathrm{Heis}}_{\mathrm{ch}}(\mathcal F)=0,
\]
$\kappa_{\mathrm{BKM}}=0$, $\kappa_{\mathrm{fiber}}=0$.
The Heisenberg projection vanishes because no abelian current sits
inside the single fermionic generator at the chart level (the
$\mathcal F\otimes\mathcal F\to V_{\mathbb Z}$ bosonisation is a
two-fermion derived statement, not a chart-level invariant).
\emph{Type signature:} (Open quadrant; chart-level scalar tuple;
Beilinson level~$5$; HP = five distinct measurements named;
collapse pattern $r_{\max}=2$). five-invariant discipline: the
$\kappa$-tuple distinguishes the free fermion class-$\mathsf{G}$
chart from the Heisenberg class-$\mathsf{G}$ chart even though
both terminate at $r_{\max}=2$, because the
$\kappa^{\mathrm{Heis}}_{\mathrm{ch}}$ entry vanishes here while it
equals $k$ on $\cH_k$.
\end{remark}

\begin{remark}[Free-field bar complex: the growth baseline]
\label{rem:free-field-bar-baseline}
\index{bar complex!free-field baseline}
The bar cohomology of the free-field archetypes fixes the first
growth baseline for subsequent families.
At genus~$0$, $\dim H^n(\barBgeom(\cA))$ grows as follows:
\begin{center}
\begin{tabular}{l c c c c c c c l}
\textbf{Family} & $n{=}0$ & $1$ & $2$ & $3$ & $4$ & $5$ & $6$
 & \textbf{GF} \\
\hline
Free fermion $\mathcal{F}$ & $1$ & $0$ & $0$ & $0$ & $0$ & $0$ & $0$
 & $1$ \\
Heisenberg $\mathcal{H}_k$ & $1$ & $0$ & $1$ & $0$ & $0$ & $0$ & $0$
 & $1 + x^2$ \\
$\beta\gamma$ & $1$ & $2$ & $4$ & $10$ & $26$ & $70$ & $192$
 & $\sqrt{(1{+}x)/(1{-}3x)}$ \\
$bc$ ghosts & $1$ & $2$ & $3$ & $6$ & $13$ & $28$ & $59$
 & $1+\sum_{n\geq1}(2^n-n+1)x^n$
\end{tabular}
\end{center}
The free fermion and Heisenberg genus-$0$ bar cohomologies are
finite-dimensional. The $\beta\gamma$ row is the algebraic
degree-$2$ branch with discriminant
$\Delta(x)=(1{-}3x)(1{+}x)$
(Theorem~\ref{thm:ds-bar-gf-discriminant}); the $bc$ row is the
rational Koszul-dual branch
$\dim H^n(\barBgeom(bc))=2^n-n+1$ for $n\ge1$
(Proposition~\ref{prop:bar-bc-system},
Corollary~\ref{cor:bar-cohomology-koszul-dual}). Thus the two
two-generator simple-pole systems share class-$\mathsf C$ shadow
depth but not bar-cohomology growth: $\beta\gamma$ has dominant
rate~$3^n$ with an $n^{-1/2}$ correction, while $bc$ has dominant
rate~$2^n$.
\end{remark}

\section{Free-field archetypes: setup and OPE structure}\label{sec:free-field-atoms-setup}

The free-field archetypes are the one-generator Gaussian systems
(the free fermion and its one-generator bosonic Koszul dual), the
two-generator first-order system $\beta\gamma$, and the two-generator
ghost system $bc$. Each definition specifies input data for the universal construction
$\Theta_\cA := D_\cA - d_0 \in \MC(\gAmod)$
(Theorem~\ref{thm:mc2-bar-intrinsic}). The bar complex
$B(\cA)=T^c(s^{-1}\bar\cA)$, its cohomology coalgebra
$\cA^{\mathrm{i}}=H^\bullet B(\cA)$, the Verdier-linear dual
algebra $\cA^!=(\cA^{\mathrm{i}})^\vee$, the genus tower, and
the shadow invariants are distinct outputs of this datum. The
cobar equivalence $\Omega B(\cA)\simeq \cA$ is inversion; it is
not the definition of $\cA^!$. For free-field archetypes, the
projection is maximally simple: the Gaussian archetype (free
fermion) has $\Theta_\cA = \Theta_\cA^{\leq 2}$, and the
contact/quartic archetype ($\beta\gamma$, $bc$) terminates at
$\Theta_\cA^{\leq 4}$.

The numerical invariants used below are not interchangeable.
The generator-pole order $p_{\max}$ records the largest pole in the
chosen generators; the collision residue depth $k_{\max}$ records the
pole-valued part that survives ordered $d\log$ absorption; the shadow
depth $r_{\max}$ records the highest nonzero scalar obstruction in the
ordered-to-symmetric shadow tower.
\begin{center}
\scriptsize
\renewcommand{\arraystretch}{1.25}
\begin{tabular}{lccccccl}
\toprule
\emph{System} & $p_{\max}$ & $k_{\max}$ & $r_{\max}$ &
$\kappa$ & $c$ & \emph{Verdier partner} & \emph{Chiral Hochschild} \\
\midrule
Free fermion $\mathcal F$ & $1$ & $0$ & $2$ &
$1/4$ & $1/2$ &
$\mathrm{Sym}^{\mathrm{ch}}(\gamma)$ &
$\ChirHoch^{>2}=0$ \\
Heisenberg $\cH_k$ & $2$ & $1$ & $2$ &
$k$ & $1$ &
$\mathrm{Sym}^{\mathrm{ch}}(V^*)$, $m_0=-k\omega_{\mathrm{cen}}$ &
$1+t+t^2$ \\
$\beta\gamma_\lambda$ & $1$ & $0$ & $4$ &
$6\lambda^2-6\lambda+1$ & $2(6\lambda^2-6\lambda+1)$ &
$bc_\lambda$ &
$\ChirHoch^{>2}=0$ \\
$bc_\lambda$ & $1$ & $0$ & $4$ &
$-(6\lambda^2-6\lambda+1)$ & $-2(6\lambda^2-6\lambda+1)$ &
$\beta\gamma_\lambda$ &
$\ChirHoch^{>2}=0$ \\
Symplectic boson $\beta\gamma_{1/2}$ & $1$ & $0$ & $4$ &
$-1/2$ & $-1$ &
$bc_{1/2}$ &
$\ChirHoch^{>2}=0$ \\
Symplectic fermion $\mathcal{SF}=bc_1$ & $1$ & $0$ & $4$ &
$-1$ & $-2$ &
$\beta\gamma_1$ &
$\ChirHoch^{>2}=0$ \\
Lattice $V_\Lambda$ & $2$ & $1$ & $2$ &
$\operatorname{rank}\Lambda$ & $\operatorname{rank}\Lambda$ &
$V_{\Lambda^\vee}$ with opposite scalar curvature &
$\ChirHoch^{>2}=0$ on Koszul finite-type sectors \\
\bottomrule
\end{tabular}
\end{center}
The corresponding singular local data are:
\begin{center}
\scriptsize
\renewcommand{\arraystretch}{1.25}
\begin{tabular}{llll}
\toprule
\emph{System} & \emph{Singular OPE} & \emph{Pole-valued residue}
& \emph{Ordered or lattice datum} \\
\midrule
$\mathcal F$ &
$\psi(z)\psi(w)\sim (z-w)^{-1}$ &
$0$ &
scalar contraction; no field-valued $r$ \\
$\cH_k$ &
$J(z)J(w)\sim k(z-w)^{-2}$ &
$k/z$ &
central oscillator curvature $m_0^!=-k\omega_{\mathrm{cen}}$ \\
$\beta\gamma_\lambda$ &
$\beta(z)\gamma(w)\sim (z-w)^{-1}$,\;
$\gamma(z)\beta(w)\sim -(z-w)^{-1}$ &
$0$ &
$\Theta_{\beta\gamma}^{\mathrm{ord}}
=\beta\otimes\gamma-\gamma\otimes\beta$ \\
$bc_\lambda$ &
$b(z)c(w)\sim (z-w)^{-1}$,\;
$c(z)b(w)\sim (z-w)^{-1}$ &
$0$ &
$\Theta_{bc}^{\mathrm{ord}}=b\otimes c+c\otimes b$ \\
$\beta\gamma_{1/2}$ &
same first-order OPE as $\beta\gamma_\lambda$ &
$0$ &
symplectic boson; $\kappa=-1/2$ \\
$\mathcal{SF}=bc_1$ &
same first-order OPE as $bc_\lambda$ &
$0$ &
symplectic fermion; $\kappa=-1$ \\
$V_\Lambda$ &
$J_a(z)J_b(w)\sim (a,b)(z-w)^{-2}$;\;
$e^\alpha e^\beta\sim \epsilon_{\alpha,\beta}
(z-w)^{(\alpha,\beta)}e^{\alpha+\beta}$ &
$(a,b)/z$ on oscillators &
simple lattice residues when $(\alpha,\beta)=-1$; scalar class~G \\
\bottomrule
\end{tabular}
\end{center}
For the first-order rows, the class-C quartic contact is not a
generator-pole assertion: $\beta(z)\gamma(w)$ and $b(z)c(w)$ have only
simple poles, while $S_4=-5/12$ is transferred from the regular
ordered contact class. For the lattice row, $p_{\max}=2$ and
$k_{\max}=1$ refer to the oscillator currents. Root vertex operators
may create simple lattice-graded residue maps, but they do not produce
a class-L scalar shadow; the universal scalar shadow remains Gaussian.

\subsection{Free fermion}\label{sec:free-fermion}%
\index{free fermion|textbf}%
\index{free fields!as archetypes}

\begin{remark}[Shadow archetype: Gaussian]
\label{rem:fermion-shadow-archetype}
\index{free fermion!shadow archetype}
Shadow depth $r_{\max} = 2$ (class~G).
Modular characteristic $\kappa(\mathcal{F}) = c/2 = 1/4$.
Central charge $c = 1/2$.
The shadow obstruction tower is purely scalar:
$\Theta_{\mathcal{F}}^{\leq r}
= \Theta_{\mathcal{F}}^{\leq 2}
= \kappa \cdot \eta \otimes \Lambda$
for all $r \geq 2$. No cubic or higher shadows.
Koszul dual: $\mathcal{F}^! = \mathrm{Sym}^{\mathrm{ch}}(\gamma)$
(one bosonic generator, Theorem~\ref{thm:single-fermion-boson-duality}).
\end{remark}

The free fermion replaces the Heisenberg double pole with a fermionic simple pole; the bar complex collapses.

\begin{definition}[Free fermion chiral algebra]
The free fermion chiral algebra $\mathcal{F}$ is generated by a single fermionic field $\psi(z)$ of
conformal weight $h = \frac{1}{2}$ with OPE:
\[
\psi(z)\psi(w) = \frac{1}{z-w} + \text{regular}
\]
The quadratic relation enforcing fermionic statistics is:
\[
R_{\text{ferm}} = \{\psi(z_1) \otimes \psi(z_2) + \psi(z_2) \otimes \psi(z_1)\} \subset
j_*j^*(\mathcal{F} \boxtimes \mathcal{F})
\]
\end{definition}
\index{conformal weight}

\begin{remark}[Fermionic sign]
The antisymmetry $\psi(z)\psi(w) = -\psi(w)\psi(z)$ away from the diagonal forces many components of the bar complex to vanish.
\end{remark}

\subsubsection{Shadow obstruction tower of the free fermion}
\label{sec:fermion-shadow-tower}%
\index{free fermion!shadow obstruction tower}%

The free fermion is the Gaussian archetype in its purest form. The
Heisenberg algebra achieves class~G through \emph{commutativity}: the
double pole generates only symmetric bar elements, and the absence of
a simple pole kills the cubic shadow. The free fermion achieves class~G
through \emph{antisymmetry}: the simple pole would produce a cubic
shadow in a bosonic algebra, but fermionic statistics force $S_3 = 0$
because cyclic permutations on three-point configuration space forms
act by cube roots of unity (proof:
Theorem~\ref{thm:fermion-bar-complex-genus-0}, degree~$2$). The two
mechanisms are complementary, and both produce the same outcome: the
shadow obstruction tower terminates at degree~$2$.

\begin{proposition}[Shadow invariants of the free fermion;
\ClaimStatusConditional]
\label{prop:fermion-shadow-invariants}
\index{free fermion!shadow invariants}
The shadow obstruction tower of the free fermion
$\mathcal{F}$ is given by:
\begin{enumerate}
\item\label{item:fermion-kappa}
 Modular characteristic:
 $\kappa(\mathcal{F}) = c/2 = 1/4$.
\item\label{item:fermion-S3}
 Cubic shadow:
 $S_3(\mathcal{F}) = 0$.
 \textup{(}Fermionic antisymmetry forces the cubic coupling to vanish.\textup{)}
\item\label{item:fermion-S4}
 Quartic shadow:
 $S_4(\mathcal{F}) = 0$.
 \textup{(}Consequence of $S_3 = 0$: no quartic contact term survives in the
 deformation complex.\textup{)}
\item\label{item:fermion-discriminant}
 Critical discriminant:
 $\Delta(\mathcal{F}) = 8\kappa S_4 = 0$.
\item\label{item:fermion-class}
 Shadow class: G \textup{(}Gaussian\textup{)}, $r_{\max} = 2$.
\item\label{item:fermion-tower}
 Tower termination:
 $\Theta_{\mathcal{F}}^{\leq r}
 = \Theta_{\mathcal{F}}^{\leq 2}
 = \tfrac{1}{4} \cdot \eta \otimes \Lambda$
 for all $r \geq 2$.
\end{enumerate}
\end{proposition}

\begin{proof}
\eqref{item:fermion-kappa}:
The free fermion has central charge $c = 1/2$ (a single real
fermion of weight $h = 1/2$; standard from the Virasoro commutation
relations).
Its modular characteristic is the stress-tensor scalar coefficient:
$\kappa(\mathcal F)=\kappa(\mathrm{Vir}_{1/2})=c/2=1/4$ by
Theorem~\ref{thm:modular-characteristic}. This is the Virasoro scalar
lane. It is not the Heisenberg convention, where
$\kappa(\mathcal H_k)=k$ is the level of the double-pole OPE rather
than $c(\mathcal H_k)/2$.

\eqref{item:fermion-S3}:
The cubic shadow $S_3$ is the degree-$3$ projection of
$\Theta_{\mathcal{F}}$ in the modular convolution algebra. A degree-$3$
element requires three fermion insertions on a three-point configuration
space. By fermionic antisymmetry, $\psi_1 \otimes \psi_2 \otimes \psi_3$
is totally antisymmetric in labels. The cyclic permutation
$(123) \mapsto (231)$ on $\overline{C}_3(X)$ acts on the logarithmic
$2$-form space $\Omega^2_{\log}(\overline{C}_3)$ with eigenvalues
$\zeta_3, \zeta_3^2$ (primitive cube roots of unity). Compatibility
with the antisymmetric tensor factor forces the form to lie in the
$\zeta_3$-eigenspace, but the bar differential requires the
$1$-eigenspace: intersection is zero.

\eqref{item:fermion-S4}:
With $S_3 = 0$, the quartic contact invariant
$Q^{\mathrm{contact}} = 0$ by the shadow obstruction tower
(Theorem~\ref{thm:recursive-existence}):
the obstruction to extending $\Theta^{\leq 3}$ to $\Theta^{\leq 4}$
lies in the cyclic deformation complex at degree~$4$, and with
$S_3 = 0$ the obstruction class $o_4 = 0$
by cubic gauge triviality
(Theorem~\ref{thm:cubic-gauge-triviality}).

\eqref{item:fermion-discriminant}--\eqref{item:fermion-class}:
$\Delta = 8 \cdot (1/4) \cdot 0 = 0$, so $Q_{\mathcal{F}}(t)$ is
a perfect square. By the single-line dichotomy
(Theorem~\ref{thm:single-line-dichotomy}), $\Delta = 0$
implies $r_{\max} \leq 3$. With $S_3 = 0$, the tower terminates
at exactly $r_{\max} = 2$: class~G.

\eqref{item:fermion-tower}:
Since $S_r = 0$ for all $r \geq 3$,
$\Theta_{\mathcal{F}}^{\leq r} = \kappa \cdot \eta \otimes \Lambda$
for all $r \geq 2$, and the inverse limit is trivially
$\Theta_{\mathcal{F}} = (1/4) \cdot \eta \otimes \Lambda$.
\end{proof}

\begin{proposition}[Shadow metric of the free fermion;
\ClaimStatusProvedHere]
\label{prop:fermion-shadow-metric}
\index{free fermion!shadow metric}
The shadow metric \textup{(}Definition~\textup{\ref{def:shadow-metric}}\textup{)}
evaluates to:
\begin{equation}\label{eq:fermion-shadow-metric}
Q_{\mathcal{F}}(t) = (2\kappa(\mathcal F))^2 = \tfrac{1}{4},
\end{equation}
a constant independent of $t$. The weighted shadow generating function
$H(t) = t^2 \sqrt{Q_{\mathcal{F}}} = \tfrac{1}{2}\,t^2$ is purely
quadratic; only $S_2 = \kappa(\mathcal F)=1/4$ contributes. The shadow
connection $\nabla^{\mathrm{sh}} = d - Q'/(2Q)\,dt$ is trivial
\textup{(}$\omega = d\log Q = 0$\textup{)}, and parallel transport
carries no monodromy.

Compare: the Heisenberg has
$Q_{\mathcal{H}}(t) = 4k^2$ \textup{(}also constant\textup{)}. Both
are Gaussian archetypes. The structural difference is that
$\kappa(\mathcal{H}_k) = k$ is tunable
\textup{(}$k$ is a continuous parameter\textup{)}, while
$\kappa(\mathcal{F}) = 1/4$ is fixed.
\end{proposition}

\begin{proof}
The shadow metric is $Q_L(t) = (2\kappa + 3\alpha t)^2 + 2\Delta\,t^2$
where $\alpha = S_3$ and $\Delta = 8\kappa S_4$. With
$S_3 = 0$ and $S_4 = 0$:
$Q_{\mathcal{F}}(t) = (2 \cdot \tfrac{1}{4})^2 = (1/2)^2 = 1/4$.
\end{proof}

\subsubsection{Genus expansion of the free fermion}
\label{sec:fermion-genus-expansion}%
\index{free fermion!genus expansion}%

\begin{theorem}[Free fermion genus expansion; \ClaimStatusConditional]
\textup{(}UNIFORM-WEIGHT\textup{})
\label{thm:fermion-genus-expansion}
For all $g \geq 1$, the genus-$g$ scalar/full free energy of the free
fermion is:
\begin{equation}\label{eq:fermion-Fg}
F_g(\mathcal{F})
=F_g^{\mathrm{sc}}(\mathcal{F})
= \frac{1}{4} \cdot \frac{2^{2g-1}-1}{2^{2g-1}}
 \cdot \frac{|B_{2g}|}{(2g)!}
\end{equation}
where $B_{2g}$ denotes the $2g$-th Bernoulli number, and
$\delta F_g^{\mathrm{cross}}(\mathcal F)=0$ on this class~G lane.
The first values are:
\begin{align}
F_1(\mathcal{F}) &= \frac{1}{96},
 \label{eq:fermion-F1} \\
F_2(\mathcal{F}) &= \frac{7}{23040},
 \label{eq:fermion-F2} \\
F_3(\mathcal{F}) &= \frac{31}{3870720},
 \label{eq:fermion-F3} \\
F_4(\mathcal{F}) &= \frac{127}{619315200}.
 \label{eq:fermion-F4}
\end{align}
The generating function is:
\begin{equation}\label{eq:fermion-generating-function}
\sum_{g=1}^{\infty} F_g(\mathcal{F})\,x^{2g}
=\sum_{g=1}^{\infty} F_g^{\mathrm{sc}}(\mathcal{F})\,x^{2g}
= \frac{1}{4}\left(\frac{x/2}{\sin(x/2)} - 1\right).
\end{equation}
\end{theorem}

\begin{proof}
The free fermion is a modular Koszul algebra
with $\kappa(\mathcal{F}) = 1/4$
\textup{(}Proposition~\textup{\ref{prop:fermion-shadow-invariants}}\textup{)},
uniform conformal weight $h = 1/2$, and a single generator. By
Theorem~\ref{thm:modular-characteristic},
$\mathrm{obs}_g(\mathcal{F})
 = \kappa(\mathcal{F}) \cdot \lambda_g$ for all
$g \geq 1$. The genus-$g$ free energy is
$F_g(\mathcal{F})=F_g^{\mathrm{sc}}(\mathcal{F})
 = \kappa(\mathcal{F}) \cdot \lambda_g^{\mathrm{FP}}$ where
\[
\lambda_g^{\mathrm{FP}}
=\int_{\overline{\mathcal M}_{g,1}}\psi_1^{2g-2}\lambda_g
= \frac{2^{2g-1}-1}{2^{2g-1}} \cdot \frac{|B_{2g}|}{(2g)!}
\]
is the marked Faber--Pandharipande scalar coefficient. This gives
$F_g(\mathcal{F})=F_g^{\mathrm{sc}}(\mathcal{F})
= (1/4) \cdot \lambda_g^{\mathrm{FP}}$, yielding the
stated values via $B_2 = 1/6$, $B_4 = -1/30$, $B_6 = 1/42$,
$B_8 = -1/30$.
The generating function is the $\hat{A}$-genus at imaginary argument
(equation~\eqref{eq:frame-generating-function} with $k$ replaced
by $\kappa(\mathcal F)=1/4$).

The explicit checks:
$F_1(\mathcal F)=F_1^{\mathrm{sc}}(\mathcal F)
=\kappa(\mathcal F)(1/24)=1/96$,
$F_2(\mathcal F)=F_2^{\mathrm{sc}}(\mathcal F)
=\kappa(\mathcal F)(7/5760)=7/23040$,
$F_3(\mathcal F)=F_3^{\mathrm{sc}}(\mathcal F)
=(1/4)\cdot31/32\cdot(1/42)/(720)=31/3870720$,
and
$F_4(\mathcal F)=F_4^{\mathrm{sc}}(\mathcal F)
=(1/4)\cdot127/128\cdot(1/30)/(40320)=127/619315200$.
\end{proof}

\begin{remark}[Fermion versus Heisenberg genus tower]
\label{rem:fermion-vs-heisenberg-genus}
\index{free fermion!comparison with Heisenberg}
Both the free fermion and the Heisenberg algebra have genus expansions
of the form
$F_g(\cA)=F_g^{\mathrm{sc}}(\cA)
=\kappa(\cA) \cdot \lambda_g^{\mathrm{FP}}$
\textup{(}UNIFORM-WEIGHT, class~G exact lane\textup{)}. The ratio at
each genus is:
\[
\frac{F_g(\mathcal{F})}{F_g(\mathcal{H}_k)}
= \frac{\kappa(\mathcal{F})}{\kappa(\mathcal{H}_k)}
= \frac{1}{4k}.
\]
In particular, $F_g(\mathcal{F}) = F_g(\mathcal{H}_{1/4})$: the free
fermion genus tower coincides with that of the Heisenberg at level
$k = 1/4$. The equality follows from
class~G termination: when the shadow obstruction tower terminates at degree~$2$,
the full genus coefficient equals the scalar coefficient and depends
only on $\kappa$, and the combinatorial
genus weights (encoded in the $\hat{A}$-class) are universal.
\end{remark}

\subsubsection{Collision residue and $r$-matrix}
\label{sec:fermion-rmatrix}%
\index{free fermion!$r$-matrix}%

\begin{proposition}[Free fermion $r$-matrix;
\ClaimStatusProvedHere]
\label{prop:fermion-rmatrix}
The collision residue $r(z) = \Res^{\mathrm{coll}}_{0,2}(\Theta_{\mathcal{F}})$ of the free fermion is \emph{regular}:
\begin{equation}\label{eq:fermion-rmatrix}
r(z) = 0.
\end{equation}
No poles survive the $d\log$ extraction.
\end{proposition}

\begin{proof}
The OPE $\psi(z)\psi(w) = 1/(z-w)$ has a simple pole of order~$1$.
The bar construction extracts residues along $d\log(z_i - z_j)$
and thereby lowers the pole order by one in the collision-residue
normalization used here. Thus an OPE term $a_n(z-w)^{-n}$
contributes $a_n(z-w)^{-(n-1)}$ to the polar trace-form residue. For
the free fermion the only singular coefficient is the scalar identity
$a_1=\mathbf 1$, so the extracted term is regular and scalar. The
scalar contraction survives in the ordered bar differential before
reduction, but its projection to the field-valued rational $r$-matrix
has no polar part. Hence $r(z)=0$.

Compare with the Heisenberg, where the double-pole OPE
$\alpha(z)\alpha(w) \sim k/(z-w)^2$ gives
$r_{\mathcal H}(z)=k\Omega_{\mathcal H}/z$ (a simple tensor pole after
$d\log$ absorption, with rank-one coefficient $k/z$).
The one-step pole reduction works uniformly: double pole
$\to$ simple pole (Heisenberg), simple pole $\to$ regular (fermion).
\end{proof}

\begin{remark}[Vanishing $r$-matrix and trivial braiding]
\label{rem:fermion-rmatrix-trivial}
The vanishing of $r(z)$ reflects the absence of nontrivial
braiding on the bar complex of $\mathcal{F}$. For a single
generator with only a simple pole, the collision divisor in
$\overline{C}_2(X)$ contributes no nontrivial $z$-dependent
structure. This is consistent with the genus-$0$ bar complex
collapsing to $\mathbb{C}$
\textup{(}Theorem~\ref{thm:fermion-bar-complex-genus-0}\textup{)}:
a trivial coalgebra admits only the trivial $R$-matrix.
\end{remark}

\subsubsection{Koszul duality and complementarity}
\label{sec:fermion-complementarity}%
\index{free fermion!Koszul duality}%
\index{free fermion!complementarity}%

\begin{proposition}[Free fermion complementarity;
\ClaimStatusConditional]
\label{prop:fermion-complementarity}
The free fermion and its Koszul dual satisfy:
\begin{enumerate}
\item\label{item:fermion-kd}
 Koszul dual:
 $\mathcal{F}^! = \mathrm{Sym}^{\mathrm{ch}}(\gamma)$
 \textup{(}Theorem~\textup{\ref{thm:single-fermion-boson-duality}}\textup{)},
 with one bosonic generator $\gamma$ of conformal weight $h_\gamma = 1/2$.
\item\label{item:fermion-kappa-dual}
 Dual modular characteristic:
 $\kappa(\mathcal{F}^!) = -1/4$.
\item\label{item:fermion-complementarity-sum}
 Free-field complementarity sum:
 $\kappa(\mathcal{F}) + \kappa(\mathcal{F}^!) = 0$ \textup{(}free-field pair\textup{)}.
\item\label{item:fermion-genus-complementarity}
 Scalar genus-$g$ cancellation:
 $\mathrm{obs}_g(\mathcal{F})+\mathrm{obs}_g(\mathcal{F}^!)=0$
 on the uniform-weight lane, equivalently
 $(1/4-1/4)\lambda_g=0$ for all $g\geq 1$.
\end{enumerate}
\end{proposition}

\begin{proof}
\eqref{item:fermion-kd}:
Theorem~\ref{thm:single-fermion-boson-duality}.

\eqref{item:fermion-kappa-dual}:
The Koszul dual $\mathrm{Sym}^{\mathrm{ch}}(\gamma)$ has
$c(\mathcal{F}^!) = -c(\mathcal{F}) = -1/2$
(Koszul duality negates the central charge for free-field families:
$c + c' = 0$), so
$\kappa(\mathcal{F}^!) = c'/2 = -1/4$.

\eqref{item:fermion-complementarity-sum}:
$1/4 + (-1/4) = 0$. This is the free-field case of
Theorem~C \textup{(}Theorem~\ref{thm:quantum-complementarity-main}\textup{)}:
exact cancellation, matching the free-field Heisenberg pair
$\kappa(\mathcal{H}_k) + \kappa(\mathcal{H}_k^!) = k + (-k) = 0$ \textup{(}free-field Heisenberg pair\textup{)}.

\eqref{item:fermion-genus-complementarity}:
Theorem~\ref{thm:quantum-complementarity-main} gives the scalar
Verdier sum for the Koszul pair
$(\mathcal{F}, \mathrm{Sym}^{\mathrm{ch}}(\gamma))$:
the original algebra contributes $\kappa=1/4$, while the
Verdier-dual algebra contributes curvature coefficient $-1/4$.
The statement is only the scalar obstruction cancellation. The
objects remain distinct: $B(\mathcal F)$ is the bar coalgebra,
$\mathcal F^{\mathrm i}=H^\bullet B(\mathcal F)$ is its cohomology
coalgebra, $\mathcal F^!=(\mathcal F^{\mathrm i})^\vee$ is the
Verdier-linear Koszul-dual algebra, and
$Z_{\mathrm{ch}}^{\mathrm{der}}(\mathcal F)$ is the derived center.
\end{proof}

\begin{remark}[Fermion--boson complementarity as statistics exchange]
\label{rem:fermion-statistics-exchange}
The fermion--boson Koszul duality $\mathcal{F}^! = \mathrm{Sym}^{\mathrm{ch}}(\gamma)$ exchanges
Fermi and Bose statistics at the level of generators, while preserving
the conformal weight $h = 1/2$. The bar complex encodes this exchange:
$\barBgeom(\mathcal{F})$ collapses (fermionic antisymmetry kills bar
elements), while $\barBgeom(\mathrm{Sym}^{\mathrm{ch}}(\gamma))$ grows
polynomially (bosonic symmetry permits all bar words). The
free-field complementarity sum $\kappa + \kappa' = 0$ is exact for this pair, as
for all free-field Koszul pairs \textup{(}this holds for
KM\slash free fields; W-algebras require the general formula
$\kappa + \kappa' = \rho \cdot K$\textup{)}.
\end{remark}

\subsubsection{Sewing and the genus-$g$ partition function}
\label{sec:fermion-sewing}%
\index{free fermion!sewing}%
\index{Pfaffian!free fermion partition function}%

\begin{theorem}[Free fermion sewing;
\ClaimStatusProvedHere]
\label{thm:fermion-sewing}
\index{free fermion!genus-$g$ partition function}
The genus-$g$ partition function of the free fermion on a closed Riemann
surface $\Sigma_g$ with period matrix $\Omega \in \mathbb{H}_g$ and
spin structure $\sigma \in (\mathbb{Z}/2)^{2g}$ is:
\begin{equation}\label{eq:fermion-partition}
Z_g(\mathcal{F};\, \sigma)
= \mathrm{Pf}\bigl(\bar{\partial}_{K^{1/2}_\sigma}\bigr)
= \bigl(\det \bar{\partial}_{K^{1/2}_\sigma}\bigr)^{1/2}
= \theta[\sigma](0,\Omega)^{1/2},
\end{equation}
where $K^{1/2}_\sigma$ is the spin bundle on $\Sigma_g$ determined by
$\sigma$, the Pfaffian $\mathrm{Pf}$ is the square root of the functional
determinant of the $\bar\partial$ operator on $K^{1/2}_\sigma$, and
$\theta[\sigma]$ is the theta function with characteristic $\sigma$.
\end{theorem}

\begin{proof}
The free fermion path integral on $\Sigma_g$ is Gaussian:
the action $S[\psi] = \int_{\Sigma_g} \psi\, \bar{\partial}\, \psi$
is quadratic, and the functional integral evaluates to the Pfaffian
of the kinetic operator $\bar{\partial}_{K^{1/2}_\sigma}$. The fermion
$\psi$ is a section of $K^{1/2}_\sigma$ (in the NS sector;
the Ramond sector replaces $\sigma$ by a shifted characteristic),
so the operator is $\bar{\partial}: \Gamma(K^{1/2}_\sigma) \to
\Gamma(K^{1/2}_\sigma \otimes \bar{K})$.

The Pfaffian, rather than the determinant, is the correct answer because
$\psi$ is a single real (Majorana) fermion: the pairing
$\psi(z)\psi(w) \sim 1/(z-w)$ is antisymmetric, and the functional
integral over an antisymmetric quadratic form is a Pfaffian. This
is the structural difference from the Heisenberg, whose bosonic
partition function is a Fredholm determinant
$Z_g(\mathcal{H}_k) = (\det' \bar{\partial}_1)^{-k/2}$
\textup{(}Theorem~\ref{thm:heisenberg-genus-g}\textup{)}.

The identification with the theta function follows from the
bosonization correspondence: $\det \bar{\partial}_{K^{1/2}_\sigma}
= \theta[\sigma](0,\Omega)$ on a genus-$g$ surface with period
matrix~$\Omega$, so the Pfaffian is $\theta[\sigma]^{1/2}$.
The $2^{2g}$ spin structures yield $2^{2g}$ distinct partition
functions; the GSO projection selects a modular-invariant combination.
\end{proof}

\begin{remark}[Pfaffian versus determinant: the fermionic signature]
\label{rem:pfaffian-vs-determinant}
\index{free fermion!Pfaffian vs determinant}
The distinction between Pfaffian \textup{(}fermion\textup{)} and
Fredholm determinant \textup{(}Heisenberg\textup{)} is the fundamental
structural dichotomy in free-field sewing:
\begin{center}
\small
\renewcommand{\arraystretch}{1.2}
\begin{tabular}{lcc}
\toprule
& \textbf{Free fermion $\mathcal{F}$}
 & \textbf{Heisenberg $\mathcal{H}_k$} \\
\midrule
Statistics & fermionic & bosonic \\
OPE pole & simple ($z^{-1}$) & double ($z^{-2}$) \\
Central charge & $c = 1/2$ & $c = 1$ per generator \\
$\kappa$ & $1/4$ \textup{(fixed)} & $k$ \textup{(tunable)} \\
Shadow class & G & G \\
Partition function & $\mathrm{Pf}(\bar\partial_{K^{1/2}})$
 & $(\det' \bar\partial)^{-k/2}$ \\
Spin structure & required ($2^{2g}$ choices)
 & not required \\
Bosonization & after pairing two real fermions: $\mathcal{H}_1$ on
 lattice $\mathbb{Z}$
 & n/a \\
\bottomrule
\end{tabular}
\end{center}
The Pfaffian is a square root of the determinant:
$\mathrm{Pf}(\bar\partial)^2 = \det(\bar\partial)$. Two free fermions pair to
a single complex fermion, whose partition function \emph{is} a
determinant; the boson-fermion correspondence then identifies this with
the Heisenberg at $k = 1$ on the $\mathbb{Z}$-lattice. This is the
genus-$g$ manifestation of the classical bosonization
$\mathcal{F} \otimes \mathcal{F} \cong \mathcal{H}_1|_{\mathbb{Z}}$.
\end{remark}

\subsubsection{The free fermion as class~G archetype}
\label{sec:fermion-class-g-archetype}%
\index{free fermion!class G archetype}%

\begin{proposition}[Free fermion characteristic data;
\ClaimStatusConditional]
\label{prop:fermion-characteristic-data}
\index{free fermion!characteristic hierarchy}
The characteristic hierarchy of the free fermion is:
\begin{center}
\small
\renewcommand{\arraystretch}{1.2}
\begin{tabular}{ll}
\toprule
\emph{Invariant} & \emph{Value for $\mathcal{F}$} \\
\midrule
Central charge $c$ & $1/2$ \\
Modular characteristic $\kappa$ & $1/4$ \\
Anomaly ratio $\rho = \kappa/c$ & $1/2$ \\
Shadow depth $r_{\max}$ & $2$ \\
Shadow class & G \textup{(Gaussian)} \\
Shadow metric $Q(t)$ & $1/4$ \textup{(constant)} \\
Critical discriminant $\Delta$ & $0$ \\
$r$-matrix $r(z)$ & $0$ \textup{(regular)} \\
Koszul dual $\mathcal{F}^!$ & $\mathrm{Sym}^{\mathrm{ch}}(\gamma)$,
 $h_\gamma = 1/2$ \\
$\kappa(\mathcal{F}^!)$ & $-1/4$ \\
Complementarity sum & $0$ \\
Genus tower & $F_g = (1/4)\,\lambda_g^{\mathrm{FP}}$ \textup{(}UNIFORM-WEIGHT\textup{}) \\
$F_1$ & $1/96$ \\
Genus partition function & $\mathrm{Pf}(\bar\partial_{K^{1/2}_\sigma})$ \\
Spin structures & $2^{2g}$ choices per genus \\
\bottomrule
\end{tabular}
\end{center}
All entries are projections of the single MC element
$\Theta_{\mathcal{F}} = (1/4) \cdot \eta \otimes \Lambda
\in \MC(\gAmod[\mathcal{F}])$
\textup{(}Theorem~\textup{\ref{thm:mc2-bar-intrinsic}}\textup{)}.
\end{proposition}

\begin{proof}
Each entry is computed above:
$\kappa$ in Proposition~\ref{prop:fermion-shadow-invariants},
$r(z)$ in Proposition~\ref{prop:fermion-rmatrix},
complementarity in Proposition~\ref{prop:fermion-complementarity},
$F_g$ in Theorem~\ref{thm:fermion-genus-expansion},
sewing in Theorem~\ref{thm:fermion-sewing}.
The anomaly ratio $\rho = \kappa/c = (1/4)/(1/2) = 1/2$ matches the
Virasoro scalar-lane value $\kappa=c/2$. The rank-one Heisenberg is the
contrasting Gaussian convention: it has fixed central charge~$1$ but
variable modular characteristic $\kappa(\mathcal H_k)=k$.
\end{proof}

\begin{remark}[Extensions: lattice VOAs and superconformal algebras]
\label{rem:fermion-extensions}
\index{free fermion!extensions}
The free fermion extends in two orthogonal directions:
\begin{enumerate}
\item \emph{Lattice extension.}
 The boson-fermion correspondence identifies
 $\mathcal{F} \otimes \mathcal{F}$ with the
 Heisenberg algebra on the lattice $\mathbb{Z}$, and more
 generally $\mathcal{F}^{\otimes d}$ with a lattice VOA on
 $\mathbb{Z}^{d/2}$ \textup{(}for $d$ even\textup{)}. The shadow
 tower of the tensor product inherits the class~G termination by
 additivity \textup{(}Proposition~\ref{prop:independent-sum-factorization}:
 $\kappa(\mathcal{F}^{\otimes d}) = d/4$\textup{)}.

\item \emph{Superconformal extension.}
 Combining $\mathcal{F}$ with a Virasoro algebra of central
 charge $c_{\mathrm{Vir}}$ yields the $N = 1$ superconformal algebra
 at $c = c_{\mathrm{Vir}} + 1/2$. The fermion sector contributes
 $\kappa_{\mathrm{ferm}} = 1/4$ to the total $\kappa$, and the
 shadow class of the combined system is determined by the Virasoro
 factor \textup{(}class~M, $r_{\max} = \infty$\textup{)}: the
 fermionic contribution does not increase shadow depth.
\end{enumerate}
\end{remark}

\subsection{\texorpdfstring{The $\beta\gamma$ system}{The beta-gamma system}}
\index{beta-gamma system|textbf}

\begin{remark}[Shadow archetype: contact/quartic]
\label{rem:betagamma-shadow-archetype-free}
\index{beta-gamma system@$\beta\gamma$ system!shadow archetype}
Shadow depth $r_{\max} = 4$ (class~C).
Modular characteristic $\kappa(\beta\gamma) = c/2 = 6\lambda^2 - 6\lambda + 1$,
where $\lambda$ is the conformal weight of $\beta$.
Central charge $c_{\beta\gamma} = 2(6\lambda^2 - 6\lambda + 1)$.
The scalar shadow $\Theta^{\leq 2} = \kappa \cdot \eta \otimes \Lambda$
is nonzero at generic $\lambda$, but on the weight-changing line
$\Theta^{\leq 2}\big|_{\mathrm{w.c.}} = 0$ and the entire
weight-changing slice tower vanishes. On the T-line, by contrast, the
recomputed one-dimensional tower is the Virasoro tower and has infinite
depth. The class~C witness is the standard conformal-weight family
$\beta\gamma_\lambda$:
$S_2 = 6\lambda^2 - 6\lambda + 1$, $S_3 = 0$,
$S_4 = -5/12$, and $S_r = 0$ for $r \geq 5$.
Tower terminates at exactly $r = 4$ by
Proposition~\ref{prop:depth-gap-trichotomy}; the
weight-changing slice is shadow-trivial by
Theorem~\ref{thm:betagamma-rank-one-rigidity};
see \S\ref{sec:betagamma-shadow-tower-free} for the full tower.
Koszul dual: $(\beta\gamma)^! = bc$
(Theorem~\ref{thm:betagamma-bc-koszul}).
See Chapter~\ref{chap:beta-gamma} for the full quartic
birth analysis.
\end{remark}

The standard $\beta\gamma_\lambda$ system combines two bosonic
fields with conformal weights
$h_\beta=\lambda$ and $h_\gamma=1-\lambda$. The physical
weight-$(1,0)$ system is the special point $\lambda=1$. Its bar
complex shares the discriminant $\Delta(x) = (1-3x)(1+x)$ with
$\widehat{\mathfrak{sl}}_2$ and the Virasoro algebra. The
Koszul dual is the $bc$ ghost system (\S\ref{sec:fermion-boson-koszul}).
Its pole-valued collision residue is
$r^{\mathrm{coll}}_{\beta\gamma}(z)=0$; the nonzero ordered datum is
the regular contact operator
$\Theta_{\beta\gamma}^{\mathrm{ord}}
=\beta\otimes\gamma-\gamma\otimes\beta$, with
$R_{\mathrm{contact}}=\id+2\pi i\,\Theta_{\beta\gamma}^{\mathrm{ord}}$.

\begin{definition}[\texorpdfstring{$\beta\gamma$}{beta-gamma} system]
The $\beta\gamma$ chiral algebra is generated by:
\begin{itemize}
\item $\beta(z)$ of conformal weight $h_\beta = \lambda$,
\item $\gamma(z)$ of conformal weight $h_\gamma = 1-\lambda$.
\end{itemize}
with OPEs:
\[
\beta(z)\gamma(w) = \frac{1}{z-w} + \text{regular}, \quad
\gamma(z)\beta(w) = -\frac{1}{z-w} + \text{regular}
\]
The relation $R_{\beta\gamma} = \beta \otimes \gamma - \gamma \otimes \beta$ enforces normal ordering.
\end{definition}
\index{normal ordering!free fields}


\subsection{\texorpdfstring{The $bc$ ghosts}{The bc ghosts}}
\index{$bc$ system|textbf}

\begin{remark}[Shadow archetype: contact/quartic, Koszul dual of $\beta\gamma$]
\label{rem:bc-shadow-archetype}
\index{$bc$ system!shadow archetype}
Shadow depth $r_{\max} = 4$ (class~C), the same as its Koszul dual
$\beta\gamma$. The $bc$ system has two fermionic generators ($b$ of
weight $\lambda$, $c$ of weight $1 - \lambda$) and a single
simple-pole OPE, giving it the same deformation-complex structure as
$\beta\gamma$: a T-line (Virasoro subalgebra at $c_{bc}$, with
$S_3 = 2$ and $S_4 = 10/[c_{bc}(5c_{bc}+22)]$) and a
weight-changing line (cubic vanishes by rank-one abelian rigidity).
The class-C witness is again the standard conformal-weight family,
for which $S_3 = 0$, $S_4 = -5/12$, and $S_r = 0$ for $r \ge 5$,
exactly as for $\beta\gamma$.

This is distinct from the single-generator free fermion
$\mathcal{F}$ (class~G, $r_{\max} = 2$), where fermionic
antisymmetry on a single generator kills all higher bar cohomology.
The $bc$ system has two generators, so the antisymmetry argument
does not apply.

Modular characteristic $\kappa(bc) = c_{bc}/2 = -(6\lambda^2 - 6\lambda + 1)$.
At the physical point $(h_b, h_c) = (2, -1)$ (i.e.\ $\lambda = 2$):
$c_{bc} = -26$, $\kappa(bc) = -13$.
Complementarity:
$\kappa(\beta\gamma) + \kappa(bc) = 0$ (exact cancellation, free-field family).
Shadow depth is preserved by Koszul duality for this pair: both
$\beta\gamma$ and $bc$ are class~C with $r_{\max} = 4$.
The pole-valued collision residue is still zero:
$r^{\mathrm{coll}}_{bc}(z)=0$. The nonzero ordered datum is the
regular graded-Clifford contact relation
$\Theta_{bc}^{\mathrm{ord}}=b\otimes c+c\otimes b$.
\end{remark}

\begin{definition}[\texorpdfstring{$bc$}{bc} ghost system]
Generated by $b(z)$ of weight $h_b = 2$ and $c(z)$ of weight $h_c = -1$,
with OPE $b(z)c(w) = \frac{1}{z-w}$ and Clifford relation $R_{bc} = b \otimes c + c \otimes b$.
\end{definition}

\begin{proposition}[$bc$ ghost system at general spin: class~C for all $\lambda$]
\label{prop:bc-general-spin-class-c}
\ClaimStatusProvedElsewhere
\index{$bc$ system!general spin $\lambda$}
\index{shadow depth!independence of spin}
Let $bc_\lambda$ denote the fermionic ghost system with generators $b$
of conformal weight $\lambda$ and $c$ of conformal weight $1 - \lambda$,
OPE $b(z)c(w) = 1/(z-w)$, and Clifford relation
$R_{bc_\lambda} = b \otimes c + c \otimes b$. For every $\lambda \in
\mathbb{C}$, the bar complex of $bc_\lambda$ has shadow depth exactly
$r_{\max} = 4$ (class~C, contact/quartic archetype), with Virasoro
subalgebra central charge
\[
 c_{bc}(\lambda) \;=\; 1 - 3(2\lambda - 1)^2
 \;=\; -2(6\lambda^2 - 6\lambda + 1),
\]
and modular characteristic
$\kappa(bc_\lambda) = c_{bc}(\lambda)/2 = -(6\lambda^2 - 6\lambda + 1)$.
The shadow coefficient depends only on $\lambda$, not on any level
parameter: the $bc$ system has no Heisenberg level, only the
conformal-weight datum $\lambda$. Koszul complementarity with $\beta\gamma_\lambda$
holds for every $\lambda$:
\[
 \kappa(\beta\gamma_\lambda) + \kappa(bc_\lambda) \;=\; 0 \quad\textup{(}free-field pair\textup{)},
 \qquad c_{\beta\gamma}(\lambda) + c_{bc}(\lambda) \;=\; 0
 \;\; \]
\end{proposition}

\begin{proof}
The PBW computation at degrees $n \leq 4$ shows the same contact/quartic
shadow stratum as $\beta\gamma_\lambda$ (with the fermion--boson roles
exchanged), and the rank-one abelian rigidity argument
(Theorem~\ref{thm:betagamma-rank-one-rigidity}) applies verbatim once
the single simple-pole OPE is fixed. Stratum separation at degree~$5$
kills the quintic obstruction for every $\lambda$, so $r_{\max} = 4$
holds as a function of $\lambda$. Numerical verification:
at $\lambda = 2$ (standard $bc$), $c_{bc} = 1 - 3 \cdot 9 = -26$ and
$\kappa(bc) = -13$, recovering the physical ghost point;
at $\lambda = 1/2$ (the weight-symmetric point), $c_{bc} = 1$ and
$\kappa(bc) = 1/2$. This is a complex fermion point, not a
Koszul self-duality point: the Verdier-linear dual is still
$\beta\gamma_{1/2}$, with $\kappa=-1/2$.
Cross-check against the $bc$ shadow engines in
\texttt{compute/tests/} (the full family of $bc$-shadow tests all
evaluate a weight-independent shadow depth of $4$).
\end{proof}


\subsection{\texorpdfstring{Free fermion $\leftrightarrow$ free boson: single-generator Koszul duality}{Free fermion--free boson: single-generator Koszul duality}}

\begin{theorem}[Single-generator fermion-boson duality; \ClaimStatusProvedHere]\label{thm:single-fermion-boson-duality}
The free fermion $\mathcal{F} = \Lambda^{\mathrm{ch}}(\psi)$ (one fermionic generator) and the free boson $\mathrm{Sym}^{\mathrm{ch}}(\gamma)$ (one bosonic generator) form a Koszul pair:
\[\mathcal{F}^! \cong \mathrm{Sym}^{\mathrm{ch}}(\gamma), \qquad
 (\mathrm{Sym}^{\mathrm{ch}}(\gamma))^! \cong \mathcal{F}.\]
This is the chiral incarnation of the operadic duality $\mathrm{Lie}^! = \mathrm{Com}$
in a single generator.
\end{theorem}

\begin{remark}
This single-generator duality is distinct from the \emph{two-generator} duality $(\beta\gamma)^! \cong bc$ of Theorem~\ref{thm:betagamma-bc-koszul}. Here both sides have $\dim V = 1$; in the $\beta\gamma$--$bc$ duality, both sides have $\dim V = 2$.
\end{remark}

\begin{proof}[Verification of Koszul pair conditions]
\noindent\emph{Generators and weights.}
\begin{itemize}
\item $\mathcal{F}$: one fermionic generator $\psi$ with $h_\psi = 1/2$, so $V = \mathbb{C}\psi$ ($\dim V = 1$)
\item $\mathrm{Sym}^{\mathrm{ch}}(\gamma)$: one bosonic generator $\gamma$ with $h_\gamma = 1/2$, so $V^* = \mathbb{C}\gamma$ ($\dim V^* = 1$)
\end{itemize}

\noindent\emph{Pairing.} The residue pairing $V \times V^* \to \mathbb{C}$ is:
\[\langle\psi, \gamma\rangle = \mathrm{Res}_{z=w}\left[\psi(z)\gamma(w) \cdot 1\right] = 1\]
This is a non-degenerate pairing between 1-dimensional spaces.

\noindent\emph{Relations.}
\begin{itemize}
\item $R_{\mathrm{ferm}} = V^{\otimes 2}$, generated by the
 graded-symmetric tensor $\psi\otimes\psi$; in characteristic
 zero this is the quadratic relation $\psi^2=0$ of the exterior
 algebra on the odd generator.
\item $R_{\mathrm{bos}} = 0$; the polynomial algebra on the
 even generator has no quadratic relation.
\end{itemize}

\noindent\emph{Orthogonality.}
The quadratic dual relation is
\[
R_{\mathrm{ferm}}^\perp
=\{\rho\in (V^*)^{\otimes 2}\mid
\rho(R_{\mathrm{ferm}})=0\}.
\]
Since $R_{\mathrm{ferm}}=V^{\otimes 2}$ is one-dimensional and
$\langle \psi\otimes\psi,\gamma\otimes\gamma\rangle=1$,
one has $R_{\mathrm{ferm}}^\perp=0=R_{\mathrm{bos}}$.
Conversely, $R_{\mathrm{bos}}=0$ gives
$R_{\mathrm{bos}}^\perp=V^{\otimes 2}=R_{\mathrm{ferm}}$.
Thus the exterior relation on the fermion and the absence of a
quadratic relation on the boson are orthogonal complements.

\noindent\emph{Acyclicity.} The quadratic Koszul complex is the
one-generator exterior--symmetric Koszul complex
\[
\Lambda(\psi)\otimes \mathrm{Sym}(\gamma)
\]
with differential induced by the pairing
$\langle\psi,\gamma\rangle=1$; it is acyclic off the vacuum.
The chiral statement is the corresponding Verdier-linear
realization on the Ran space. The cobar counit
$\Omega B(\mathcal F)\simeq\mathcal F$ is the inversion
statement, not the definition of the dual algebra.
\end{proof}





\section{Bar complexes of free-field archetypes}\label{sec:bar-complexes-free-fields}

The bar complex of each free-field archetype is determined by OPE
residues along collision divisors in the Fulton--MacPherson
compactification; the results range from genus-$0$ cohomological
collapse (free fermion), through a finite Heisenberg shadow with
level central class, to the infinite two-generator growth of
$\beta\gamma$ and the lattice-graded structure of lattice VOAs.

\subsection{Free fermion bar complex}

Fermionic antisymmetry
($\psi(z)\psi(w) = -\psi(w)\psi(z)$) forces the reduced
genus-$0$ bar cohomology on $\mathbb P^1$ to collapse. The mechanism
is not a literal prohibition on several insertions; it is the
incompatibility between the alternating tensor factor and the cyclic
eigenspaces of logarithmic configuration-space forms. Cyclic
permutations of the relevant Arnold-form spaces have nontrivial roots
of unity as eigenvalues, leaving no invariant class in positive
reduced bar degree.

\begin{theorem}[Free fermion bar complex at genus 0; \ClaimStatusProvedHere]
\label{thm:fermion-bar-complex-genus-0}
For the free fermion $\mathcal{F}$ on $\mathbb{P}^1$, the reduced
geometric bar complex has only the vacuum cohomology:

\[
H^n(\bar{B}_{geom}(\mathcal{F})) =
\begin{cases}
\mathbb{C} & n = 0\\
0 & n \geq 1
\end{cases}.
\]

At genus $g \geq 1$, additional contributions arise from the topology of $\Sigma_g$; see Theorem~\ref{thm:elliptic-fermion-bar} for genus~1.
\end{theorem}

\begin{proof}
\emph{Degree~0.} $\bar{B}^0_{geom} = \mathbb{C} \cdot 1$ (vacuum state).

\emph{Degree~1.} Elements have form
$\alpha = \int_{C_2(X)} \psi(z_1) \otimes \psi(z_2) \otimes f(z_1,z_2)\eta_{12}$

The differential:
\begin{align}
d\alpha &= \text{Res}_{D_{12}}[\mu_{12}(\psi \otimes \psi) \otimes f\eta_{12}]\\
&= \text{Res}_{z_1=z_2}\left[\frac{1}{z_1-z_2} \cdot f(z_1,z_2) \cdot \frac{dz_1-dz_2}{z_1-z_2}\right]
\end{align}

Expanding $f$ near the diagonal:
$f(z_1, z_2) = f(z, z) + (z_1 - z_2)\partial_1 f|_z + (z_2 - z_1)\partial_2 f|_z + O((z_1 - z_2)^2)$

Since $\psi(z_1)\psi(z_2) = -\psi(z_2)\psi(z_1)$, the function $f$ must be antisymmetric: $f(z_1, z_2) = -f(z_2, z_1)$. This implies $f(z, z) = 0$ and $\partial_2 f = -\partial_1 f$.

The residue extracts the coefficient of $(z_1 - z_2)^{-1}$ in:
$\frac{1}{z_1 - z_2} \cdot [(z_1 - z_2)\partial_1 f|_z + (z_1 - z_2)\partial_1 f|_z] \cdot \frac{dz_1 - dz_2}{z_1 - z_2}$
$= \frac{2(z_1 - z_2)\partial_1 f|_z \cdot (dz_1 - dz_2)}{(z_1 - z_2)^2}$
$= \frac{2\partial_1 f|_z \cdot (dz_1 - dz_2)}{z_1 - z_2}$

The residue gives $2\partial_1 f|_z \cdot dz = df|_{\text{diagonal}}$ (the factor of 2 cancels with the $1/2$ from symmetrization).

Thus the only possible degree-$1$ class maps to
$H^1(X,\mathbb{C})$. For $X=\mathbb{P}^1$ this group vanishes, so
$H^1(\bar B_{\mathrm{geom}}(\mathcal F))=0$.

\emph{Degree~2.} Elements would be $\psi_1 \otimes \psi_2 \otimes \psi_3 \otimes \omega$ with $\omega \in \Omega^2(C_3(X))$.

By fermionic antisymmetry:
$\psi_1 \otimes \psi_2 \otimes \psi_3 = -\psi_2 \otimes \psi_1 \otimes \psi_3 = -\psi_1 \otimes \psi_3 \otimes \psi_2 = \psi_3 \otimes \psi_1 \otimes \psi_2$

Under cyclic permutation $(123) \to (312)$:
$\omega = g(z_1,z_2,z_3)\eta_{12} \wedge \eta_{23} \mapsto g(z_3,z_1,z_2)\eta_{31} \wedge \eta_{12}$

By Arnold, the space of $2$-forms $\Omega^2_{\log}(\overline{C}_3)$ is $2$-dimensional, spanned by $\eta_{12} \wedge \eta_{23}$ and $\eta_{23} \wedge \eta_{31}$. The cyclic permutation $\sigma: (123) \mapsto (231)$ acts on this space with matrix
$\left[\begin{smallmatrix} 0 & -1 \\ 1 & -1 \end{smallmatrix}\right]$
(in this basis), whose eigenvalues are the primitive cube roots of unity $\zeta_3, \zeta_3^2$. Since neither eigenvalue equals~$1$, the fixed subspace under $\sigma$ is $\{0\}$. But fermionic antisymmetry forces the form component to be $\sigma$-invariant, hence the degree-$2$ class is zero.

\emph{Higher degrees.} The reduced genus-$0$ fermion complex is
generated by totally antisymmetric tensor factors tensored with
logarithmic Arnold forms. For $n\geq 3$ the alternating tensor
requires the Arnold-form component to lie in the invariant part of a
non-trivial cyclic eigenspace. The Arnold relation has no such
invariant component in top logarithmic degree, so the same argument
as in degree~$2$ kills every higher class. Hence the only surviving
cohomology on $\mathbb{P}^1$ is the vacuum.
\end{proof}



\subsection{Free fermion coalgebra structure}

\begin{theorem}[Fermion bar complex coalgebra; \ClaimStatusProvedHere]\label{thm:fermion-bar-coalg}
The coaugmented bar complex
$B^{\text{ch}}(\mathcal{F})=\mathbb{C}\oplus
\bar{B}^{\text{ch}}(\mathcal{F})$ carries the chiral coalgebra structure:
\begin{enumerate}
\item \emph{Comultiplication.} For $\alpha = \psi_1 \otimes \cdots \otimes \psi_n \otimes \omega \in \bar{B}^n$:
\[
\Delta(\alpha) = \sum_{I \sqcup J = [n], 1 \in I} \text{sign}(\sigma) \cdot \alpha_I \otimes \alpha_J
\]
where $\alpha_I = \bigotimes_{i \in I} \psi_i \otimes \omega|_{C_{|I|}(X)}$ and $\sigma$ is the shuffle permutation.

\item \emph{Counit.} The coaugmented bar coalgebra
$B^{\mathrm{ch}}(\mathcal{F})=\mathbb{C}\oplus
\bar{B}^{\mathrm{ch}}(\mathcal{F})$ has counit
$\epsilon(1)=1$ and
$\epsilon(\alpha)=0$ for every positive reduced bar word
$\alpha\in\bar{B}^{\mathrm{ch}}(\mathcal{F})$.

\item \emph{Graded symmetry.} The fermionic sign structure satisfies:
\[
\tau(\psi_1 \otimes \cdots \otimes \psi_n) = (-1)^{n(n-1)/2} \psi_n \otimes \cdots \otimes \psi_1
\]
expressing the reversal of tensor factors with Koszul signs.
\end{enumerate}
\end{theorem}

\begin{proof}[Geometric construction]
The coalgebra structure arises from the stratification of $\overline{C}_n(X)$ by collision patterns.

\noindent\emph{Comultiplication from boundary strata.}
The FM boundary $\partial\overline{C}_n(X)$ decomposes into divisors $D_{I,J}$ indexed by ordered bi-partitions $[n] = I \sqcup J$ with $1 \in I$ and $|I|, |J| \geq 1$. Restricting the degree-$n$ generator $\alpha = \psi_1 \otimes \cdots \otimes \psi_n \otimes \omega$ to the divisor $D_{I,J}$ gives:
\[
\Delta(\alpha)|_{D_{I,J}} = \mathrm{sign}(\sigma_{I,J})\;\alpha_I \otimes \alpha_J
\]
where $\sigma_{I,J}$ is the shuffle sending $(1,\ldots,n)$ to $(i_1,\ldots,i_{|I|},j_1,\ldots,j_{|J|})$, and $\alpha_I = \bigotimes_{i\in I}\psi_i \otimes \omega|_{C_{|I|}(X)}$ uses restriction of the log-form $\omega$ to the appropriate stratum. Summing over all bi-partitions gives the formula stated in (1).

\noindent\emph{Coassociativity.}
Coassociativity $(\Delta \otimes \mathrm{id}) \circ \Delta
= (\mathrm{id} \otimes \Delta) \circ \Delta$ follows from the
associativity of the FM stratification: the codimension-$2$ boundary
of $\overline{C}_n(X)$ is covered by iterated $D_{I,J}$-strata in two
orderings, and the two resulting expressions for
$(\Delta \otimes \mathrm{id}) \circ \Delta(\alpha)$ agree by the
Arnold relation on orientation forms
(Theorem~\ref{thm:arnold-relations-appendix}).

\noindent\emph{Counit.}
The counit is the standard counit of the cofree conilpotent coalgebra:
it projects to the coaugmentation summand $\mathbb{C}\cdot 1$ and
vanishes on the reduced tensor coalgebra. Hence
$(\epsilon\otimes\mathrm{id})\circ\Delta
=\mathrm{id}=(\mathrm{id}\otimes\epsilon)\circ\Delta$ on
$B^{\mathrm{ch}}(\mathcal{F})$ by the deconcatenation formula.

\noindent\emph{Chain map property.}
The differential $d_{\bar{B}}$ and the comultiplication $\Delta$ are
compatible (i.e.\ $\Delta$ is a chain map
$(\bar{B}^{\mathrm{ch}}(\mathcal{F}), d_{\bar{B}}) \to
(\bar{B}^{\mathrm{ch}}(\mathcal{F}) \otimes
\bar{B}^{\mathrm{ch}}(\mathcal{F}), d \otimes 1 + 1 \otimes d)$) by
the Leibniz rule for the bar differential on FM strata: the
differential and comultiplication both arise from boundary residues,
and Stokes' theorem $\partial \circ \partial = 0$ ensures they commute
up to sign in the standard way for dg coalgebras, with the relevant
Arnold relations supplied by
Theorem~\ref{thm:arnold-relations-appendix}.

\noindent\emph{Graded symmetry.}
The sign formula in (3) encodes the Koszul sign rule for the fermion system: swapping $\psi_i$ and $\psi_j$ in the tensor product introduces a sign $(-1)^{|\psi_i||\psi_j|}$. Since each fermion has odd degree, this produces $(-1)^1 = -1$ for each transposition, so reversing all $n$ factors gives $(-1)^{n(n-1)/2}$, matching the stated formula.
\end{proof}


\subsection{\texorpdfstring{$\beta\gamma$}{beta-gamma} bar complex computation}

\begin{theorem}[\texorpdfstring{$\beta\gamma$}{beta-gamma} bar complex; \ClaimStatusProvedHere]\label{thm:betagamma-bar-complex}
Let $\barB^{\mathrm{lw}}(\beta\gamma_\lambda)$ denote the
lowest-weight generator-type quotient of the chiral bar complex,
obtained by keeping the two primary generators
$\beta,\gamma$ and the first derivative sector created by the
simple-pole residue, and by discarding higher conformal descendants.
Then
\begin{align}
\dim(\bar{B}^{\mathrm{lw},1}) &= 2, \\
\dim(\bar{B}^{\mathrm{lw},2}) &= 6.
\end{align}
The unreduced bar group $\barB^n(\beta\gamma_\lambda)$ is
infinite-dimensional for every $n\geq 1$, because
$\overline{\beta\gamma_\lambda}$ contains arbitrarily high
conformal descendants. The displayed numbers count only the
finite generator-type sectors used in the shadow computation.
\end{theorem}

\begin{proof}
The quotient keeps the two primary generator types in bar degree~$1$,
so $\bar B^{\mathrm{lw},1}$ has basis represented by
$\beta$ and $\gamma$. In bar degree~$2$ the retained pure tensor
types are
\[
\beta\beta,\qquad \beta\gamma,\qquad
\gamma\beta,\qquad \gamma\gamma.
\]
The simple-pole contractions
$\beta(z)\gamma(w)\sim (z-w)^{-1}$ and
$\gamma(z)\beta(w)\sim -(z-w)^{-1}$ create exactly the first
descendant sectors $\partial\beta$ and $\partial\gamma$ after
collision extraction. Higher descendants are excluded by the
definition of $\bar B^{\mathrm{lw}}$. Hence the retained degree-$2$
sector has the four pure tensor types plus two first-derivative
types, giving $\dim \bar B^{\mathrm{lw},2}=6$.
\end{proof}

\begin{theorem}[\texorpdfstring{$\beta\gamma$}{beta-gamma} bar complex rank; \ClaimStatusConditional]\label{thm:betagamma-bar-dim}
For all $n \geq 1$,
\[
\operatorname{rank}\bigl(\bar{B}^{\mathrm{lw},n}_{\mathrm{geom}}
(\beta\gamma_\lambda)\bigr)=2\cdot 3^{n-1}.
\]
This rank is the number of independent lowest-weight generator-type
sectors. It is neither the vector-space dimension of the full
geometric bar group nor the dimension of bar cohomology.
\end{theorem}

\begin{proof}
We prove by induction on~$n$.

\emph{Base cases.} The cases $n = 1, 2, 3$ are verified explicitly in the
complete bar complex computation (Theorem~\ref{thm:betagamma-bar-cohomology},
Chapter~\ref{chap:beta-gamma}):
$\operatorname{rank}(\bar{B}^{\mathrm{lw},1}) = 2$ ($\beta$- and $\gamma$-sectors),
$\operatorname{rank}(\bar{B}^{\mathrm{lw},2}) = 6$ ($2^2 = 4$ pure tensor types $+$ $2$
derivative types $V_{\partial\beta}, V_{\partial\gamma}$),
$\operatorname{rank}(\bar{B}^{\mathrm{lw},3}) = 18$ ($2^3 = 8$ pure tensor $+$ $10$ derivative).

\emph{Type decomposition.} The lowest-weight geometric bar model at
bar degree~$n$
decomposes by \emph{type assignment}: at each position
$i \in \{1, \ldots, n\}$, assign a type $t_i$ from the alphabet
$\{\beta, \gamma, \partial\}$, subject to the constraint $t_1 \in \{\beta, \gamma\}$
(no descendant at the root). Here $\partial$ denotes the conformal
descendant sector arising from the simple-pole OPE
$\beta(z)\gamma(w) \sim 1/(z{-}w)$: the first-order pole contributes
exactly one new generator type per non-root position (the residue
along each collision divisor $D_{ij}$ in the NBC spanning tree).

\emph{Independence.} The type sectors carry distinct quantum numbers
under the PBW charge and derivative filtrations. At the standard
weight-$(1,0)$ point $\lambda=1$, this can be read from conformal
weights: $\beta$ has leading weight~$1$, $\gamma$ has leading
weight~$0$, and the first descendant sector has weight~$1$ but
positive derivative degree. For generic~$\lambda$, the same separation
is made by the ordered pair
\[
(\text{charge in the }\beta/\gamma\text{ lattice},
 \text{ derivative degree}),
\]
so the exceptional weight coincidences at $\lambda=0,1/2,1$ do not
identify the retained PBW type sectors. Hence the sectors are linearly
independent in the lowest-weight generator-type quotient.

\emph{Exhaustion.} Every lowest-weight type element belongs
to a unique type sector, since the PBW basis of the $\beta\gamma$ vertex
algebra provides a weight-graded decomposition of
$(s^{-1}\bar{A})^{\otimes n}$, and each retained PBW basis element
tensored with an NBC form belongs to a
unique type assignment via its charge and derivative degree at each
position.

\emph{Count.} $|t_1| = 2$ and $|t_i| = 3$ for $i > 1$ gives
$\operatorname{rank}(\bar{B}^{\mathrm{lw},n}) = 2 \cdot 3^{n-1}$.
\end{proof}

\begin{remark}[\texorpdfstring{$\beta\gamma$}{beta-gamma} bar \emph{cohomology} closed form]\label{rem:betagamma-bar-cohomology-gf}
Complementing the chain-space rank formula (Theorem~\ref{thm:betagamma-bar-dim}),
the bar \emph{cohomology} dimensions have a proved closed form:
\[
\dim H^n(\barBgeom(\beta\gamma))
= [x^n]\sqrt{(1{+}x)/(1{-}3x)},
\]
The generating function shares the discriminant
$\Delta = (1{-}3x)(1{+}x)$
with $\widehat{\mathfrak{sl}}_2$ and Virasoro;
see Theorem~\ref{thm:ds-bar-gf-discriminant}.
\end{remark}

\subsection{\texorpdfstring{Shadow obstruction tower of the $\beta\gamma$ system}{Shadow obstruction tower of the beta-gamma system}}
\label{sec:betagamma-shadow-tower-free}
\index{shadow obstruction tower!beta-gamma@$\beta\gamma$!full tower}

The $\beta\gamma$ system is the bosonic free-field archetype outside
the Gaussian class; its $bc$ Koszul dual has the same finite depth in
the opposite statistical sector. The full shadow obstruction tower is controlled by the
interaction between two primary directions in the modular cyclic
deformation complex (Definition~\ref{def:modular-cyclic-deformation-complex}),
and terminates at degree~$4$ on the standard conformal-weight family
$\beta\gamma_\lambda$, where
$S_2 = 6\lambda^2 - 6\lambda + 1$, $S_3 = 0$, $S_4 = -5/12$, and
$S_r = 0$ for $r \ge 5$. The internal weight-changing line is
shadow-trivial, so the class-C witness is not a one-dimensional slice.

\subsubsection{Multi-channel structure}
\label{sec:betagamma-multi-channel}

\begin{proposition}[$\beta\gamma$ deformation complex;
\ClaimStatusProvedHere]
\label{prop:betagamma-deformation-channels}
\index{beta-gamma system@$\beta\gamma$ system!deformation channels}
The modular cyclic deformation complex of the $\beta\gamma$ system at
conformal weight~$\lambda$ has (at least) two primary directions:
\begin{enumerate}[label=\textup{(\roman*)}]
\item \emph{The stress tensor line}
 (T-line): the one-dimensional slice corresponding to
 the Virasoro subalgebra at central charge
 $c = c_{\beta\gamma}(\lambda) = 2(6\lambda^2 - 6\lambda + 1)$.
 On this line the shadow data coincides with that of
 $\mathrm{Vir}_{c}$.
\item \emph{The weight-changing line}: the deformation that shifts the
 conformal weight $\lambda \mapsto \lambda + \epsilon$. On this line
 all transferred higher brackets vanish
 ($\ell_n^{\mathrm{tr}}|_{L} = 0$ for $n \geq 2$),
 so the shadow obstruction tower is purely quadratic by rank-one abelian rigidity
 (Theorem~\textup{\ref{thm:betagamma-rank-one-rigidity}}).
\end{enumerate}
The $\beta\gamma$ system is \emph{not} a scalar-lane algebra:
the two directions carry independent shadow data, and the
global tower depends on their interaction.
\end{proposition}

\begin{proof}
Direction~(i): the stress tensor $T(z)$ of the $\beta\gamma$ system
generates a Virasoro subalgebra at central charge
$c = 2(6\lambda^2 - 6\lambda + 1)$. The shadow obstruction tower restricted to
this one-dimensional slice is the Virasoro shadow obstruction tower at this value
of~$c$, with $S_3 = 2$ (the universal Virasoro cubic shadow) and
$S_4 = 10/(c(5c+22))$ (the quartic contact invariant). Direction~(ii):
the weight-changing deformation class $\eta$ generates a one-dimensional
slice on which the beta-gamma OPE contributes no higher-order poles
beyond the simple pole $\beta(z)\gamma(w) \sim 1/(z{-}w)$. All
transferred higher brackets $\ell_n^{\mathrm{tr}}$ for $n \geq 2$
vanish on this slice, placing it in the hypothesis of
Theorem~\ref{thm:betagamma-rank-one-rigidity}.\qedhere
\end{proof}

\subsubsection{Shadow data on each line}

\begin{proposition}[$\beta\gamma$ shadow obstruction tower: T-line data;
\ClaimStatusConditional]
\label{prop:betagamma-T-line-shadows}
\index{beta-gamma system@$\beta\gamma$ system!T-line shadow data}
On the stress tensor line, the shadow obstruction tower of $\beta\gamma$ at
weight~$\lambda$ is
\begin{alignat}{2}
 S_2 &= \kappa(\beta\gamma) &&= 6\lambda^2 - 6\lambda + 1,
 \label{eq:bg-S2-Tline} \\
 S_3 &= \alpha &&= 2
 \quad \textup{(independent of~$\lambda$)},
 \label{eq:bg-S3-Tline} \\
 S_4 &= \frac{10}{c(5c+22)}
 &&= \frac{10}{2(6\lambda^2{-}6\lambda{+}1)
 \bigl(10(6\lambda^2{-}6\lambda{+}1)+22\bigr)},
 \label{eq:bg-S4-Tline}
\end{alignat}
and $S_r \neq 0$ for all $r \geq 2$.
The restriction to the T-line has infinite shadow depth
\textup{(}class~M on the one-dimensional slice\textup{)}.
\end{proposition}

\begin{proof}
The shadow invariants on the T-line are the Virasoro shadow invariants
at central charge $c = c_{\beta\gamma}(\lambda)$: $S_2 = \kappa = c/2$,
$S_3 = 2$ from the universal Virasoro cubic (the coefficient of
$z^{-3}$ in the $r$-matrix, which is one pole order below the $z^{-4}$
OPE pole by pole absorption), and $S_4 = Q^{\mathrm{contact}}_{\mathrm{Vir}}(c) =
10/(c(5c+22))$. The critical discriminant
$\Delta = 8\kappa S_4 = 40/(5c+22)$ is nonzero for all real~$\lambda$
(the denominator $5c+22 = 60\lambda^2 - 60\lambda + 32$ has negative
discriminant $-4080 < 0$), so the T-line tower does not terminate
(Theorem~\ref{thm:single-line-dichotomy}).\qedhere
\end{proof}

\begin{proposition}[$\beta\gamma$ shadow obstruction tower: weight-changing line;
\ClaimStatusConditional]
\label{prop:betagamma-weight-line-shadows}
\index{beta-gamma system@$\beta\gamma$ system!weight-changing line}
On the pure weight-changing line, the shadow obstruction tower vanishes at all degrees:
\[
 S_r\big|_{\mathrm{wc}} = 0
 \qquad (r \geq 2).
\]
This statement concerns only the one-dimensional slice generated by the
weight-changing class. It does not describe either the Virasoro T-line
of Proposition~\ref{prop:betagamma-T-line-shadows} or the standard
conformal-weight family that realizes the class-$\mathbf{C}$ witness in
Theorem~\ref{thm:betagamma-global-depth}.
\end{proposition}

\begin{proof}
The weight-changing line satisfies the hypothesis of rank-one abelian
rigidity (Theorem~\ref{thm:betagamma-rank-one-rigidity}): all
transferred higher brackets $\ell_n^{\mathrm{tr}}|_L = 0$ for
$n \geq 2$. The shadow obstruction tower is therefore purely quadratic, and since
$\kappa|_{\mathrm{wc}} = 0$ (the weight-changing deformation does not
produce genus-$0$ curvature), the entire tower
vanishes.\qedhere
\end{proof}

\subsubsection{Stratum separation and the quintic vanishing}

\begin{theorem}[$\beta\gamma$ global shadow depth; \ClaimStatusConditional]
\label{thm:betagamma-global-depth}
\index{beta-gamma system@$\beta\gamma$ system!standard conformal-weight family}
\index{beta-gamma system@$\beta\gamma$ system!shadow depth}
The standard conformal-weight family $\beta\gamma_\lambda$ has
global shadow depth $r_{\max} = 4$
\textup{(}class~C, the contact/quartic archetype\textup{)} with
\[
 S_2 = 6\lambda^2 - 6\lambda + 1,\qquad
 S_3 = 0,\qquad
 S_4 = -5/12,\qquad
 S_r = 0 \quad (r \geq 5).
\]
The internal weight-changing line has $S_r|_{\mathrm{wc}} = 0$ for all
$r \geq 2$, while the separately recomputed one-dimensional T-line tower
is class~M. Hence neither one-dimensional slice is the class-C witness
by itself.
\end{theorem}

\begin{proof}
This is the class-$\mathbf{C}$ clause of
Proposition~\ref{prop:depth-gap-trichotomy} and of the shadow
archetype classification in
Chapter~\ref{chap:higher-genus-modular-koszul}. The final sentence
combines Proposition~\ref{prop:betagamma-weight-line-shadows} with
Proposition~\ref{prop:betagamma-T-line-shadows}: the internal
weight-changing slice is trivial, while the T-line is class~M.\qedhere
\end{proof}

\begin{remark}[Why neither internal slice is the class-C witness]
\label{rem:betagamma-stratum-separation-mechanism}
The corrected class-$\mathbf{C}$ witness is the standard
conformal-weight family $\beta\gamma_\lambda$, not either internal
one-dimensional slice. The weight-changing line is trivial by
rank-one abelian rigidity, whereas the T-line is the Virasoro
recomputation at central charge
$c_{\beta\gamma}(\lambda) = 2(6\lambda^2 - 6\lambda + 1)$ and is
therefore class~M. The family-level shadow data
$S_2 = 6\lambda^2 - 6\lambda + 1$, $S_3 = 0$,
$S_4 = -5/12$, $S_r = 0$ for $r \ge 5$ is what isolates the contact
archetype in the standard landscape.
\end{remark}

\subsubsection{Shadow metric and growth rate on the T-line}

\begin{proposition}[$\beta\gamma$ shadow metric; \ClaimStatusConditional]
\label{prop:betagamma-shadow-metric}
\index{shadow metric!beta-gamma@$\beta\gamma$}
On the T-line, the shadow metric
\textup{(}Definition~\textup{\ref{def:shadow-metric})} is
\begin{equation}\label{eq:bg-shadow-metric}
 Q_L(t)
 = (2\kappa + 6t)^2 + 2\Delta\, t^2,
\end{equation}
where $\kappa = 6\lambda^2 - 6\lambda + 1$, the cubic shadow
$\alpha = S_3 = 2$, and the critical discriminant is
\begin{equation}\label{eq:bg-critical-discriminant}
 \Delta
 = 8\kappa\, S_4
 = \frac{40}{5c + 22}
 = \frac{40}{60\lambda^2 - 60\lambda + 32}.
\end{equation}
The discriminant $\Delta \neq 0$ for all real~$\lambda$, confirming
that $Q_L$ is irreducible quadratic and the T-line slice is
class~M.
\end{proposition}

\begin{proof}
Substituting $\alpha = 2$ and writing $d_c=5c+22$, so that
$S_4=10/(c\,d_c)$, into the shadow
metric formula $Q_L(t) = (2\kappa + 3\alpha t)^2 + 2\Delta t^2$
(Theorem~\ref{thm:riccati-algebraicity}) gives
$(2\kappa + 6t)^2 + 2\Delta t^2$ with
$\Delta = 80\kappa/(c\,d_c)=40/d_c$
(using $\kappa = c/2$). The denominator
$d_c=5c+22=60\lambda^2-60\lambda+32$
has discriminant $3600 - 4 \cdot 60 \cdot 32 = -4080 < 0$, so it has
no real roots.\qedhere
\end{proof}

\subsubsection{Evaluation at special weights}

\begin{computation}[$\beta\gamma$ shadow obstruction tower: special weight table;
\ClaimStatusProvedHere]
\label{comp:betagamma-shadow-weights}
\index{beta-gamma system@$\beta\gamma$ system!shadow invariants at special weights}
\begin{center}
\small
\renewcommand{\arraystretch}{1.35}
\begin{tabular}{l c c c c c c l}
\toprule
$\lambda$ & $c$ & $\kappa$ & $S_3$ & $S_4$ & $\Delta$ & $\rho$
 & \emph{Description} \\
\midrule
$0$ & $2$ & $1$ & $2$ & $\frac{5}{32}$ & $\frac{5}{4}$ & $\approx 3.1$
 & $\beta =$ function, $\gamma =$ diff. \\[2pt]
$\frac{1}{2}$ & $-1$ & $-\frac{1}{2}$ & $2$ & $-\frac{10}{17}$ & $\frac{40}{17}$ & $\approx 6.4$
 & symplectic boson \\[2pt]
$1$ & $2$ & $1$ & $2$ & $\frac{5}{32}$ & $\frac{5}{4}$ & $\approx 3.1$
 & $\beta =$ diff., $\gamma =$ function \\[2pt]
$2$ & $26$ & $13$ & $2$ & $\frac{5}{1976}$ & $\frac{5}{19}$ & $\approx 0.23$
 & quad.\ diff.\ / vector field \\
\bottomrule
\end{tabular}
\end{center}
Here $S_3$ and $S_4$ are the T-line shadow invariants,
$\Delta = 40/(5c+22)$ is the critical discriminant, and
$\rho$ is the shadow growth rate on the T-line
(Definition~\textup{\ref{def:shadow-growth-rate}}).
The columns $\lambda = 0$ and $\lambda = 1$ coincide by
weight symmetry ($\kappa(\lambda) = \kappa(1{-}\lambda)$,
Proposition~\textup{\ref{prop:betagamma-weight-symmetry}}).
The point $\lambda = 2$ gives the bosonic string pair
$(\beta\gamma, bc)$ at $(c, c') = (26, -26)$.
\end{computation}

\subsubsection{Weight symmetry and Koszul duality}

\begin{proposition}[Weight symmetry $\neq$ Koszul duality;
\ClaimStatusProvedHere]
\label{prop:betagamma-weight-symmetry}
\index{beta-gamma system@$\beta\gamma$ system!weight symmetry}
\index{Koszul duality!beta-gamma weight symmetry}
The modular characteristic satisfies
$\kappa(\beta\gamma_\lambda) = \kappa(\beta\gamma_{1-\lambda})$
for all~$\lambda$. This weight symmetry is \emph{distinct} from
Koszul duality.
\end{proposition}

\begin{proof}
The formula $\kappa(\lambda) = 6\lambda^2 - 6\lambda + 1$ is invariant
under $\lambda \mapsto 1 - \lambda$, since
$6(1{-}\lambda)^2 - 6(1{-}\lambda) + 1
= 6\lambda^2 - 6\lambda + 1$.
Weight symmetry exchanges the roles of $\beta$ and $\gamma$ (weights
$\lambda \leftrightarrow 1{-}\lambda$) within the same
\emph{bosonic} algebra. Koszul duality, by contrast, maps
$\beta\gamma_\lambda$ to $bc_\lambda$
(Theorem~\ref{thm:betagamma-bc-koszul}): same conformal weight, opposite
statistics. Concretely,
$\kappa(\beta\gamma_\lambda) + \kappa(bc_\lambda) = 0$ \textup{(}free-field pair\textup{)}
(Koszul complementarity) while
$\kappa(\beta\gamma_\lambda) = \kappa(\beta\gamma_{1-\lambda})$
(weight symmetry); these are independent identities.\qedhere
\end{proof}

\subsubsection{Mumford isomorphism connection}

\begin{remark}[$\beta\gamma$ modular characteristic and the Mumford isomorphism]
\label{rem:betagamma-mumford}
\index{Mumford isomorphism!beta-gamma connection@$\beta\gamma$ connection}
The Mumford exponent $e(n) = 6n^2 - 6n + 1$ satisfies
$e(\lambda) + e(1{-}\lambda) = c_{\beta\gamma}(\lambda)$ and
$\kappa(\beta\gamma_\lambda)
= \tfrac{1}{2}\bigl(e(\lambda) + e(1{-}\lambda)\bigr)$.
At integer~$\lambda = n$, the exponent $e(n)$ governs the
isomorphism
$\det R\Gamma(\Sigma_g, K_{\Sigma}^{\otimes n})
\cong \lambda_1^{\otimes e(n)}$,
so the modular characteristic of $\beta\gamma$ at integer weight
is the half-sum of the two Mumford exponents corresponding to
its two generators.
At the special values $e(0) = e(1) = 1$, $e(2) = 13$, $e(3) = 37$;
the identity $e(\lambda) = \kappa(\lambda)$ holds because
$e(\lambda) = 6\lambda^2 - 6\lambda + 1 = \kappa(\lambda)$.
\end{remark}

\medskip
\noindent\emph{Computational verification:}
\texttt{compute/lib/betagamma\_shadow\_full.py}
and
\texttt{compute/tests/test\_betagamma\_shadow\_full.py}
(103~tests).



\subsection{Heisenberg bar complex}\label{sec:heisenberg-bar-complex-sec2}

\begin{remark}[Shadow archetype: Gaussian]
\label{rem:heisenberg-shadow-archetype}
\index{Heisenberg algebra!shadow archetype}
Shadow depth $r_{\max} = 2$ (class~G).
Modular characteristic $\kappa(\cH_k) = k$.
Central charge $c = 1$.
Shadow obstruction tower terminates at the scalar level:
$\Theta^{\leq r} = k \cdot \eta \otimes \Lambda$ for all $r \geq 2$.
No cubic, quartic, or higher shadows exist.
The modular convolution algebra is strictly formal:
all transferred $L_\infty$ brackets vanish ($\ell_k^{\mathrm{tr}} = 0$
for $k \geq 3$).
The Heisenberg is the archetype of class~G: the simplest curved
algebra with nontrivial genus tower ($F_g = k \cdot \lambda_g^{\mathrm{FP}}$
\textup{(}UNIFORM-WEIGHT\textup{)})
but trivial nonlinear shadow structure.
\end{remark}

The Heisenberg algebra $\mathcal{H}_k$ (Chapter~\ref{ch:heisenberg-frame})
has no field-valued simple-pole bracket: the double-pole OPE
$\alpha(z)\alpha(w) \sim k(z-w)^{-2}$ kills the bracket differential
and produces the level central two-cochain. The corresponding
bar-side curvature is $+k\,\omega_{\mathrm{cen}}$; after
Verdier-linear Koszul duality the dual commutative algebra carries
the opposite curvature $-k\,\omega_{\mathrm{cen}}$.
The Koszul dual is the commutative chiral
algebra $\mathrm{Sym}^{\mathrm{ch}}(V^*)$, not a level-inverted
Heisenberg.

\subsubsection{Bar complex computation}
\label{sec:heisenberg-bar-complex}

\begin{theorem}[Heisenberg bar complex at genus 0; \ClaimStatusProvedHere]\label{thm:heisenberg-bar}
For $\mathcal{H}_k$ on $\mathbb{P}^1$:
\[
H^n(\bar{B}_{\text{geom}}(\mathcal{H}_k)) =
\begin{cases}
\mathbb{C} & n = 0 \\
0 & n = 1 \\
\mathbb{C} \cdot \omega_{\mathrm{cen},k} & n = 2 \\
0 & n > 2
\end{cases}
\]
where $\omega_{\mathrm{cen},k}$ is the level central class. For genus
$g \geq 1$, additional classes arise from $H^1(\Sigma_g)$; see
Theorem~\ref{thm:heisenberg-higher-genus}.
\end{theorem}

\begin{proof}
The degree-by-degree computation is the frame calculation of
Theorem~\ref{thm:frame-heisenberg-bar}:
the double-pole-only OPE kills the degree-$1$ differential,
the degree-$2$ triple-collision residue produces the level central
class $\omega_{\mathrm{cen},k}$, and degrees $\geq 3$ vanish by
dimension counting.
See Theorem~\ref{thm:heisenberg-bar-complete} for explicit formulas
and $d^2 = 0$ verification.

\begin{lemma}[Orientation consistency; \ClaimStatusProvedHere]\label{lem:orientation-freefields}
For the Fulton--MacPherson compactification $\overline{C}_{n+1}(X)$, the orientation on codimension-2 strata satisfies:
$\text{or}_{D_{ijk}} = \text{or}_{D_{ij}} \wedge \text{or}_{D_{jk}} = -\text{or}_{D_{ik}} \wedge \text{or}_{D_{jk}}$
\end{lemma}

\begin{proof}
In blow-up coordinates near $D_{ijk}$, let $\epsilon_{ij} = |z_i - z_j|$ and $\theta_{ij} = \arg(z_i - z_j)$. The blow-up of $\Delta_{ij}$ followed by $\Delta_{jk}$ gives coordinates:
\begin{align}
z_i &= u + \frac{\epsilon_{ij}}{2}e^{i\theta_{ij}} + \frac{\epsilon_{ijk}}{4}e^{i\phi_i}\\
z_j &= u - \frac{\epsilon_{ij}}{2}e^{i\theta_{ij}} + \frac{\epsilon_{ijk}}{4}e^{i\phi_j}\\
z_k &= u + \frac{\epsilon_{ijk}}{4}e^{i\phi_k}
\end{align}
where $\epsilon_{ijk}$ measures the scale of the triple collision. The orientation form is:
$\text{or}_{D_{ijk}} = d\epsilon_{ij} \wedge d\theta_{ij} \wedge d\epsilon_{jk} \wedge d\theta_{jk} \wedge \text{sgn}(\sigma)$
where $\sigma \in S_3$ is the permutation relating different blow-up orders. Computing the Jacobian:
$J = \frac{\partial(\epsilon_{ij}, \theta_{ij}, \epsilon_{jk}, \theta_{jk})}{\partial(\epsilon_{ik}, \theta_{ik}, \epsilon_{jk}, \theta_{jk})} = -1$
This gives the required sign relation, ensuring consistency of orientation across all strata.
\end{proof}

Stokes' theorem on $\overline{C}_{n+1}(X)$ as a manifold with corners is then rigorously justified: the boundary operator squares to zero because the orientation signs from different paths to codimension-2 strata cancel (Lemma~\ref{lem:orientation-freefields}).
\end{proof}


\subsubsection{Central terms and curved structure}

The Heisenberg bar complex carries a curved $A_\infty$ structure
(Definition~\ref{def:curved-ainfty-hg}), with curvature, convergence,
and monodromy finiteness guaranteed by the general theory of
Chapter~\ref{chap:higher-genus}.

\begin{remark}[Physical meaning of curvature]
\index{central charge!free fields}
The curvature class is the homological manifestation of the
Heisenberg central extension. In the rank-one bar-side convention it
is
\[
m_0=k\,\omega_{\mathrm{cen}},
\]
where $\omega_{\mathrm{cen}}$ is the central two-cochain dual to the
pairing $\langle J,J\rangle=k$. The Verdier-linear Koszul dual
algebra uses the opposite sign $m_0^!=-k\,\omega_{\mathrm{cen}}$.
This is a level statement, not the central charge $c=1$ of the
Heisenberg vertex algebra.
\end{remark}

\begin{remark}[Energy-momentum tensor]
\index{energy-momentum tensor!Heisenberg}
The stress tensor
\[
T=\frac{1}{2k}\normord{\alpha(z)\alpha(z)}
\]
has central charge~$1$, while the OPE
$\alpha(z)\alpha(w)\sim k(z-w)^{-2}$ records the level. The
normal-ordering point-splitting creates the central two-cochain
$\omega_{\mathrm{cen}}$ on the collision diagonal; multiplying it by
$k$ gives the curved cobar term.
\end{remark}

\begin{theorem}[Heisenberg curved cobar structure; \ClaimStatusProvedHere]
\label{thm:heisenberg-curved-structure}
The ordered bar complex of the Heisenberg algebra $\mathcal{H}_k$
has ordinary square-zero differential. Its curved cobar model records
the level curvature, and the Verdier-linear Koszul-dual algebra
carries the opposite scalar:
\begin{enumerate}
\item Bar-side curvature: $m_0^{\mathrm{bar}} = k\,\omega_{\mathrm{cen}}$, where
 $\omega_{\mathrm{cen}}$ is the central collision two-cochain.
\item Koszul-dual curvature:
 $m_0^! = -k\,\omega_{\mathrm{cen}}$ on
 $\mathcal H_k^!=(\mathrm{Sym}^{\mathrm{ch}}(V^*),m_0^!)$.
\item Binary: $m_2(J \otimes J) = 0$ after removing the scalar
 central collision term.
\item Curved relation: $m_1(m_0^{\mathrm{bar}}) = 0$ and
 $m_1(m_0^!)=0$ (the central class is closed)
\item Higher: $m_k = 0$ for $k \geq 3$
\end{enumerate}
\end{theorem}

\begin{proof}
The OPE $\alpha(z)\alpha(w)=k(z-w)^{-2}$ has no simple-pole
field-valued bracket, so the bracket part of the bar differential
vanishes on degree~$1$. The double pole is not discarded: with the
trace-form collision convention, the Heisenberg scalar is the level
$k$, not the Virasoro scalar \(c/2\). It produces the ordered residue
$r_{\cH_k}^{\mathrm{coll}}(z)=k/z$ and the central two-cochain
$\omega_{\mathrm{cen}}$.

At triple collision the two boundary orders give opposite Arnold
orientation signs and leave the scalar central class:
\[
\operatorname{Res}_{D_{123}}
\bigl[J_1J_2J_3\otimes\eta_{12}\wedge\eta_{23}\bigr]
=k\,\omega_{\mathrm{cen}}.
\]
This is the bar-side curved cobar curvature
$m_0^{\mathrm{bar}}=k\,\omega_{\mathrm{cen}}$. Verdier-linear
duality reverses the scalar and gives the Koszul-dual curvature
$m_0^!=-k\,\omega_{\mathrm{cen}}$ on
$\mathrm{Sym}^{\mathrm{ch}}(V^*)$.
The genus-$0$ bar differential itself still satisfies $d_{\bar B}^2=0$.

The curved $A_\infty$ relation at $n = 0$ gives
$m_1(m_0^{\mathrm{bar}}) = 0$ (curvature is closed), and at $n = 1$:
\[m_1^2(a) = m_2(m_0^{\mathrm{bar}}, a) - m_2(a, m_0^{\mathrm{bar}}) = [m_0^{\mathrm{bar}}, a]_{m_2}\]

Since $m_0^{\mathrm{bar}}=k\,\omega_{\mathrm{cen}}$ is central,
$[m_0^{\mathrm{bar}},a]=0$ and so $m_1^2=0$, consistent with
$\dzero^2=0$ on the genus-$0$ bar complex. The same centrality
argument applies after the sign reversal on the Koszul-dual side.
\end{proof}

% [Duplicate content removed]


\subsection{Lattice vertex operator algebras}

For an even lattice $L$ with bilinear form $(\cdot, \cdot)$:

\subsubsection{Setup}

\begin{definition}[Lattice VOA]
The lattice vertex algebra $V_L$ has vertex operators $e^\alpha$ for $\alpha \in L$ with:
\[
e^\alpha(z)e^\beta(w) \sim (z-w)^{(\alpha,\beta)} e^{\alpha+\beta}(w) + \cdots
\]
Conformal weight: $h_{e^\alpha} = \frac{(\alpha,\alpha)}{2}$.
\end{definition}

\subsubsection{Bar complex structure}

\begin{theorem}[Lattice VOA bar complex; \ClaimStatusProvedHere]
\label{thm:lattice-voa-bar}
The bar complex $\bar{B}_{\text{geom}}(V_L)$ has:
\begin{enumerate}
\item Grading by total lattice degree: $\sum_i \alpha_i \in L$
\item Differential preserves lattice grading
\item Simple poles occur only when $(\alpha_i, \alpha_j) = -1$
\end{enumerate}
These are lattice-graded residue channels. They do not change the
scalar shadow census: the oscillator current sector of an even
lattice VOA remains Gaussian with $r_{\max}=2$.
\end{theorem}

\begin{proof}
An element in degree $n$:
\[
e^{\alpha_1}(z_1) \otimes \cdots \otimes e^{\alpha_{n+1}}(z_{n+1}) \otimes \omega
\]
has lattice degree $\alpha_1 + \cdots + \alpha_{n+1}$.

The differential:
\[
d_{\text{fact}} = \sum_{(\alpha_i,\alpha_j)=-1} \text{Res}_{D_{ij}}\left[e^{\alpha_i+\alpha_j} \otimes \eta_{ij} \wedge -\right]
\]
preserves the total lattice degree.

Only pairs with $(\alpha_i, \alpha_j) = -1$ contribute simple poles and hence nontrivial residues.
\end{proof}

\subsubsection{\texorpdfstring{Example: root lattice $A_2$}{Example: root lattice A-2}}

For the $A_2$ root lattice with simple roots $\alpha_1, \alpha_2$ and $(\alpha_1, \alpha_2) = -1$:

\begin{proposition}[\texorpdfstring{$A_2$}{A2} lattice computation; \ClaimStatusProvedHere]\label{prop:A2-lattice-bar}
Key residue channels:
\begin{align}
d(e^{\alpha_1} \otimes e^{\alpha_2} \otimes \eta_{12}) &= -e^{\alpha_1+\alpha_2} \\
\pi_{\mathrm{vac}}\,
d(e^{\alpha_1} \otimes e^{-\alpha_1-\alpha_2} \otimes e^{\alpha_2}
\otimes \eta_{12} \wedge \eta_{23}) &= e^0 = 1
\end{align}
Here $\pi_{\mathrm{vac}}$ is projection to the vacuum scalar in the
oscillator factor. The unprojected second residue has oscillator-current
terms; the displayed equality is the scalar projection. The residue
composites are equivariant for the $A_2$ Weyl group action.
\end{proposition}

\begin{proof}
For the $A_2$ root lattice, $(\alpha_1, \alpha_1) = (\alpha_2, \alpha_2) = 2$ and $(\alpha_1, \alpha_2) = -1$.

\emph{First differential.}
The bar complex element $e^{\alpha_1} \otimes e^{\alpha_2} \otimes \eta_{12}$ lives in bar degree 1 (two tensor factors, one form). The differential extracts the OPE residue:
\[
d(e^{\alpha_1} \otimes e^{\alpha_2} \otimes \eta_{12}) = \operatorname{Res}_{z_1 = z_2}\left[e^{\alpha_1}(z_1) e^{\alpha_2}(z_2)\right].
\]
The lattice VOA OPE gives $e^{\alpha_1}(z) e^{\alpha_2}(w) = \epsilon(\alpha_1, \alpha_2)(z-w)^{(\alpha_1,\alpha_2)} e^{\alpha_1+\alpha_2}(w) + \cdots = \epsilon(\alpha_1, \alpha_2)(z-w)^{-1} e^{\alpha_1+\alpha_2}(w) + \cdots$ where $\epsilon(\alpha_1, \alpha_2) = -1$ is the 2-cocycle sign (for the standard bimultiplicative choice on $A_2$). Taking the residue of the simple pole gives $d = -e^{\alpha_1 + \alpha_2}$.

\emph{Second differential.}
The element $e^{\alpha_1} \otimes e^{-\alpha_1-\alpha_2} \otimes e^{\alpha_2} \otimes \eta_{12} \wedge \eta_{23}$ lives in bar degree 2. The differential is the sum of residues at the two collision divisors $D_{12}$ and $D_{23}$. For $D_{12}$: $(\alpha_1, -\alpha_1 - \alpha_2) = -2 + 1 = -1$, so $e^{\alpha_1}(z_1) e^{-\alpha_1-\alpha_2}(z_2) \sim \epsilon \cdot (z_1-z_2)^{-1} e^{-\alpha_2}(z_2)$. The remaining collision with $e^{\alpha_2}$ has exponent $-2$ and therefore passes through the Heisenberg oscillator current before taking its scalar part. For $D_{23}$: $(-\alpha_1-\alpha_2, \alpha_2) = -(\alpha_1,\alpha_2) - (\alpha_2,\alpha_2) = 1-2 = -1$, yielding the symmetric channel. After accounting for signs from form orientations and the cocycle, the vacuum projection of the two-channel composite is $e^0 = 1$.

\emph{Weyl group action.}
The Weyl group $W(A_2) \cong S_3$ acts on the lattice, permuting the roots $\{\alpha_1, \alpha_2, -\alpha_1-\alpha_2\}$. The bar complex is $W$-equivariant: the simple reflections $s_1, s_2$ act on bar complex elements by permuting lattice labels and adjusting cocycle signs. The residue composites above realize the rank-two Weyl combinatorics on lattice labels. They are not additional scalar shadow generators and do not change the lattice row from class~G.
\end{proof}

\subsection{Genus 1 examples: elliptic bar complexes}

\subsubsection{Free fermion on the torus}

\begin{theorem}[Elliptic free fermion bar complex; \ClaimStatusProvedHere]\label{thm:elliptic-fermion-bar}
Fix a spin structure $\sigma$ on the elliptic curve $E_\tau$.
The zero-mode projection of the free-fermion elliptic bar complex has
\[
H^n(\bar{B}_{\text{ell},0}(\mathcal{F},\sigma)) = \begin{cases}
\mathbb{C} & n = 0 \\
\mathbb{C}^2 & n = 1 \\
\mathbb{C} \cdot \hat{c}_\sigma & n = 2 \\
0 & n > 2
\end{cases}
\]
where the degree-$1$ term is $H^1(E_\tau,\mathbb C)$ and
$\hat c_\sigma$ is the Pfaffian anomaly line for the chosen spin
structure. Spin structures label the Pfaffian lines; they are not an
additional vector-space summand in the bar cohomology.
\end{theorem}

\begin{proof}
The genus-$1$ propagator for the spin structure $\sigma$ is the
Szeg\H{o} kernel. Its local singular part is the same
$(z-w)^{-1}$ singularity as on $\mathbb P^1$, so the positive
reduced bar differential is unchanged locally. The new classes come
from harmonic one-forms on the elliptic curve:
$H^1(E_\tau,\mathbb C)\cong \mathbb C^2$.

The theta functions enter through the determinant line, not as
ordinary cohomology generators. In degree~$2$ the modular anomaly
gives the Pfaffian line
\[
\hat{c}_\sigma = \frac{c}{24}\,\omega_{\mathcal{M}_1}
\qquad (c=1/2),
\]
with vanishing Pfaffian for the odd spin structure because of the
fermion zero mode. No further zero-mode cohomology occurs on an
elliptic curve.
\end{proof}

\subsubsection{Heisenberg algebra on higher genus}

\begin{theorem}[Higher genus Heisenberg; \ClaimStatusProvedHere]\label{thm:heisenberg-higher-genus}
For $\mathcal{H}_k$ on $\Sigma_g$ ($g \geq 1$), let
$\bar{B}_{\mathrm{geom},0}^{(g)}(\mathcal H_k)$ denote the
zero-mode scalar quotient of the geometric bar complex: the quotient
which retains the vacuum, the harmonic one-form current modes, the
fundamental class of $\Sigma_g$, and the level central collision
class. Its cohomology is
\[
H^n(\bar{B}_{\mathrm{geom},0}^{(g)}(\mathcal{H}_k)) = \begin{cases}
\mathbb{C} & n = 0 \\
\mathbb{C}^{2g} & n = 1 \\
\mathbb{C} \oplus \mathbb{C} \cdot \omega_{\mathrm{cen},k}^{(g)} & n = 2 \\
0 & n > 2
\end{cases}
\]
Here $H^1 \cong H^1(\Sigma_g, \mathbb{C})$ reflects the $2g$ first-cohomology classes of~$\Sigma_g$ (realized as zero modes of the current $\alpha(z)$), $H^2$ contains the fundamental class $[\Sigma_g]$ and an additional level class $\omega_{\mathrm{cen},k}^{(g)}$ arising from the central extension, measuring the obstruction to trivializing the Heisenberg bundle over $\mathcal{M}_g$.
This is the scalar obstruction projection used by the genus tower; it
is not a computation of every configuration-space class in the full
geometric bar complex.
\end{theorem}

\begin{proof}
The projection factors the full configuration-space bar complex
through the de~Rham cohomology of $\Sigma_g$ and the central
collision class. The truncation in this quotient is therefore the
ordinary surface truncation, not a vanishing theorem for the full
configuration-space complex.

\emph{Degree~0.} The bar complex $\bar{B}^0 = \mathbb{C}$ (vacuum), and $d^0 = 0$, so $H^0 = \mathbb{C}$.

\emph{Degree~1.} The degree-1 bar complex is:
\[
\bar{B}^1 = \Gamma(\overline{C}_2(\Sigma_g),\, V \boxtimes V \otimes \Omega^1_{\log})
\]
where $V$ is the weight-1 space spanned by the current $\alpha(z)$. On the closed surface~$\Sigma_g$, the zero modes of $\alpha(z)$, i.e., the elements of $\ker(d: \bar{B}^1 \to \bar{B}^0)$ modulo $\mathrm{im}(d: \bar{B}^2 \to \bar{B}^1)$, are classified by $H^1(\Sigma_g, \mathbb{C})$.

Concretely, the current $\alpha(z) = \partial\phi(z)$ on $\Sigma_g$ has zero modes $\oint_{A_i} \alpha(z)\, dz$ and $\oint_{B_j} \alpha(z)\, dz$ for the $A$- and $B$-cycles of $\Sigma_g$ ($i,j = 1, \ldots, g$). By Hodge theory, these span a $2g$-dimensional space isomorphic to $H^1(\Sigma_g, \mathbb{C}) \cong H^{1,0}(\Sigma_g) \oplus H^{0,1}(\Sigma_g)$.

\emph{Degree~2.} The degree-2 bar differential involves the collision residue:
\[
d(\alpha(z_1) \otimes \alpha(z_2) \cdot \eta_{12}) = \mathrm{Res}_{z_1 = z_2}\!\left[\frac{k}{(z_1-z_2)^2}\, \eta_{12}\right] = k\,1
\]
On $\Sigma_g$, this residue maps to $H^2(\Sigma_g, \mathbb{C}) \cong \mathbb{C}$ (the fundamental class~$[\Sigma_g]$). The class $\omega_{\mathrm{cen},k}^{(g)} \in H^2$ arises as follows: the Heisenberg algebra at level~$k$ defines a line bundle $\mathcal{L}_k$ on $\mathcal{M}_g$ whose first Chern class is $c_1(\mathcal{L}_k) = k\,\lambda$, where $\lambda$ is the Hodge class. The class $\omega_{\mathrm{cen},k}^{(g)}$ is the image of this Chern class in the bar complex cohomology, and it is nonzero whenever $k \neq 0$.

\emph{Degree $n > 2$.} The zero-mode scalar quotient contains only
the cohomology of the surface and the central collision line. Since
$H^n(\Sigma_g,\mathbb C)=0$ for $n>2$, no further class survives in
this quotient. The full configuration-space bar complex may contain
classes outside this scalar projection; those classes are not used in
the genus-universality scalar obstruction.
\end{proof}



\section{Koszul duality for free fields}\label{sec:koszul-duality-free-fields}

The free-field bar complexes determine Koszul-dual algebras by
Verdier-linear duality on their bar cohomology coalgebras. The
algebra-level dualities exchange statistics (bosonic
$\leftrightarrow$ fermionic) and invert the propagator. The same
bar-coalgebra mechanism controls the module-theoretic and
twisted-sector comparisons for the free-field families.

\subsection{\texorpdfstring{$\beta\gamma$--$bc$}{beta-gamma--bc} orthogonality and duality}

\begin{proposition}[\texorpdfstring{$bc$}{bc}--\texorpdfstring{$\beta\gamma$}{beta-gamma} orthogonality; \ClaimStatusProvedHere]\label{prop:bc-betagamma-orthogonality}
The relations $R_{bc} \perp R_{\beta\gamma}$ under the residue pairing: the Clifford relation of the $bc$ system is orthogonal to the Weyl relation of the $\beta\gamma$ system.
\end{proposition}

\begin{proof}
The generators pair via the residue pairing $\langle b, \beta \rangle = 1$, $\langle c, \gamma \rangle = 1$, $\langle b, \gamma \rangle = 0$, $\langle c, \beta \rangle = 0$ (dual bases).

The quadratic relations are $R_{bc} = b \otimes c + c \otimes b$ (Clifford/symmetric) and $R_{\beta\gamma} = \beta \otimes \gamma - \gamma \otimes \beta$ (Weyl/antisymmetric).

Computing:
\begin{align*}
\langle R_{bc}, R_{\beta\gamma} \rangle &= \langle b \otimes c + c \otimes b, \; \beta \otimes \gamma - \gamma \otimes \beta \rangle \\
&= \langle b, \beta \rangle \langle c, \gamma \rangle - \langle b, \gamma \rangle \langle c, \beta \rangle \\
&\quad + \langle c, \beta \rangle \langle b, \gamma \rangle - \langle c, \gamma \rangle \langle b, \beta \rangle \\
&= (1)(1) - (0)(0) + (0)(0) - (1)(1) = 0.
\end{align*}

The cancellation is exact: the symmetric relation $(+)$ and
antisymmetric relation $(-)$ are orthogonal complements under the
residue pairing. This is the quadratic coalgebra datum whose
finite Verdier dual gives the Koszul duality
$(\beta\gamma)^! \cong bc$.
\end{proof}


\subsection{\texorpdfstring{$\beta\gamma$--$bc$}{beta-gamma--bc} Koszul duality}\label{sec:fermion-boson-koszul}
\index{free fermion!Koszul dual}

\begin{theorem}[\texorpdfstring{$\beta\gamma$}{beta-gamma}--\texorpdfstring{$bc$}{bc} Koszul duality; \ClaimStatusProvedHere]\label{thm:betagamma-bc-koszul}
The $\beta\gamma$ system (two bosonic generators with antisymmetric relation) and the $bc$ ghost system (two fermionic generators with symmetric relation) are Koszul dual:
\begin{equation}\label{eq:betagamma-bc-koszul}
(\beta\gamma)^! \cong bc \quad \text{and} \quad (bc)^! \cong \beta\gamma.
\end{equation}
At the level of bar complexes:
\[
H^*(\bar{B}_{\mathrm{geom}}(bc))=:bc^{\mathrm i},\qquad
H^*(\bar{B}_{\mathrm{geom}}(\beta\gamma))
=:(\beta\gamma)^{\mathrm i}.
\]
These are intrinsic bar-dual coalgebras. The algebra statements in
\eqref{eq:betagamma-bc-koszul} are obtained after continuous
finite-window Verdier duality:
\[
\mathbb D\bigl((\beta\gamma)^{\mathrm i}\bigr)\cong bc,\qquad
\mathbb D\bigl(bc^{\mathrm i}\bigr)\cong \beta\gamma.
\]
\end{theorem}

\begin{remark}
This is a \emph{two-generator} duality: $\beta\gamma$ has
$\dim V = 2$ (bosonic) and $bc$ has $\dim V^* = 2$ (fermionic). It is
not the \emph{single-generator} duality
$\Lambda^{\mathrm{ch}}(\psi) \leftrightarrow
\mathrm{Sym}^{\mathrm{ch}}(\gamma)$ between one free fermion and one
free boson.
\end{remark}

\begin{proof}
The key computation uses the Gui--Li--Zeng quadratic duality framework. The quadratic datum for $\beta\gamma$ has generators $N = \mathbb{C}\beta \oplus \mathbb{C}\gamma$ ($\dim N = 2$, bosonic) and relation $P = \mathrm{span}\{\beta \otimes \gamma - \gamma \otimes \beta\} \subset \Lambda^2(N)$ (antisymmetric).

The dual datum $(s^{-1}N^\vee\omega^{-1}, P^\perp)$ is computed via the residue pairing:
\[\langle P \otimes \omega_{X^2}, P^\perp \otimes s^2\omega_{X^2}\rangle = 0\]

Under this pairing, the antisymmetric relation dualizes to:
\[P^\perp = \mathrm{span}\{b \otimes c + c \otimes b\} \subset \mathrm{Sym}^2(N^*)\]
which is the anticommutation relation for the $bc$ ghost system (two fermionic generators $b = \beta^*$, $c = \gamma^*$).

Therefore
$\mathcal{A}(s^{-1}N^\vee\omega^{-1}, P^\perp)=bc$. This is the
post-Verdier algebra associated to the intrinsic bar-dual coalgebra,
hence $(\beta\gamma)^! = bc$.

The geometric manifestation is that Verdier duality on configuration spaces
exchanges symplectic pairings with Clifford pairings. Here ``symplectic''
means antisymmetric ($\beta\gamma - \gamma\beta = 1$), while ``Clifford'' means anticommuting ($bc + cb = 1$).
See Section~\ref{subsec:NAP-betagamma} for the full NAP perspective.
\end{proof}

\begin{remark}[The twisting morphism for free-field dualities]
\label{rem:free-field-twisting}
\index{twisting morphism!free fields}
The Koszul duality $(\beta\gamma)^! \cong bc$ is mediated by the
universal twisting morphism
$\tau\colon \bar{B}(\beta\gamma) \to \beta\gamma$,
which takes a concrete form for free fields.

For the $\beta\gamma$ system, $\tau$ acts on a single-insertion
bar element by returning the corresponding current:
$\tau([\beta^*]) = \beta$, $\tau([\gamma^*]) = \gamma$.
The Maurer--Cartan equation $d\tau + \tau \star \tau = 0$
encodes the free-field OPE: the convolution
$(\tau \star \tau)([\beta^* \otimes \gamma^* \otimes \eta_{12}])
= \operatorname{Res}_{z_1 = z_2}[\beta(z_1)\,\gamma(z_2)\,\eta_{12}]
= 1$
recovers the contraction $\beta(z)\gamma(w) \sim 1/(z-w)$,
and the MC equation states that this is the \emph{only}
obstruction to exactness of the bar complex.

Dually, the twisting morphism for the $bc$ system
$\tau_{bc}\colon \bar{B}(bc) \to bc$ satisfies the
\emph{same} MC equation with the Clifford
OPE $b(z)c(w) \sim 1/(z-w)$. Verdier duality on
$\overline{C}_2(X)$ exchanges the two: the propagator
$\eta_{12} = d\log(z_1 - z_2)$ is Verdier self-dual, so
$\mathbb{D}_{\mathrm{Ran}}(\tau_{\beta\gamma})
= \tau_{bc}$
(Theorem~\ref{thm:bar-cobar-isomorphism-main},
Theorem~A$_2$).
The Koszul duality $(\beta\gamma)^! \cong bc$ is the
\emph{statement that Verdier duality intertwines the two
twisting morphisms}. Acyclicity of the twisted tensor product
$K(\tau) = \beta\gamma \otimes_\tau \bar{B}(\beta\gamma)$
follows from the quadratic Koszulness of the antisymmetric
relation $\beta\gamma - \gamma\beta = 1$, confirming the
Koszul property via the fundamental theorem
(Theorem~\ref{thm:fundamental-twisting-morphisms}).
\end{remark}


\subsection{Derived completion and extended duality}

\begin{definition}[Derived \texorpdfstring{$\beta\gamma$}{beta-gamma}-\texorpdfstring{$bc$}{bc} datum]\label{def:derived-bg-bc}
A derived $\beta\gamma$-$bc$ datum is a $\mathbb Z$-graded
finite-type chiral complex
$(\mathcal B^\bullet,Q)$ whose degree-zero fields contain the
first-order pair $\beta,\gamma$ and the ghost pair $b,c$, whose
weight spaces are finite-dimensional, and whose BRST operator has
cohomological degree~$1$ and satisfies
\[
Q^2=0,\qquad
Q\bigl(a(z)b(w)\bigr)
= (Qa)(z)b(w)+(-1)^{|a|}a(z)(Qb)(w)
\]
on singular OPE coefficients. The datum is Verdier-dualizable only
after choosing a nondegenerate residue pairing between the finite
weight pieces. It is not an alternating chain
$\beta\gamma\to bc\to\beta'\gamma'$ without these nilpotence,
finite-type, and pairing conditions.
\end{definition}

\begin{remark}[Geometric origin]
For a holomorphic vector bundle $E$ on a curve, a first-order
$\beta\gamma$ system is attached to sections of $E$ and
$K\otimes E^\vee$; a $bc$ system is attached to the parity-reversed
cotangent fields. A BRST operator becomes the Dolbeault differential
only after the gauge symmetry, ghost sector, and anomaly cancellation
make $Q^2=0$. The geometric model therefore supplies examples of
Definition~\ref{def:derived-bg-bc}; it does not classify all derived
free-field extensions.
\end{remark}

\begin{conjecture}[BV-compatible derived free-field extension;
\ClaimStatusConjectured]\label{conj:extended-ferm-ghost}
Let $(\mathcal B^\bullet,Q)$ be a derived $\beta\gamma$-$bc$ datum
with finite-dimensional weight spaces and nondegenerate residue
pairing. If the bar spectral sequence of $\mathcal B^\bullet$
converges and the induced pairing identifies
$H^\bullet B(\mathcal B^\bullet)$ with its continuous Verdier dual,
then the Verdier-linear Koszul dual
\[
(\mathcal B^\bullet)^!
=\bigl(H^\bullet B(\mathcal B^\bullet)\bigr)^\vee
\]
carries a derived free-field model whose degree-shifted symplectic
pairing is the BV field-antifield pairing. This belongs to the
derived/super extension governed by
Conjecture~\ref{conj:v1-master-infinite-generator}, separate from the
$W_\infty$/Yangian comparison.
\end{conjecture}

\begin{remark}[Structural constraints on the derived extension]
The conjecture has four independent obligations:
\begin{enumerate}[label=\textup{(\roman*)}]
\item the jet-level BRST operator must be square-zero, including its
anomaly term;
\item the residue pairing must be nondegenerate on each finite weight
piece, so that continuous Verdier duality is defined;
\item the bar spectral sequence must be complete and convergent;
\item the BV symplectic form must agree with the Verdier pairing after
the Koszul shift.
\end{enumerate}
The $A_\infty$ structure on the dual side is recovered from
$\mathbb D_{\mathrm{Ran}}(\bar B(\mathcal B^\bullet))$. The separate
equivalence
$\Omega(\bar B(\mathcal B^\bullet))\simeq\mathcal B^\bullet$
is bar-cobar inversion: it recovers $\mathcal B^\bullet$ itself, not
the dual algebra.
\end{remark}

\begin{example}[BV interpretation]
When the four obligations above hold for a gauge-fixed free theory,
the BRST complex computes the gauge-invariant boundary states and the
Verdier-linear Koszul dual carries the shifted BV pairing. This is a
comparison between two chain complexes. It is not the assertion that
bar-cobar inversion is itself the BV field-antifield transform.
\end{example}

\begin{remark}[Derived extension boundary]\label{rem:extended-ferm-ghost-scope}
The derived/super extension proof consists of one functorial argument
assembling the jet bundle structure, stratified residue pairing, and
cobar reconstruction. It is not one of the $W_\infty$/Yangian
finite-window comparisons; it has no analogue of
$\mathcal{I}_N$ or $\Delta_{a,0}(N)$ in this chapter.
\end{remark}


\subsection{Module Koszul duality for \texorpdfstring{$\beta\gamma$--$bc$}{betagamma--bc}}
\label{sec:betagamma-bc-modules}
\index{beta-gamma system@$\beta\gamma$ system!module Koszul duality}
\index{bc system@$bc$ system!module Koszul duality}

The algebra-level duality $(\beta\gamma)^! \cong bc$
(Theorem~\ref{thm:betagamma-bc-koszul}) extends to module
categories. Both systems have natural families of modules
parameterized by their zero-mode eigenvalue.

\begin{definition}[\texorpdfstring{$\beta\gamma$}{beta-gamma} and \texorpdfstring{$bc$}{bc} Fock modules]
\label{def:bg-bc-fock}
\index{Fock module!beta-gamma@$\beta\gamma$|textbf}
\index{Fock module!bc@$bc$|textbf}
For $\lambda \in \mathbb{C}$:
\begin{enumerate}[label=\textup{(\roman*)}]
\item The \emph{$\beta\gamma$ Fock module} $M_\lambda^{\beta\gamma}$
 is generated by $|\lambda\rangle$ with
 $\gamma_0|\lambda\rangle = \lambda|\lambda\rangle$,
 $\beta_n|\lambda\rangle = \gamma_n|\lambda\rangle = 0$ for
 $n > 0$. Character (with $q = e^{2\pi i\tau}$,
 $y = e^{2\pi i z}$):
 \[
 \operatorname{ch}(M_\lambda^{\beta\gamma})(\tau, z)
 = \frac{y^{\lambda}}{\prod_{n \geq 1}(1 - q^n)^2}.
 \]
\item The \emph{$bc$ Fock module} $M_q^{bc}$ at ghost number
 $q \in \mathbb{Z}$ is generated by $|q\rangle$ with
 ghost number $J_0^{\mathrm{gh}}|q\rangle = q|q\rangle$,
 $b_n|q\rangle = c_n|q\rangle = 0$ for appropriate
 $n > 0$ (depending on ghost number vacuum).
 For general weights $(j, 1{-}j)$, the conformal weight is
 $h_q = -\tfrac{1}{2}q(q - (2j{-}1))$. For $(2, -1)$:
 \[
 \operatorname{ch}(M_q^{bc})
 = q^{h_q}\,\prod_{n \geq 1}(1 + q^n)^2,
 \qquad h_q = -\tfrac{1}{2}q(q-3).
 \]
\end{enumerate}
\end{definition}

\begin{proposition}[Module Koszul duality for \texorpdfstring{$\beta\gamma$}{beta-gamma}--\texorpdfstring{$bc$}{bc};
\ClaimStatusConditional]
\label{prop:bg-bc-module-kd}
\index{module Koszul duality!$\beta\gamma$--$bc$}
Under the Koszul duality $(\beta\gamma)^! \cong bc$, the
module Koszul duality functor
\textup{(}Theorem~\textup{\ref{thm:e1-module-koszul-duality})}
yields an equivalence
\begin{equation}\label{eq:bg-bc-module-equiv}
D(\mathrm{Mod}^{\Eone,\,\mathrm{compl}}_{\beta\gamma})
\;\simeq\;
D(\mathrm{Mod}^{\Eone,\,\mathrm{compl}}_{bc}).
\end{equation}
Under this equivalence:
\begin{enumerate}[label=\textup{(\roman*)}]
\item The $\beta\gamma$ Fock module $M_\lambda^{\beta\gamma}$
 maps to the $bc$ Fock module $M_{q(\lambda)}^{bc}$ with
 $q(\lambda) = \lambda$ \textup{(}ghost number
 $=$ $\gamma_0$-eigenvalue\textup{)}.
\item The symplectic pairing
 $\beta(z)\gamma(w) \sim (z-w)^{-1}$ on the $\beta\gamma$ side
 corresponds to the Clifford pairing
 $b(z)c(w) \sim (z-w)^{-1}$ on the $bc$ side: the
 symmetric/antisymmetric exchange of the algebra-level duality
 extends to modules.
\item The spectral flow
 $\sigma_n^{\beta\gamma}: \beta_m \mapsto \beta_{m+n}$,
 $\gamma_m \mapsto \gamma_{m-n}$ dualizes to
 $\sigma_n^{bc}: b_m \mapsto b_{m+n}$,
 $c_m \mapsto c_{m-n}$, exchanging twist sectors.
\end{enumerate}
\end{proposition}

\begin{proof}
Apply Theorem~\ref{thm:e1-module-koszul-duality} to the
Koszul pair $(\beta\gamma, bc)$. The hypotheses are satisfied:
$\beta\gamma$ is Koszul (the bar spectral sequence collapses
at $E_2$ by Theorem~\ref{thm:betagamma-bc-koszul}), and
both algebras have finite-type grading (finite-dimensional
weight spaces), so the bar coalgebra is graded-dualizable.

For (i): the bar resolution of $M_\lambda^{\beta\gamma}$
has the form
$\bar{B}_n(\beta\gamma, M_\lambda^{\beta\gamma})
= \Gamma(\overline{C}_{n+2}(X),\,
j_*j^*(\beta\gamma \boxtimes
\overline{\beta\gamma}^{\boxtimes n} \boxtimes
M_\lambda^{\beta\gamma}) \otimes \Omega^n_{\log})$.
The $\gamma_0$-eigenvalue $\lambda$ is preserved by the
bar differential (which acts on the
$z_i \neq z_j$ locus and does not affect the zero mode).
On the dual side, $\lambda$ becomes the ghost number $q$.

For (ii): the exchange of symmetric/antisymmetric pairings
is the quadratic duality of Proposition~\ref{prop:bc-betagamma-orthogonality},
which extends to modules by functoriality of the bar construction.

For (iii): follows from Example~\ref{ex:bc-spectral-flow} and
the functoriality of bar-cobar for algebra automorphisms.
\end{proof}

\begin{remark}[BRST complexes and module Koszul duality]
\label{rem:brst-module-kd}
\index{BRST!module Koszul duality}
The $bc$ ghosts provide a BRST resolution of gauge constraints only
on the anomaly-free surface where the total BRST charge is nilpotent.
Module Koszul duality for the $\beta\gamma$--$bc$ pair is a separate
finite-type Verdier statement: it applies the bar construction to a
module and then dualizes the resulting comodule. On proved comparison
surfaces these two complexes can be related by a chain map. Without
that map, one must not identify the BRST differential with the bar
differential, nor the BRST cohomology with the Koszul-dual module.
\end{remark}


\subsection{Heisenberg Koszul dual: symmetric algebra}

\begin{theorem}[Heisenberg Koszul dual; \ClaimStatusProvedHere]\label{thm:heisenberg-koszul-dual-early}
\index{Heisenberg algebra!Koszul dual}
The Heisenberg chiral algebra $\mathcal{H}_k$ at level $k$ has Koszul dual:
\[\mathcal{H}_k^! \simeq \bigl(\mathrm{Sym}^{\mathrm{ch}}(V^*),\, d = 0,\, m_0 = -k\,\omega\bigr)\]
where $\mathrm{Sym}^{\mathrm{ch}}(V^*)$ is the chiral symmetric (commutative) algebra on the dual generator space, and $\omega \in \mathrm{Sym}^2(V^*)$ is the element dual to the pairing $\langle \alpha, \alpha \rangle = k$. This is a \emph{curved} commutative chiral algebra in the sense of Positselski: the curvature $m_0 = -k\,\omega$ encodes the level, and $m_0 = 0$ only at level $k = 0$.

More explicitly, the bar complex decomposes as:
\begin{align}
\bar{B}^{\mathrm{ch}}(\mathcal{H}_k) &\simeq \mathrm{coLie}^{\mathrm{ch}}(V^*) \quad \text{(conilpotent chiral coalgebra)} \\
\Omega^{\mathrm{ch}}(\mathrm{coLie}^{\mathrm{ch}}(V^*)) &\simeq \bigl(\mathrm{Sym}^{\mathrm{ch}}(V^*),\, m_0 = -k\,\omega\bigr) \quad \text{(curved cobar)}
\end{align}
\end{theorem}

\begin{proof}
\emph{Proof~1: Via factorization envelopes.}
The Heisenberg chiral algebra is the factorization envelope of the abelian Lie$^*$ algebra $\mathfrak{h}_k = (\mathcal{O}_X, 0)$ with trivial bracket and pairing $\langle \alpha, \alpha \rangle = k$. This is a special case of Nishinaka's general construction~\cite{Nish26} (see also Vicedo~\cite{Vic25}): the enveloping vertex algebra $\Uvert(\mathfrak{h}_k)$ recovers $\cH_k$. By the Francis--Gaitsgory theory of chiral Lie$\leftrightarrow$Com duality \cite{FG12}, for an abelian chiral Lie algebra $\mathfrak{a}$:
\[\mathrm{Fact}(\mathfrak{a})^! \simeq C^*_{\mathrm{Lie}^*}(\mathfrak{a}) = \mathrm{Sym}^{\mathrm{ch}}(\mathfrak{a}^*[1])\]
Since $\mathfrak{h}_k$ is abelian, the Chevalley--Eilenberg complex is formal: $C^*_{\mathrm{Lie}^*}(\mathfrak{h}_k) = \mathrm{Sym}^{\mathrm{ch}}(\mathfrak{h}_k^*[1])$, which is the chiral symmetric algebra on the dual generators.

\medskip
\emph{Proof~2: Explicit bar complex computation}
(Theorem~\ref{thm:frame-heisenberg-bar}).
The double-pole-only OPE makes the coproduct on $\bar{B}^{\mathrm{ch}}(\mathcal{H}_k)$ cocommutative (the residue $\mathrm{Res}_{z_i = z_j}[\omega/(z_i - z_j)^2]$ is symmetric under $z_i \leftrightarrow z_j$). By Milnor--Moore, a conilpotent cocommutative coalgebra is a coLie coalgebra, whose cobar is $\mathrm{Sym}^{\mathrm{ch}}(V^*)$.

The level $k$ determines the curvature of the Koszul dual:
$m_0 = -k\,\omega \in \mathrm{Sym}^2(V^*)$. The bar differential
satisfies $d^2 = 0$; the curvature manifests instead in the
$A_\infty$ relation $m_1^2(a) = [m_0, a]_{m_2}$ on the cobar side.
\end{proof}

\begin{remark}[Level inversion vs Koszul duality]
Two distinct phenomena must be distinguished. \emph{Koszul duality}:
the bar coalgebra $\bar{B}(\mathcal{H}_k)$, dualized, gives
$\mathcal{H}_k^! = \mathrm{Sym}(V^*)$ with curvature $-k\omega$.
\emph{Level inversion}: over $\mathbb C$, the oscillator relation at
level $-k$ is obtained from level $k$ by rescaling $J\mapsto iJ$;
with real or unitary structures this is only an opposite-curvature
comparison. Level inversion does not compute $\mathcal H_k^!$.
\end{remark}

\begin{corollary}[Heisenberg module-comodule equivalence; \ClaimStatusProvedHere]
\label{cor:heisenberg-module-equivalence}
The bar--cobar adjunction restricts to an equivalence of derived categories
between \emph{complete} $\Eone$-chiral $\mathcal{H}_k$-modules and
\emph{conilpotent} $\Eone$-chiral $\bar{B}^{\mathrm{ch}}(\mathcal{H}_k)$-comodules:
\begin{equation}
D(\mathrm{Mod}^{\Eone,\,\mathrm{compl}}_{\mathcal{H}_k})
\;\simeq\;
D(\mathrm{CoMod}^{\Eone,\,\mathrm{conil}}_{\bar{B}^{\mathrm{ch}}(\mathcal{H}_k)}).
\end{equation}
Since $\mathcal{H}_k$ has finite-type grading (finite-dimensional weight spaces),
the coalgebra $\bar{B}^{\mathrm{ch}}(\mathcal{H}_k)$ is graded-dualizable with
linear dual $\mathcal{H}_k^! \simeq \mathrm{Sym}^{\mathrm{ch}}(V^*)$
(Theorem~\ref{thm:heisenberg-koszul-dual-early}). Under this dualization,
conilpotent comodules correspond to complete modules, yielding:
\begin{equation}
D(\mathrm{Mod}^{\Eone,\,\mathrm{compl}}_{\mathcal{H}_k})
\;\simeq\;
D(\mathrm{Mod}^{\Eone,\,\mathrm{compl}}_{\mathrm{Sym}^{\mathrm{ch}}(V^*)}).
\end{equation}
Under this equivalence, Ext groups are exchanged:
$\mathrm{Ext}^i_{\mathcal{H}_k}(L, M) \cong
\mathrm{Ext}^i_{\mathrm{Sym}^{\mathrm{ch}}(V^*)}(\mathrm{Ind}(M), \mathrm{Ind}(L))$.

Since $\mathrm{Sym}^{\mathrm{ch}}(V^*)$ is commutative, its complete module theory
is considerably more tractable than Heisenberg Fock module theory. This equivalence
``linearizes'' Heisenberg representation theory via Koszul duality.
\end{corollary}

\begin{remark}
The equivalence is between \emph{complete} module categories,
\emph{not} the full bounded derived categories $D^b$
The intermediate
coalgebra step is essential: modules become comodules via bar resolution,
and only under the finite-type dualizability of $\mathcal{H}_k$
do conilpotent comodules further correspond to complete modules over the
dual algebra $\mathrm{Sym}^{\mathrm{ch}}(V^*)$.
\end{remark}

\begin{proof}
Apply Theorem~\ref{thm:e1-module-koszul-duality} (module Koszul duality) to
$\mathcal{A} = \mathcal{H}_k$. The hypotheses are satisfied: $\mathcal{H}_k$
is Koszul (the bar spectral sequence collapses at $E_2$ by the computation above), and
the bar coalgebra is $\bar{B}^{\mathrm{ch}}(\mathcal{H}_k)$, whose
cohomology coalgebra $\mathcal{H}_k^i = H^*(\barB^{\mathrm{ch}}(\mathcal{H}_k))$
is the dual coalgebra; its graded linear dual is
$\mathcal{H}_k^! \simeq \mathrm{Sym}^{\mathrm{ch}}(V^*)$ by
Theorem~\ref{thm:heisenberg-koszul-dual-early}. The passage from comodules to
modules over the dual algebra uses the finite-type grading of $\mathcal{H}_k$
(finite-dimensional weight spaces ensure graded dualizability).
\end{proof}

\begin{computation}[Bar complex of the Koszul dual \texorpdfstring{$\mathrm{Sym}^{\mathrm{ch}}(V^*)$}{Symch(V^*)}]
\label{comp:bar-sym-ch}
\index{commutative chiral algebra!bar complex}
We compute $\bar{B}^{\mathrm{ch}}(\mathrm{Sym}^{\mathrm{ch}}(V^*))$
explicitly, verifying the Koszul duality inverse:
$\Omega(\bar{B}(\mathrm{Sym}^{\mathrm{ch}}(V^*)))
\simeq \mathrm{Sym}^{\mathrm{ch}}(V^*)$.

\emph{The bar complex.}
For a commutative chiral algebra $\mathcal{C} = \mathrm{Sym}^{\mathrm{ch}}(W)$
generated by a vector space~$W$, the operadic Koszul duality
$\mathsf{Com}^{\mathrm{ch},!} = \chirLie$ gives:
\[
\bar{B}^{\mathrm{ch}}(\mathrm{Sym}^{\mathrm{ch}}(W))
\;\simeq\;
\mathrm{coLie}^{\mathrm{ch}}(s^{-1}\bar{W})
\]
where $\bar{W} = W$ (no vacuum component in the augmentation ideal
of a free algebra) and $\mathrm{coLie}^{\mathrm{ch}}$ is the
conilpotent Lie chiral coalgebra cogenerated by~$s^{-1}W$.

\emph{Rank~$1$ case ($W = V^* \cong \mathbb{C}$).}
The free Lie algebra on one generator is abelian: $\mathrm{Lie}(1) = \mathbb{C}$
and $\mathrm{Lie}(n) = 0$ for $n \geq 2$
(there are no nontrivial brackets on a single element).
Therefore:
\begin{equation}\label{eq:bar-sym-rank1}
\dim \bar{B}^n(\mathrm{Sym}^{\mathrm{ch}}(V^*)) =
\begin{cases}
1 & n = 1, \\
0 & n \geq 2.
\end{cases}
\end{equation}
The bar complex is concentrated in degree~$1$, with
$\bar{B}^1 = s^{-1}V^*$. This is the bar complex of a
\emph{free} algebra; it detects only the generators and has
no higher bar cohomology.

\emph{General rank ($W = V^* \cong \mathbb{C}^m$).}
For $\dim V^* = m \geq 2$, the free Lie algebra dimensions
are given by the Witt formula (necklace counting):
\[
\dim \bar{B}^n(\mathrm{Sym}^{\mathrm{ch}}(V^*))
= \ell(n, m) = \frac{1}{n}\sum_{d \mid n} \mu(d)\, m^{n/d}
\]
where $\mu$ is the M\"obius function. For $m = 2$
(the $\beta\gamma$-system; see Theorem~\ref{thm:betagamma-bc-koszul}):
\[
\ell(1,2) = 2,\quad
\ell(2,2) = 1,\quad
\ell(3,2) = 2,\quad
\ell(4,2) = 3,\quad
\ell(5,2) = 6.
\]
These are the dimensions of the degree-$n$ component of the
free Lie algebra on two generators.

\emph{Koszul inversion check.}
Taking cobar of $\mathrm{coLie}^{\mathrm{ch}}(s^{-1}V^*)$
recovers $\mathrm{Sym}^{\mathrm{ch}}(V^*)$ by the
bar-cobar inversion theorem
(Theorem~\ref{thm:bar-cobar-inversion-qi}), confirming the
Koszul pair $(\mathcal{H}_k,\,
\mathrm{Sym}^{\mathrm{ch}}(V^*))$ is self-consistent.
\end{computation}

\subsection{Heisenberg module Koszul duality}

\begin{definition}[Fock module; \ClaimStatusDefinitional]\label{def:fock-module}
\index{Fock module|textbf}
For $\lambda \in \mathbb{C}$ and nonzero level~$k$, the \emph{Fock module}
$F_\lambda$ over $\mathcal{H}_k$ is generated by a vector
$|\lambda\rangle$ with:
\begin{itemize}
\item $J_0|\lambda\rangle = \lambda|\lambda\rangle$,\quad
 $J_n|\lambda\rangle = 0$ for $n > 0$;
\item $F_\lambda \cong \mathbb{C}[J_{-1}, J_{-2}, \ldots]$ as a
 vector space, with character
  $\operatorname{ch}(F_\lambda) = q^{\lambda^2/2k}\,
  \prod_{n \geq 1}(1-q^n)^{-1}$.
\end{itemize}
\end{definition}

\begin{lemma}[Fock module simplicity; \ClaimStatusProvedHere]
\label{lem:fock-module-simplicity}
For nonzero level~$k$, the Heisenberg Fock module $F_\lambda$ is simple
for every $\lambda\in\mathbb C$.
\end{lemma}

\begin{proof}
Use the oscillator relations
$[J_n,J_m]=n k\,\delta_{n+m,0}$ with $k\neq0$.
Let $0\neq v\in F_\lambda$ and write it in PBW monomials in the negative
modes $J_{-n}$. Choose a highest monomial in a lexicographic order
refining total oscillator degree. Applying the matching product of
positive modes gives a nonzero scalar multiple of $|\lambda\rangle$ plus
terms of lower order; descending induction puts $|\lambda\rangle$ in
the submodule generated by~$v$. The negative modes generate all of
$F_\lambda$ from $|\lambda\rangle$, so any nonzero submodule is the
whole module.
\end{proof}

\begin{proposition}[Fock module bar resolution; \ClaimStatusProvedHere]
\label{prop:fock-bar-resolution}
The geometric bar resolution
(Theorem~\ref{thm:geometric-bar-module}) of $F_\lambda$ over
$\mathcal{H}_k$ takes the form:
\begin{equation}\label{eq:fock-bar-resolution}
\overline{B}_n^{\mathrm{ch}}(\mathcal{H}_k, F_\lambda)
= \Gamma\bigl(\overline{C}_{n+2}(X),\;
j_*j^*(\mathcal{H}_k \boxtimes \overline{\mathcal{H}}_k^{\boxtimes n}
\boxtimes F_\lambda) \otimes \Omega^n_{\log}\bigr).
\end{equation}
The bar differential $d_{\mathrm{bar}} = \sum_{i<j} d_{ij}$ has the
following structure:
\begin{enumerate}[label=\textup{(\roman*)}]
\item Collisions $z_i = z_j$ between two algebra points
 ($1 \leq i < j \leq n$) produce the residue
 $k \cdot \mathrm{Res}_{z_i = z_j}[\omega/(z_i - z_j)^2]$,
 which is \emph{symmetric} in $z_i, z_j$ (as in
 Theorem~\textup{\ref{thm:heisenberg-koszul-dual-early})};
\item Collision of an algebra point $z_i$ with the module point
 $z_{n+1}$ produces the Heisenberg action
 $\mathrm{Res}_{z_i = z_{n+1}} \alpha(z_i) \cdot F_\lambda$;
\item Mixed terms vanish by the abelian OPE
 $\alpha(z)\alpha(w) \sim k/(z-w)^2$ (no simple pole).
\end{enumerate}
\end{proposition}

\begin{proof}
The bar complex definition is Theorem~\ref{thm:geometric-bar-module}
specialized to $\mathcal{A} = \mathcal{H}_k$,
$\mathcal{M} = F_\lambda$. The structure of the bar differential
follows from the Heisenberg OPE:

\emph{Algebra--algebra collisions.} For $1 \leq i < j \leq n$, the OPE
$\alpha(z_i)\alpha(z_j) \sim k/(z_i - z_j)^2$ contributes only a double pole.
The residue $\mathrm{Res}_{z_i = z_j}[\omega/(z_i - z_j)^2]
= \partial_{z_i}\omega|_{z_i = z_j}$ is symmetric in $z_i, z_j$
(as computed in Step~2--3 of the proof of
Theorem~\ref{thm:heisenberg-koszul-dual-early}).

\emph{Algebra--module collisions.} The OPE
$\alpha(z_i) \cdot |\lambda\rangle$ produces the Heisenberg action:
$J_0|\lambda\rangle = \lambda|\lambda\rangle$ from the
residue of the simple pole, and $J_{-m}|\lambda\rangle$
from higher poles. This is the standard module action,
realized geometrically by the residue along the collision divisor
$D_{i,n+1} = \{z_i = z_{n+1}\}$.

\emph{No simple poles between algebra points.} Since the Heisenberg OPE
has no simple pole, there are no ``mixed'' contributions where an
algebra collision produces a non-scalar result. Consequently the algebra part of the bar complex is purely
cocommutative, and the module point interacts with algebra points
only through the standard action.
\end{proof}

\begin{proposition}[Koszul dual module; \ClaimStatusConditional]
\label{prop:fock-koszul-dual}
Under the module Koszul duality equivalence of
Corollary~\textup{\ref{cor:heisenberg-module-equivalence}}, the Fock module
$F_\lambda$ over $\mathcal{H}_k$ corresponds to a module
$F_\lambda^!$ over
$\mathcal{H}_k^! = \mathrm{Sym}^{\mathrm{ch}}(V^*)$
characterized by:
\begin{enumerate}[label=\textup{(\roman*)}]
\item $F_\lambda^!$ is generated by a vector $|\lambda\rangle$ with
 $\alpha_0|\lambda\rangle = \lambda|\lambda\rangle$;
\item as a graded vector space,
 $F_\lambda^! \cong \mathbb{C}[\alpha_{-1}, \alpha_{-2}, \ldots]
 \cdot |\lambda\rangle$;
\item the $A_\infty$ module structure on $F_\lambda^!$ inherits the
 curvature $m_0 = -k\,\omega \in \mathrm{Sym}^2(V^*)$
 from the algebra level
 \textup{(}Theorem~\textup{\ref{thm:heisenberg-koszul-dual-early})},
 so the strict differential satisfies
 $m_1^2 = [m_0, -]_{m_2} \neq 0$ when $k \neq 0$.
\end{enumerate}
\end{proposition}

\begin{proof}
Apply the module Koszul duality functor of
Theorem~\ref{thm:e1-module-koszul-duality} to
$\mathcal{M} = F_\lambda$. The bar complex
$\bar{B}^{\mathrm{ch}}(\mathcal{H}_k, F_\lambda)$ is a comodule
over $\bar{B}^{\mathrm{ch}}(\mathcal{H}_k)$, which is
cocommutative (Proposition~\ref{prop:fock-bar-resolution}(i)).

Under graded linear dualization (valid because $\mathcal{H}_k$
has finite-type grading), the comodule structure becomes a module
structure over $\mathcal{H}_k^! = \mathrm{Sym}^{\mathrm{ch}}(V^*)$.
The charge is preserved: the $J_0$-eigenvalue $\lambda$ of
$|\lambda\rangle$ dualizes to the $\alpha_0$-eigenvalue $\lambda$ on
$F_\lambda^!$, because the Koszul pairing $\langle J, \alpha \rangle = 1$
identifies $J_0$ with $\alpha_0$ at the Zhu algebra level
($A(\mathcal{H}_k) \cong \mathbb{C}[J_0]
\cong \mathbb{C}[\alpha_0] \cong
A(\mathrm{Sym}^{\mathrm{ch}}(V^*))$; see
Example~\ref{ex:zhu-algebras}(iv)).

The curvature on $F_\lambda^!$ is inherited from the algebra-level
curvature $m_0 = -k\,\omega$. The module action of $m_0$
on $F_\lambda^!$ is well-defined because $m_0$ lies in the center of
$\mathrm{Sym}^{\mathrm{ch}}(V^*)$ (it is a quadratic element of a
commutative algebra).
\end{proof}

\begin{corollary}[Fock module character from Koszul resolution; \ClaimStatusProvedHere]
\label{cor:fock-character-koszul}
The character of $F_\lambda$ can be derived from the Koszul resolution:
\begin{equation}
\operatorname{ch}(F_\lambda)
= \sum_{n \geq 0} (-1)^n\,
 \operatorname{ch}(\mathcal{H}_k) \cdot
 \operatorname{ch}((\mathcal{H}_k^!)_n) \cdot
 \operatorname{ch}_{\mathrm{naive}}(F_\lambda)
= q^{\lambda^2/2k}\,\prod_{n \geq 1}\frac{1}{1-q^n}.
\end{equation}
\end{corollary}

\begin{proof}
By Theorem~\ref{thm:koszul-resolution-module},
the Koszul resolution of $F_\lambda$ has terms
$\mathcal{P}_n = \mathcal{H}_k \otimes (\mathcal{H}_k^!)_n
\otimes F_\lambda$. The alternating character sum is:
\[
\operatorname{ch}(F_\lambda)
= \operatorname{ch}(\mathcal{H}_k) \cdot
 \operatorname{ch}_{\mathrm{naive}}(F_\lambda) \cdot
 \sum_{n \geq 0} (-1)^n \operatorname{ch}((\mathcal{H}_k^!)_n)
= \operatorname{ch}(\mathcal{H}_k) \cdot
 \operatorname{ch}_{\mathrm{naive}}(F_\lambda) \cdot
 h_{\mathcal{H}_k^!}(-t).
\]
The Koszul relation $h_{\mathcal{H}_k}(t) \cdot
h_{\mathcal{H}_k^!}(-t) = 1$
(equivalent to Koszulity of $\mathcal{H}_k$;
Theorem~\ref{thm:heisenberg-koszul-dual-early}) cancels
$\operatorname{ch}(\mathcal{H}_k)$, giving
$\operatorname{ch}(F_\lambda) =
\operatorname{ch}_{\mathrm{naive}}(F_\lambda)
= q^{\lambda^2/2k}\,\prod_{n \geq 1}(1-q^n)^{-1}$.

The non-trivial content is the \emph{exactness} of the resolution: the
Koszul property of $\mathcal{H}_k$ (established by the spectral
sequence collapse in Theorem~\ref{thm:heisenberg-koszul-dual-early})
guarantees that the resolution is acyclic.
\end{proof}

\begin{proposition}[Ext groups between Fock modules; \ClaimStatusProvedHere]
\label{prop:fock-ext}
\index{Fock module!Ext computation}
\index{Ext groups!Fock modules}
Let $\dim V = d$ and let $F_\lambda$, $F_\mu$ be Fock modules for
$\mathcal{H}_k$ with $J_0$-eigenvalues $\lambda, \mu \in V^*$.
\begin{enumerate}[label=\textup{(\roman*)}]
\item If $\lambda \neq \mu$, then
$\mathrm{Ext}^i_{\mathcal{H}_k}(F_\lambda, F_\mu) = 0$
for all $i \geq 0$.
\item If $\lambda = \mu$, the graded Ext algebra is an exterior
algebra on countably many generators:
\begin{equation}\label{eq:fock-ext-generating}
\sum_{i \geq 0} \sum_{j \geq 0}
\dim \mathrm{Ext}^i_{\mathcal{H}_k}(F_\lambda, F_\lambda)_j
\; t^i q^j
\;=\; \prod_{n=1}^{\infty}(1 + tq^n)^d.
\end{equation}
In particular,
$\dim \mathrm{Ext}^i_{\mathcal{H}_k}(F_\lambda, F_\lambda)_j$
equals the number of $d$-coloured partitions of~$j$ into exactly~$i$
distinct parts.
\end{enumerate}
\end{proposition}

\begin{proof}
Module Koszul duality
(Corollary~\ref{cor:heisenberg-module-equivalence}) sends
Fock modules (standard modules over~$\mathcal{H}_k$) to
simple modules over the Koszul dual~$\mathrm{Sym}^{\mathrm{ch}}(V^*)$:
the image of $F_\lambda$ under the Koszul duality functor is the
one-dimensional module~$\mathbb{C}_\lambda$ on which $\alpha_0$ acts
by~$\lambda$ and all other modes act by~$0$. The Ext exchange gives
\begin{equation}\label{eq:fock-ext-exchange}
\mathrm{Ext}^i_{\mathcal{H}_k}(F_\lambda, F_\mu)
\;\cong\;
\mathrm{Ext}^i_{\mathrm{Sym}^{\mathrm{ch}}(V^*)}
(\mathbb{C}_\mu, \mathbb{C}_\lambda),
\end{equation}
with reversal of arguments (contravariant Koszul duality).

\emph{Part~(i).}
The simple modules $\mathbb{C}_\lambda$ and $\mathbb{C}_\mu$ lie in
different $\alpha_0$-eigenspaces. Since any projective resolution
preserves the $\alpha_0$-eigenvalue, $\mathrm{Hom}$ between
resolutions in distinct eigenspaces vanishes.

\emph{Part~(ii).}
For $\lambda = \mu$, the simple module $\mathbb{C}_\lambda$ over
$\mathrm{Sym}^{\mathrm{ch}}(V^*)$ admits a Koszul resolution by
free modules:
\[
\cdots \to \mathrm{Sym}^{\mathrm{ch}}(V^*) \otimes \Lambda^2(V^*_{<0})
\to \mathrm{Sym}^{\mathrm{ch}}(V^*) \otimes V^*_{<0}
\to \mathrm{Sym}^{\mathrm{ch}}(V^*) \to \mathbb{C}_\lambda \to 0,
\]
where $V^*_{<0} = \bigoplus_{n \geq 1} V^* \cdot z^n$ is the
space of negative modes (with $\alpha_{-n}^{(a)}$ contributing
conformal weight~$n$ for $a = 1, \ldots, d$), and the Koszul
differential is
$d(s \otimes v_1 \wedge \cdots \wedge v_{i+1})
= \sum_j (-1)^{j+1} s \cdot v_j \otimes
v_1 \wedge \cdots \widehat{v_j} \cdots \wedge v_{i+1}$.

Applying $\mathrm{Hom}_{\mathrm{Sym}^{\mathrm{ch}}}
(-, \mathbb{C}_\lambda)$:
\[
\mathrm{Hom}_{\mathrm{Sym}^{\mathrm{ch}}}
(\mathrm{Sym}^{\mathrm{ch}} \otimes \Lambda^i(V^*_{<0}),\;
\mathbb{C}_\lambda)
\;\cong\; \Lambda^i(V^*_{<0})^*
\]
with induced differential $d^*(f) = f \circ d$. Since $d$ involves
multiplication by elements of $V^*_{<0} \subset
\ker(\mathrm{Sym}^{\mathrm{ch}} \twoheadrightarrow \mathbb{C}_\lambda)$
(the augmentation ideal), the induced differential $d^* = 0$.

Therefore
\[
\mathrm{Ext}^i_{\mathrm{Sym}^{\mathrm{ch}}(V^*)}
(\mathbb{C}_\lambda, \mathbb{C}_\lambda)
\;\cong\; \Lambda^i(V^*_{<0})^*
\;\cong\; \Lambda^i(V^*_{<0})
\]
as graded vector spaces (the last isomorphism uses
finite-dimensionality of each weight space).
The generator $\alpha_{-n}^{(a)} \in V^*_{<0}$ contributes weight~$n$,
so the graded generating function is
\[
\sum_{i,j} \dim \Lambda^i(V^*_{<0})_j\; t^i q^j
= \prod_{n=1}^{\infty}\prod_{a=1}^{d} (1 + t q^n)
= \prod_{n=1}^{\infty}(1 + tq^n)^d. \qedhere
\]
\end{proof}

\begin{remark}[Ext Poincar\'e series and Jacobi functions]
\label{rem:fock-ext-jacobi}
For $d = 1$, the Ext generating function
$\prod_{n=1}^\infty (1 + tq^n)$ at $t = -1$ gives
$\prod(1 - q^n) = q^{-1/24}\,\eta(\tau)$, recovering the
Dedekind eta function. At $t = 1$, the series
$\prod(1 + q^n) = q^{-1/24}\,\eta(2\tau)/\eta(\tau)$
counts the total dimension $\sum_i \dim \mathrm{Ext}^i$
in each conformal weight. For general~$d$, the specialization
$t = -1$ gives $\eta(\tau)^d$, consistent with the Euler
characteristic of the Ext algebra being the ratio of Heisenberg
to dual characters.
\end{remark}

\begin{remark}[Ext algebra and the genus universality invariant]
\label{rem:ext-genus-universality}
\index{genus universality!Ext algebra}
The Euler characteristic
$\chi(\mathrm{Ext}^*)=\prod(1-q^n)^d=q^{-d/24}\eta(\tau)^d$
satisfies
$\chi(\mathrm{Ext}^*)\cdot Z_{\mathrm{osc}}(\mathcal{H}_k)=1$ for
the oscillator-normalized partition function
$Z_{\mathrm{osc}}=q^{d/24}/\eta(\tau)^d$. The conformal anomaly
$q^{d/24}$ records the oscillator central charge $c=d$. The genus
universality invariant records the level:
$\kappa(\mathcal{H}_k)=k$ in rank one, and
$\kappa(\mathcal{H}_k^{\oplus d})=dk$. Hence the genus-$g$
obstruction is
$\mathrm{obs}_g=\kappa(\mathcal{H}_k^{\oplus d})\,\lambda_g$
\textup{(}UNIFORM-WEIGHT\textup{}), with rank-one genus-$1$ value
$F_1=k/24$. This equals $c/24$ only in the standard normalization
$d=1$, $k=1$; in general central charge and level are distinct.
\end{remark}


\subsection{Twisted modules: Ramond and Neveu--Schwarz sectors}
\label{sec:twisted-modules-free-fields}

The free fermion and Heisenberg algebras possess \emph{twisted modules}
arising from automorphisms of the chiral algebra. Module Koszul duality
interacts non-trivially with these twisted sectors.

\begin{definition}[Ramond and Neveu--Schwarz sectors]\label{def:ramond-ns}
\index{Ramond sector|textbf}\index{Neveu--Schwarz sector|textbf}
Let $\mathcal{F}$ be the free fermion chiral algebra with generator
$\psi(z)$ and automorphism $\sigma: \psi \mapsto -\psi$ (the fermion
parity involution).
\begin{enumerate}[label=\textup{(\roman*)}]
\item The \emph{Neveu--Schwarz (NS) sector} is the untwisted vacuum
 module: modes $\psi_{r}$ with $r \in \mathbb{Z} + \tfrac{1}{2}$,
 $\psi_r|0\rangle_{\mathrm{NS}} = 0$ for $r > 0$. Character:
 \[
 \operatorname{ch}(\mathrm{NS})
 = q^{-1/48}\,\prod_{n \geq 1}(1 + q^{n-1/2}).
 \]
\item The \emph{Ramond (R) sector} is the $\sigma$-twisted module:
 modes $\psi_n$ with $n \in \mathbb{Z}$, $\psi_n|0\rangle_{\mathrm{R}} = 0$
 for $n > 0$. The zero mode $\psi_0$ satisfies $\psi_0^2 = \tfrac{1}{2}$,
 so the ground states form a $1$-dimensional Clifford module
 ($|+\rangle$ with $\psi_0|+\rangle = \tfrac{1}{\sqrt{2}}|+\rangle$
 for a single fermion).
 Character:
 \[
 \operatorname{ch}(\mathrm{R})
 = q^{1/24}\,\prod_{n \geq 1}(1 + q^n).
 \]
\end{enumerate}
The spectral flow automorphism $U_\theta$ with $\theta = 1/2$
interpolates between the two sectors:
$U_{1/2}: \psi_r \mapsto \psi_{r+1/2}$, sending NS to R.
\end{definition}

\begin{proposition}[Twisted module Koszul duality for fermions; \ClaimStatusConditional]
\label{prop:twisted-fermion-kd}
\index{free fermion!twisted module Koszul duality}
Under the fermion--boson Koszul duality
$\mathcal{F}^! \simeq \mathrm{Sym}^{\mathrm{ch}}(\gamma)$:
\begin{enumerate}[label=\textup{(\roman*)}]
\item The NS sector of $\mathcal{F}$ corresponds to the vacuum module
 of $\mathrm{Sym}^{\mathrm{ch}}(\gamma)$.
\item The Ramond sector of $\mathcal{F}$ corresponds to the
 $\sigma^!$-twisted module of $\mathrm{Sym}^{\mathrm{ch}}(\gamma)$,
 where $\sigma^!: \gamma \mapsto -\gamma$ is the dual involution.
\item The spectral flow $U_{1/2}$ on $\mathcal{F}$ dualizes to a
 shift functor $\gamma_n \mapsto \gamma_{n+1/2}$ on
 $\mathrm{Sym}^{\mathrm{ch}}(\gamma)$-modules.
\end{enumerate}
\end{proposition}

\begin{proof}
The automorphism $\sigma: \psi \mapsto -\psi$ induces an automorphism
$\sigma^!$ of the Koszul dual by functoriality of the bar construction
and graded dualization. Since the bar construction is functorial for
algebra automorphisms
(Theorem~\ref{thm:bar-cobar-isomorphism-main}), and since $\sigma$ acts
on the generator space $V = \mathrm{Span}\{\psi\}$ by $-1$, the induced
automorphism $\sigma^! = \bar{B}(\sigma)^\vee$ acts on
$V^* = \mathrm{Span}\{\gamma\}$ by $-1$ as well (the Koszul pairing
$\langle \psi, \gamma \rangle = 1$ is preserved under the simultaneous
sign change).

For (i): the NS sector is the vacuum module of $\mathcal{F}$, and
module Koszul duality (Theorem~\ref{thm:e1-module-koszul-duality}) sends
the vacuum module to the vacuum module.

For (ii): a $\sigma$-twisted module $M$ of $\mathcal{F}$ is, by
definition, a module where the monodromy around the insertion point
acts by $\sigma$. Under the bar construction, the monodromy operator
is the holonomy of the flat connection on
$\overline{C}_{n+1}(X) \setminus D$; the twist by $\sigma$ on the
algebra side becomes a twist by $\sigma^!$ on the dual side.

For (iii): spectral flow shifts the mode index by $1/2$. Since the
Koszul pairing identifies $\psi_r$ with $\gamma_r$
(at the level of generator modes), the shift
$\psi_r \mapsto \psi_{r+1/2}$ induces
$\gamma_r \mapsto \gamma_{r+1/2}$.
\end{proof}

\begin{remark}[Character identity under twisted KD]
\label{rem:twisted-character-identity}
The Koszul character relation
$h_{\mathcal{F}}(t) \cdot h_{\mathcal{F}^!}(-t) = 1$ restricts to
each twisted sector. For the Ramond sector:
\[
\prod_{n \geq 1}(1 + q^n) \cdot \prod_{n \geq 1}(1 - q^n)^{-1}
= \frac{1}{\prod_{n \geq 1}(1 - q^n)/(1 + q^n)}
= \prod_{n \geq 1}\frac{1 + q^n}{1 - q^n},
\]
which is the character of the $\sigma^!$-twisted module of
$\mathrm{Sym}^{\mathrm{ch}}(\gamma)$ (the symmetric algebra on
integer-moded generators).
\end{remark}

\begin{definition}[Heisenberg spectral flow]\label{def:heisenberg-spectral-flow}
\index{spectral flow!Heisenberg}
The Heisenberg algebra $\mathcal{H}_k$ admits a one-parameter
family of spectral flow automorphisms $\sigma_\lambda$ for
$\lambda \in \mathbb{C}$:
\[
\sigma_\lambda(J_n) = J_n + k\lambda\,\delta_{n,0},
\qquad
\sigma_\lambda(L_n) = L_n + \lambda J_n
+ \tfrac{k\lambda^2}{2}\,\delta_{n,0}.
\]
The spectral flow acts on Fock modules by charge shift:
$\sigma_\lambda^*(F_\mu) \cong F_{\mu + k\lambda}$.
\end{definition}

\begin{proposition}[Spectral flow under Koszul duality; \ClaimStatusConditional]
\label{prop:spectral-flow-kd}
Under the Heisenberg Koszul duality
$\mathcal{H}_k^! \simeq \mathrm{Sym}^{\mathrm{ch}}(V^*)$:
\begin{enumerate}[label=\textup{(\roman*)}]
\item The spectral flow $\sigma_\lambda$ on $\mathcal{H}_k$
 dualizes to the translation automorphism
 $\tau_{k\lambda}: \alpha_0 \mapsto \alpha_0 + k\lambda$ on
 $\mathrm{Sym}^{\mathrm{ch}}(V^*)$.
\item The Fock module correspondence
 \textup{(}Proposition~\textup{\ref{prop:fock-koszul-dual})} is
 equivariant: $\sigma_\lambda^*(F_\mu)^! \cong
 \tau_{k\lambda}^*(F_\mu^!)$.
\item The curvature $m_0 = -k\,\omega$ is invariant under
 spectral flow: $\sigma_\lambda^*(m_0) = m_0$.
\end{enumerate}
\end{proposition}

\begin{proof}
For (i): the spectral flow shifts the zero mode
$J_0 \mapsto J_0 + k\lambda$. Under the Zhu algebra
isomorphism $A(\mathcal{H}_k) \cong \mathbb{C}[J_0]
\cong \mathbb{C}[\alpha_0] \cong A(\mathrm{Sym}^{\mathrm{ch}}(V^*))$
(Example~\ref{ex:zhu-algebras}(iv)), this becomes the translation
$\alpha_0 \mapsto \alpha_0 + k\lambda$.

For (ii): $\sigma_\lambda^*(F_\mu) \cong F_{\mu + k\lambda}$
by Definition~\ref{def:heisenberg-spectral-flow}, and
$F_{\mu + k\lambda}^!$ has $\alpha_0$-eigenvalue
$\mu + k\lambda$ by Proposition~\ref{prop:fock-koszul-dual}.
On the dual side, $\tau_{k\lambda}^*(F_\mu^!)$ shifts
$\alpha_0 \mapsto \alpha_0 + k\lambda$, giving eigenvalue
$\mu + k\lambda$. The module correspondence is therefore equivariant
with the same charge shift. The level is not inverted here; it enters
only as the scalar multiplying the spectral-flow parameter.

For (iii): the curvature $m_0 = -k\,\omega$ lies in
$\mathrm{Sym}^2(V^*)$ and is independent of $J_0$ (it comes from
the quadratic pole, not the zero mode). Since $\sigma_\lambda$
shifts only $J_0$, the curvature is invariant.
\end{proof}

\begin{example}[bc system: twisted sectors and spectral flow]
\label{ex:bc-spectral-flow}
\index{bc system!spectral flow}
The $bc$ ghost system (weights $(j, 1-j)$) has spectral flow
$\sigma_n: b_m \mapsto b_{m+n}$, $c_m \mapsto c_{m-n}$, which shifts
the ghost number vacuum by $n$ units. Under the
$bc \leftrightarrow \beta\gamma$ Koszul duality
(Theorem~\ref{thm:betagamma-bc-koszul}), this spectral flow dualizes
to the corresponding $\beta\gamma$ spectral flow
$\sigma_n^!: \beta_m \mapsto \beta_{m+n}$,
$\gamma_m \mapsto \gamma_{m-n}$.

At $j = 2$ (conformal ghosts), the Ramond-type twisted sector
($n = 1$) gives the ``picture-changed'' BRST complex of string
theory: the picture-changing operator $\mathcal{X}(z) = \{Q, \xi(z)\}$
is the bar complex differential in the spectrally flowed sector.
\end{example}

\subsection{Koszul duality summary table}\label{sec:kd-computations-free-fields}

\subsubsection{Koszul duality table}

\begin{center}
\footnotesize
\begin{tabular}{|l|l|l|l|l|}
\hline
\textbf{Algebra $\mathcal{A}$} & \textbf{Koszul Dual $\mathcal{A}^!$} & \textbf{Type} & \textbf{Context} & \textbf{Status} \\
\hline
Free fermion $\psi$ & $\mathrm{Sym}^{\mathrm{ch}}(\gamma)$ & Exact & Lie/Com & \ClaimStatusProvedHere \\
$bc$ ghosts & $\beta\gamma$ system & Exact & D-branes & \ClaimStatusProvedHere \\
Heisenberg $\mathcal{H}_k$ & $\mathrm{Sym}^{\mathrm{ch}}(V^*)$ (curved) & Exact & Lie/Com & \ClaimStatusProvedHere \\
Free boson $\partial\phi$ & Symplectic bosons & Exact & Open-closed & \ClaimStatusProvedElsewhere~\cite{FBZ04} \\
$\widehat{\mathfrak{g}}_k$ & $\mathrm{CE}^{\mathrm{ch}}(\widehat{\mathfrak{g}}_{-k-2h^\vee})$ & Filt./Curved & Feigin--Frenkel & \ClaimStatusProvedHere\textsuperscript{a} \\
Virasoro$_c$ & $\mathrm{Vir}_{26-c}$ (dynamic shadow) & Curved & $c \mapsto 26-c$ & \ClaimStatusProvedHere\textsuperscript{c} \\
$\mathcal{W}_N$ & $\mathcal{W}_N^!$ (cumulant-enlarged) & Filt./Curved & Higher spin & \ClaimStatusConjectured\textsuperscript{d} \\
Super-Virasoro & Super-$W_\infty$ & Curved & AdS$_3$ sugra & \ClaimStatusConjectured \\
$Y(\mathfrak{g})$ & $Y_{R^{-1}}(\mathfrak{g})$ & Self-dual & Integrability & \ClaimStatusProvedHere\textsuperscript{b} \\
Virasoro$_{26}\otimes bc$ & Critical total string vertex algebra & Classical & Genus-$0$ BRST/bar; matter-only comparison conditional & \ClaimStatusConditional \\
$W^{-h^\vee}(\mathfrak{g})$ & Wakimoto & Classical & DS reduction $\leftrightarrow$ free field & \ClaimStatusProvedElsewhere \\
Lattice $V_L$ & Lattice $V_{L^*}$ & Classical & Form duality & \ClaimStatusProvedElsewhere \\
\hline
\end{tabular}
\end{center}
\medskip
\noindent\textsuperscript{a} The Feigin--Frenkel isomorphism
$\mathfrak{z}(\hat{\mathfrak{g}}_{-h^\vee}) \cong
\mathcal{W}(\hat{\mathfrak{g}}^\vee)$ is a theorem of
Feigin--Frenkel. The level involution
$k \leftrightarrow -k-2h^\vee$ supplies the reflected current
presentation inside the chiral CE Koszul dual. This is proved in
Chapter~\ref{chap:kac-moody-koszul}: the low-degree bar
calculation
\textup{(}Theorem~\textup{\ref{thm:sl2-koszul-dual}} for
$\mathfrak{sl}_2$\textup{)} together with the Serre-duality argument on
configuration spaces
\textup{(}Theorem~\textup{\ref{thm:universal-kac-moody-koszul})}
establishes the level shift.

\noindent\textsuperscript{c} The same-family shadow
$\mathrm{Vir}_{26-c}$ is the \emph{dynamic} Koszul dual at the M/S-level
(Example~\ref{ex:virasoro-koszul-dual}). The full \emph{kinematic}
dual has infinitely many primitive cumulants
($\Delta_{\mathrm{Vir}}(t) = t^3 + 2t^5 + \cdots \neq 0$);
the H-level identification with a $\mathcal{W}_\infty$-type object
is the remaining infinite-generator identification
(Theorem~\ref{thm:completed-bar-cobar-strong} provides the structural
framework; the specific target identification is open).

\noindent\textsuperscript{d} The naive dual on generators $L^*, W^*$
is insufficient: the composite field $\Lambda = :(TT): -
\tfrac{3}{10}\partial^2 T$ forces a cumulant cogenerator $\Lambda^*$
with nonprimitive coproduct
(Remark~\ref{rem:cumulant-dual-nonquadratic}).

\noindent\textsuperscript{b} The Yangian $\Eone$-chiral Koszul duality
$Y(\mathfrak{g})^! \cong Y_{R^{-1}}(\mathfrak{g})$ is proved in
Theorem~\ref{thm:yangian-koszul-dual}: the RTT presentation is
quadratic in the $\Eone$-chiral sense, and the quadratic dual
relations are the RTT relations with the inverse $R$-matrix
$R^{-1}(u)$. At the classical level ($\hbar = 0$), this reduces to
exact self-duality
\textup{(}Corollary~\textup{\ref{cor:yangian-classical-self-dual})}.

\begin{remark}[Koszul duality identifications]\label{rem:koszul-table-status}
The entries fall into two categories.
\emph{Proved}: the free-field and Heisenberg dualities (Chapter~\ref{sec:betagamma-koszul-dual}, Theorem~\ref{thm:heisenberg-koszul-dual-early}); the Kac--Moody CE dual with reflected current presentation $\widehat{\mathfrak{g}}_{-k-2h^\vee}$ (Theorems~\ref{thm:sl2-koszul-dual} and~\ref{thm:universal-kac-moody-koszul}).
\emph{Infinite-generator}: the identifications
Virasoro $\rightsquigarrow W_\infty$,
$\mathcal{W}_N \rightsquigarrow Y(\mathfrak{gl}_N)$, and
Super-Virasoro $\rightsquigarrow$ Super-$W_\infty$ require completed
bar complexes. The M-level theory (bar complex, inversion,
Koszulness, complementarity) is proved for all principal
$\mathcal{W}_N$ at generic level. On the $\mathcal{W}_\infty$ side,
the coefficient identities
$\mathsf{C}^{\mathrm{res}}=\mathsf{C}^{\mathrm{DS}}$ on
$\mathcal{I}_N$ follow from $\mathcal{W}_N$ rigidity
(Theorem~\ref{thm:winfty-all-stages-rigidity-closure}) and the
strong completion theorem
(Theorem~\ref{thm:completed-bar-cobar-strong}). On the Yangian side,
the evaluation-core comparison
$\mathsf{K}^{\mathrm{line}}=\mathsf{K}^{\mathrm{RTT}}$ on
$\Delta_{a,0}(N)$ is a computed invariant; the remaining categorical
identification is the compact-completion problem in type~A.
\end{remark}

\begin{remark}[Infinite-generator table entries]
The three infinite-generator identifications each involve curved
Koszul duality ($m_0 \neq 0$) with infinitely many generators. Two
structural features: (1)~non-quadratic OPE prevents the quadratic
framework of Chapter~\ref{chap:koszul-pairs} from applying directly;
(2)~the infinite-generator comparison requires the
coefficient-identification theorem, proved on the
$\mathcal{W}_\infty$ side by $\mathcal{W}_N$ rigidity. Conditional
on the DS/bar transport theorem
\textup{(}Theorem~\textup{\ref{thm:bar-semi-infinite-w}}\textup{)},
the same-family involution $\mathrm{Vir}_c \leftrightarrow
\mathrm{Vir}_{26-c}$ gives the M/S-level shadow. The H-level
realization requires an infinite-generator dual whose bar cohomology
matches the semi-infinite growth (Motzkin numbers,
Theorem~\ref{thm:virasoro-chiral-koszul}). The super row remains a
separate infinite-generator extension
(Conjecture~\ref{conj:v1-master-infinite-generator}).
\end{remark}

\begin{remark}[Heisenberg vs Yangian: self-duality contrast]\label{rem:heisenberg-yangian-contrast}
The Heisenberg algebra is \emph{not} self-dual: $\mathcal{H}^! = \mathrm{Sym}^{\mathrm{ch}}(V^*) \not\cong \mathcal{H}$, since the central extension and non-commutative oscillator modes dualize to the commutative chiral algebra on the dual space (Lie $\leftrightarrow$ Com duality, not to be confused with the boson-fermion correspondence). By contrast, the Yangian satisfies $Y(\mathfrak{g})^! \cong Y_{R^{-1}}(\mathfrak{g})$ (Theorem~\ref{thm:yangian-koszul-dual}): the RTT presentation is quadratic in the $\Eone$-chiral sense, the dual relations use the inverse R-matrix $R^{-1}(u)$, and at the classical level ($\hbar = 0$) this reduces to exact self-duality (Corollary~\ref{cor:yangian-classical-self-dual}), visible in 3d mirror symmetry (Higgs $\leftrightarrow$ Coulomb).

The Heisenberg OPE $\alpha(z)\alpha(w) \sim (z-w)^{-2}$ has a double pole that produces symmetric (bosonic) coproduct structure in the bar construction. By contrast, the Yangian's RTT relations have built-in $R$-matrix self-duality that preserves the structure under bar-cobar.
\end{remark}

\begin{remark}[Additional dualities]\label{rem:additional-dualities}
Some algebras in this table have \emph{additional} duality structures beyond standard Koszul duality:

\begin{enumerate}
\item \emph{Heisenberg.} Level inversion $\mathcal{H}_k \leftrightarrow \mathcal{H}_{-k}$ (curved duality, not Koszul).

\item \emph{Affine Kac--Moody.} Langlands duality $\widehat{\mathfrak{g}}_k \leftrightarrow \widehat{\mathfrak{g}^\vee}_{k'}$ at critical/conformal levels.

\item \emph{W-algebras.} At critical level, the external quantum-Langlands picture relates $\mathcal{W}^{-h^\vee}(\mathfrak{g}, f)$ to a Langlands-dual $\mathcal{W}$-algebra, while the manuscript-internal proved statement is the critical fixed-point inversion package for the principal algebra. At general central charge: $\mathcal{W}_N^c \leftrightarrow \mathcal{W}_N^{c'}$ with $c + c' = 2(N-1)(2N^2+2N+1)$.

\item \emph{Virasoro.} Exceptional modular invariance at certain central charges.
\end{enumerate}

These additional structures are curved or filtered dualities in
extended categories. They are distinct from the standard Koszul
duality listed in the main table column.
\end{remark}

\begin{remark}[Completion kinematics of free-field archetypes]
\label{rem:free-field-completion-kinematics}
\index{completion kinematics!free fields}
\index{completion entropy!free fields}
The free-field archetypes occupy the simplest corner of the completion
geometry
(Table~\ref{tab:completion-kinematics} in the landscape census):
\begin{center}
\small
\begin{tabular}{lcccl}
\textbf{Algebra} & $G_\cA(t)$ & $\rho_K$ & $h_K$ & \textbf{Class} \\
\hline
Heisenberg $\mathcal{H}_k$
 & $t$ & $1$ & $0$ & G \\
Free fermion $\psi$
 & $t$ & $1$ & $0$ & G \\
$\beta\gamma$
 & $t/(1-t)$ & $1/2$ & $\log 2$ & C \\
$\hat{\mathfrak{g}}_k$
 & $(\dim\mathfrak{g})\,t/(1{-}t)$ & $1/(\dim\mathfrak{g}+1)$ & $\log(\dim\mathfrak{g}+1)$ & L \\
\end{tabular}
\end{center}
All have $\Delta_\cA(t) = 0$: the declared generators span the
full primitive sector, with no cumulant enlargement needed. The
G-class families (Heisenberg, free fermion) have $h_K = 0$, meaning
the completed bar complex does not grow exponentially; the
completion is trivial. The L- and C-class families have positive
entropy controlled by the number of generator flavours.

The reduced-weight windows $K_q(\cA)$
(Definition~\ref{def:finite-window}) make the bar complex finite at
each level: for Heisenberg, $\dim K_q = q$ (one bar word per reduced
weight); for $\hat{\mathfrak{g}}_k$, $\dim K_q \sim
(\dim\mathfrak{g}+1)^q$. This is the kinematic underpinning of the
finite-stage bar computations in the sections above.
\end{remark}

\subsubsection{Computing the Koszul dual from bar cohomology}

\begin{remark}[Procedure]\label{rem:koszul-dual-procedure}
Given a quadratic chiral algebra $\mathcal{A}$ with generators
$\{a_i\}$ and relations $\{R_j\}$, the strict Koszul dual algebra
$\mathcal{A}^!$ is computed as follows.
\begin{enumerate}
\item Write $\mathcal{A} = T(V)/(R)$ where $R \subset V^{\otimes 2}$.
\item Define the pairing $\langle \cdot, \cdot \rangle\colon V \otimes V^* \to \mathbb{C}$ and compute $R^\perp \subset (V^*)^{\otimes 2}$.
\item Compute the Koszul dual coalgebra
$\mathcal{A}^{\mathrm{i}}=H^\bullet B(\mathcal{A})$; on the
quadratic Koszul locus this is the coalgebra cogenerated by $sV^*$
with corelations $s^2R^\perp$.
\item Set $\mathcal{A}^! = (\mathcal{A}^{\mathrm{i}})^\vee$,
equivalently $T(V^*)/(R^\perp)$ in finite type after the Verdier
linear duality convention is fixed.
\item Check Koszulity: if
$\mathrm{Tor}_{\mathcal{A}}^{i,j}(\mathbb{C}, \mathbb{C})=0$ for
$i\neq j$, the strict quadratic Koszul duality is exact. If this
diagonal condition fails, strict Koszulity fails; a curved duality
requires a complete filtered category, a central curvature element
$m_0$, and the Positselski finiteness hypotheses.
\end{enumerate}
The cobar counit
$\Omega B(\mathcal{A})\xrightarrow{\sim}\mathcal{A}$ is the
bar-cobar inversion theorem. It recovers $\mathcal{A}$; it is not the
operation that defines $\mathcal{A}^!$.
\end{remark}

\subsubsection{\texorpdfstring{Explicit example: $\beta\gamma$ $\leftrightarrow$ $bc$}{Explicit example: beta-gamma bc}}

This calculation is carried out in detail in Theorem~\ref{thm:betagamma-bc-koszul} above (see \S\ref{sec:fermion-boson-koszul}): the antisymmetric relation $\beta\gamma - \gamma\beta$ dualizes under the residue pairing to the Clifford relation $bc + cb$, establishing $(\beta\gamma)^! \cong bc$.



\section{Heisenberg higher-genus and Fourier duality}
\label{sec:heisenberg-koszul}


The Koszul dual $\mathcal{H}_k^! \simeq \mathrm{Sym}^{\mathrm{ch}}(V^*)$ was computed in Theorem~\ref{thm:heisenberg-koszul-dual-early}; see Chapter~\ref{ch:heisenberg-frame}\index{Heisenberg algebra|textbf}\index{Fock space} for conventions. The bar complex $\bar{B}^{\mathrm{ch}}(\mathcal{H}_k) \simeq \mathrm{coLie}^{\mathrm{ch}}(V^*)$ is cocommutative because the double-pole OPE produces symmetric coproducts.

\subsection{Failure of self-duality}

\begin{theorem}[Heisenberg is not self-dual; \ClaimStatusProvedHere]
\label{thm:heisenberg-not-self-dual}
For any level $k \neq 0$, the Heisenberg algebra $\mathcal{H}_k$ is not Koszul
self-dual. Instead:
\[\text{Koszul}(\mathcal{H}_k) = \mathrm{Sym}^{\mathrm{ch}}(V^*) \neq \mathcal{H}_k\]

The algebras have different structure:
\begin{center}
\begin{tabular}{c|c}
\textbf{Heisenberg } $\mathcal{H}_k$ & \textbf{Symmetric } $\mathrm{Sym}^{\mathrm{ch}}(V^*)$ \\ \hline
Centrally extended & Free commutative \\
$\{J_\lambda J\} = k\lambda \neq 0$ & $\{J^*_\lambda J^*\} = 0$ \\
Double-pole OPE & Regular OPE \\
Level datum $k\,\omega_{\mathrm{cen}}$ & Curved dual ($m_0 = -k\,\omega$) \\
Nontrivial $\pi_1$ & Trivial $\pi_1$
\end{tabular}
\end{center}
\end{theorem}

\begin{proof}
By Theorem~\ref{thm:heisenberg-bar}, the bar complex $\bar{B}^{\mathrm{ch}}(\mathcal{H}_k)$ has cocommutative coproduct: the Heisenberg OPE $\alpha(z)\alpha(w) \sim k(z-w)^{-2}$ produces only the double-pole term, whose residue gives a symmetric (primitive) coproduct $\Delta(\alpha) = \alpha \otimes 1 + 1 \otimes \alpha$. Cocommutative coalgebras are Koszul dual to commutative algebras ($\mathrm{Com}^! = \mathrm{Lie}$, see \cite{LV12}, Theorem~7.6.5), so $\mathcal{H}_k^! = \mathrm{Sym}^{\mathrm{ch}}(V^*)$. This is commutative with regular OPE, while $\mathcal{H}_k$ is non-commutative with singular OPE, so $\mathcal{H}_k \not\cong \mathcal{H}_k^!$.
\end{proof}

\subsection{Three different ``dualities'' for Heisenberg}\label{thm:three-dualities-compared}

Three separate mathematical structures govern Heisenberg, and
confusing them is a reliable source of error.

\begin{enumerate}
\item \emph{Koszul duality} (algebra $\to$ dual coalgebra $\to$ dual algebra).
 \[\mathcal{H}_k \xrightarrow{\;\bar{B}\;} \bar{B}^{\mathrm{ch}}(\mathcal{H}_k) \xrightarrow{\;(-)^\vee\;} \mathrm{Sym}^{\mathrm{ch}}(V^*)\]
 The first arrow is the bar construction: extraction of OPE residues
 along collision divisors in~$\overline{C}_n(X)$ produces the Koszul
 dual \emph{coalgebra}~$\mathrm{coLie}^{\mathrm{ch}}(V^*)$, a
 conilpotent factorization coalgebra on~$\mathrm{Ran}(X)$. The
 second arrow is linear duality on that coalgebra, producing the
 strict Koszul dual \emph{algebra}~$\mathrm{Sym}^{\mathrm{ch}}(V^*)$.
 The separate Verdier image
 $\mathbb{D}_{\operatorname{Ran}}\bar{B}^{\mathrm{ch}}(\mathcal{H}_k)$
 is the homotopy Koszul-dual factorization algebra. The cobar
 construction is \emph{not involved}: $\Omega(\bar{B}(\mathcal{H}_k))
 \simeq \mathcal{H}_k$ recovers the original algebra, not its dual
 (Theorem~\ref{thm:bar-cobar-inversion-qi}).

 The distinction is mediated by the convolution Lie algebra
 $\mathrm{Hom}(\bar{B}(\mathcal{H}_k), \mathcal{H}_k)$: the
 canonical twisting morphism
 $\tau \colon \mathrm{coLie}^{\mathrm{ch}}(V^*) \to \mathcal{H}_k$,
 $v^* \mapsto J_v$, is its Maurer--Cartan element. The convolution
 product $\tau \star \tau = 0$ because the cogenerator~$v^*$ is
 primitive ($\bar{\Delta}(v^*) = 0$), so the MC equation
 $\partial\tau + \tau \star \tau = 0$ holds on generators for
 degree reasons
 (cf.~Remark~\ref{rem:frame-twisting-morphism}).

 Level behavior: $k$ appears as $+k$ in the bar-side curved cobar
 convention and as $-k$ in the Verdier-linear Koszul-dual algebra;
 it does not change the duality type.

 Structure exchanged: Lie (non-commutative) $\leftrightarrow$ Com (commutative).

\item \emph{Level inversion} (opposite curvature convention).
 Over $\mathbb{C}$ the unstarred vertex algebras
 $\mathcal H_k$ and $\mathcal H_{-k}$ are identified by the rescaling
 $J\mapsto iJ$; with a real or unitary structure this becomes an
 opposite-curvature convention rather than an equality of unitary
 theories.

 Level behavior: $k \leftrightarrow -k$ only in this rescaled or
 opposite-curvature sense.

 Structure exchanged: the sign of the central pairing, not the
 bar coalgebra with its Verdier-linear Koszul dual.

\item \emph{Boson-fermion correspondence} (categorical equivalence).
 \[\mathrm{Rep}(\mathcal{H}_k) \simeq \mathrm{Rep}(\mathcal{F}^{\otimes 2})\]
 where $\mathcal{F}$ is the free fermion algebra.

 Level behavior: Relates Heisenberg to fermions, not to itself.

 Structure exchanged: Module categories have equivalent structure.
\end{enumerate}

Only~(1) is the subject of this manuscript. The three dualities
act on different mathematical categories: (1)~on algebras and
coalgebras (via the bar-cobar adjunction on~$\mathrm{Ran}(X)$),
(2)~on representation categories (via autoequivalences),
(3)~on module categories (via bosonization). The bar construction
is the engine of~(1); it has no direct role in~(2) or~(3).

\subsection{Costello--Gwilliam's construction}

\begin{theorem}[Heisenberg duality \cite{CG17}; \ClaimStatusProvedElsewhere]
\label{thm:rhagavendran-heisenberg}
The Heisenberg chiral algebra arises from:
\begin{itemize}
\item \emph{Factorization algebra.} Lie algebra cochains $C^*(\mathcal{O}_X)$ on the
 abelian Lie algebra $\mathcal{O}_X$ of holomorphic functions.
\item \emph{Koszul dual.} Factorization envelope = Lie algebra chains $C_*(\mathcal{O}_X^c)$
 on compactly supported sections.
\end{itemize}

Under the factorization homology $\int_X$:
\begin{align*}
C^*(\mathcal{O}_X) &\rightsquigarrow \mathcal{H}_k \quad \text{(Heisenberg)} \\
C_*(\mathcal{O}_X^c) &\rightsquigarrow \text{Sym}(\mathcal{O}_X) \quad \text{(Symmetric algebra)}
\end{align*}

The Koszul duality $C^* \leftrightarrow C_*$ for Lie algebra (co)homology induces:
\[\mathcal{H}_k^! \simeq \mathrm{Sym}^{\mathrm{ch}}(V^*)\]
\end{theorem}

\subsection{BV quantization and the bar complex for free fields}
\label{sec:bv-bar-free-fields}
\index{BV quantization!free fields}

\begin{proposition}[Free BV comparison criterion; \ClaimStatusConditional]
\label{prop:bar-bv-free-fields}
Let $\cA$ be a free field chiral algebra \textup{(}Heisenberg, free fermion,
$\beta\gamma$, or $bc$ ghost system\textup{)} with finite-dimensional
weight spaces and nondegenerate kinetic pairing on a Riemann surface~$X$.
Assume that the BV propagator is the Fulton--MacPherson logarithmic
propagator used in the chiral bar construction and that the anomaly
curvature is central. Then the reduced chiral bar complex
$\bar{B}^{\mathrm{ch}}(\cA)$ maps by a filtered quasi-isomorphism to
the chain-level BV complex of the corresponding Gaussian theory. The
curved term is part of this comparison datum; it is not produced by
bar-cobar inversion alone.
\end{proposition}

\begin{proof}
The genus-$0$ bar complex is computed by configuration-space
integrals on $\overline{C}_n(X)$ with the Fulton--MacPherson
logarithmic propagator. Under the stated hypothesis this is the same
kernel used to write the BV homotopy for the Gaussian action. The
free-field OPEs have no interaction vertices beyond the quadratic
pairing: Heisenberg contributes the double-pole level class, while
the free fermion, $\beta\gamma$, and $bc$ contribute simple-pole
first-order pairings or their regular ordered contact kernels. Wick's
theorem then identifies the bar contractions with the BV contractions.
Central curvature is carried as a scalar curved term in the completed
comparison, so the filtered map is a morphism of curved coalgebras.
\end{proof}

\begin{remark}[Curvature and one-loop determinants]
\label{rem:curvature-one-loop}
The ordinary Heisenberg bar differential is square-zero. The
bar-side curvature appears after passing to the curved cobar model:
for $\cH_k$ it is the level class
$m_0^{\mathrm{bar}}=k\,\omega_{\mathrm{cen}}$, while the
Verdier-linear Koszul dual carries
$m_0^!=-k\,\omega_{\mathrm{cen}}$
(Theorem~\ref{thm:heisenberg-curved-structure}). Anomaly
cancellation (Theorem~\ref{thm:anomaly-koszul}) says that the
curved cobar structure attached to $\cH^{26}\otimes bc$ has vanishing
scalar central anomaly because
$c_{\mathrm{eff}}=26\cdot 1+(-26)=0$. This cancellation is a
central-charge statement, not a $\kappa$-identity: the modular
characteristic $\kappa(\cH_1)=1$ differs from $c(\cH_1)/2=1/2$.
\end{remark}



\subsection{Higher genus: quantum complementarity}

\begin{theorem}[Quantum complementarity for Heisenberg; \ClaimStatusProvedHere]\label{thm:heisenberg-genus-g}
For the Heisenberg algebra at genus $g$ on a Riemann surface
$\Sigma_g$ with period matrix $\Omega \in \mathbb{H}_g$ (Siegel upper
half-space), the zero-mode scalar projection of the genus-$g$ chiral
bar complex is modeled by the Chevalley--Eilenberg complex of the
level-$k$ Heisenberg central extension:
\[
\bar{B}^{\text{ch}}_{g,0}(\mathcal{H}_k)
\simeq \mathrm{CE}(\hat{\mathfrak{h}}_k),
\qquad
\hat{\mathfrak h}_k=H^1(\Sigma_g,\mathbb C)\oplus\mathbb C c,
\quad [v,w]=k\,\omega(v,w)c.
\]
After descent to the Jacobian this zero-mode complex computes the
theta line cohomology:
\[
H^*(\bar{B}^{\text{ch}}_{g,0}(\mathcal{H}_k))
\cong H^*(\mathrm{Jac}(\Sigma_g), \mathcal{L}_k).
\]

The modular action on period matrices and the Koszul-dual curvature
are separate structures:
\begin{enumerate}
\item \emph{Modular transformation.} The symplectic transformation
$\Omega \to -\Omega^{-1}$ exchanges $A$- and $B$-cycles and acts on
theta sections by Poisson summation. It preserves the Heisenberg OPE
level $k$ as the coefficient of $J(z)J(w)\sim k(z-w)^{-2}$.

\item \emph{Theta functions.} The level-$k$ theta function is:
\[\Theta_k(\Omega) = \sum_{n \in \mathbb{Z}^g} e^{i\pi k \, n^T \Omega \, n}\]
(The level $k$ scales the quadratic form, not the characteristic.)
Under $\Omega \to -\Omega^{-1}$, Poisson summation gives:
\[\Theta_k(-\Omega^{-1}) = \left(\frac{\det(\Omega/i)}{k^g}\right)^{1/2} \cdot \Theta_{1/k}(\Omega)\]
in this scalar normalization. For integral positive level this is the
finite Fourier transform on the level-$k$ theta vector. It is not a
replacement of $\mathcal H_k$ by $\mathcal H_{1/k}$.

\item \emph{Koszul-dual curvature.} Verdier-linear Koszul duality
sends $\mathcal H_k$ to
$\mathcal H_k^!=(\mathrm{Sym}^{\mathrm{ch}}(V^*),
m_0^!=-k\,\omega_{\mathrm{cen}})$. This sign reversal is the
curvature of the dual algebra, not a modular level shift of the
original Heisenberg algebra.

\item \emph{Geometric interpretation.} The zero-mode bar cohomology computes:
\[H^*(\bar{B}_{g,0}^{\text{ch}}(\mathcal{H}_k)) \cong H^*(\text{Jac}(\Sigma_g), \mathcal{L}_k)\]
where $\mathcal{L}_k$ is the line bundle of level $k$ theta functions.
\end{enumerate}
\end{theorem}

\begin{proof}
We prove the theorem for arbitrary genus~$g \geq 1$ in five steps,
then illustrate Step~3 explicitly for genus~1.

\emph{Step~1: Abelian factorization of the bar differential.}
The Heisenberg OPE $\alpha(z)\alpha(w) \sim k(z-w)^{-2}$ has no
simple pole. By
Proposition~\ref{prop:abelian-bar-factorization}, the bar differential
on $\mathcal{H}_k^{\boxtimes n} \otimes
\Omega^p_{\log}(\overline{C}_n(\Sigma_g))$ decomposes as a sum of
commuting pairwise operators $d_{\mathrm{bar}} = \sum_{i<j} d_{ij}$,
each extracting the double-pole residue at the collision divisor
$D_{ij}$. In particular, the bar complex is cocommutative:
$\bar{B}^{\mathrm{ch}}(\mathcal{H}_k)
\simeq \mathrm{coLie}^{\mathrm{ch}}(V^*)$
(Theorem~\ref{thm:heisenberg-koszul-dual-early}).

\emph{Step~2: Identification with Chevalley--Eilenberg cohomology.}
On a closed Riemann surface $\Sigma_g$, the current
$\alpha(z) = \partial\phi(z)$ has $2g$ zero modes: the
integrals $\oint_{A_i} \alpha\, dz$ and $\oint_{B_j} \alpha\, dz$
($i,j = 1, \ldots, g$) over a symplectic basis of
$H_1(\Sigma_g, \mathbb{Z})$. By Hodge theory these span
$\mathfrak{h}_g := H^1(\Sigma_g, \mathbb{C}) \cong \mathbb{C}^{2g}$,
an abelian Lie algebra. The level-$k$ central extension equips
$\mathfrak{h}_g$ with the symplectic form
$\omega_k = k\,\omega$, where
$\omega \in \Lambda^2(\mathfrak{h}_g^*)$ is the intersection pairing
(normalized so that $\omega(A_i, B_j) = \delta_{ij}$).

Because the zero-mode scalar quotient of the bar complex of
$\mathcal{H}_k$ is cocommutative and the bar differential decomposes
pairwise, this quotient reduces to the Chevalley--Eilenberg complex of the one-dimensional
central extension $\hat{\mathfrak{h}}_k = \mathfrak{h}_g \oplus
\mathbb{C} \cdot c$ with bracket $[v,w] = k\,\omega(v,w)\cdot c$:
\[
\bar{B}^{\mathrm{ch}}_{g,0}(\mathcal{H}_k)
\simeq \mathrm{CE}(\hat{\mathfrak{h}}_k)
= \bigl(\Lambda^*(\mathfrak{h}_g^* \oplus \mathbb{C}),\;
 d_{\mathrm{CE}}\bigr),
\quad
d_{\mathrm{CE}}(\xi) = k\,\omega \quad
\text{on the central generator }\xi.
\]
This is the scalar-projection content of
Theorem~\ref{thm:heisenberg-higher-genus}: the pairwise
factorization gives a finite-dimensional zero-mode quotient of the
configuration-space bar complex.

\emph{Step~3: Descent to the Jacobian.}
The Abel--Jacobi map (Definition~\ref{def:jacobian-variety})
\[
\mu\colon \Sigma_g \to \mathrm{Jac}(\Sigma_g)
= \mathbb{C}^g / (\mathbb{Z}^g + \Omega\,\mathbb{Z}^g)
\]
identifies the zero-mode lattice with
$H_1(\Sigma_g, \mathbb{Z}) \cong \mathbb{Z}^{2g}$.
The level-$k$ central extension defines the theta line bundle
$\mathcal{L}_k$ on $\mathrm{Jac}(\Sigma_g)$, with Chern curvature
$k\,\omega$.
Concretely, holomorphic sections of $\mathcal{L}_k$ are functions
$f\colon \mathbb{C}^g \to \mathbb{C}$ satisfying the
quasi-periodicity conditions
\begin{align}
f(z + e_\alpha) &= f(z), \\
f(z + \Omega\, e_\beta) &= e^{-i\pi k\,\Omega_{\beta\beta}
 - 2\pi i k\, z_\beta}\, f(z),
\end{align}
which are precisely the transformation properties of
the Riemann theta function at level~$k$
(Theorem~\ref{thm:theta-properties}).

The CE differential $d_{\mathrm{CE}}$ on the zero-mode model
$\Lambda^*(\mathfrak{h}_g^*)$ maps, under this descent, to the
Dolbeault--Koszul differential computing
$H^*(\mathrm{Jac}(\Sigma_g), \mathcal{L}_k)$. The key
point is that the cochain complex
$(\Lambda^*(\mathfrak{h}_g^*),\, d_{\mathrm{CE}})$ with
$d_{\mathrm{CE}}$ determined by $k\,\omega$ is the
translation-invariant Koszul model for the theta line bundle
$\mathcal{L}_k$. The Chern curvature is $k\,\omega$; the model is
not a flat local-system computation.
This gives the isomorphism:
\[
H^*(\mathrm{CE}(\hat{\mathfrak{h}}_k))
\cong H^*(\mathrm{Jac}(\Sigma_g), \mathcal{L}_k).
\]

Combining with Step~2:
\[
H^*(\bar{B}^{\mathrm{ch}}_{g,0}(\mathcal{H}_k))
\cong H^*(\mathrm{Jac}(\Sigma_g), \mathcal{L}_k).
\]

\emph{Step~4: Modular transformation.}
The symplectic group $\mathrm{Sp}(2g, \mathbb{Z})$ acts on
$\mathrm{Jac}(\Sigma_g)$ via
$\Omega \mapsto (A\Omega + B)(C\Omega + D)^{-1}$.
For the involution $\Omega \mapsto -\Omega^{-1}$
(corresponding to $A = D = 0$, $B = -C = I_g$),
the $A$- and $B$-cycles are exchanged:
$A_i \leftrightarrow B_i$. The level-$k$ theta function
\[
\Theta_k(\Omega)
= \sum_{n \in \mathbb{Z}^g} e^{i\pi k\, n^T \Omega\, n}
\]
transforms by Poisson summation on the lattice
$\mathbb{Z}^g$:
\[
\Theta_k(-\Omega^{-1})
= \left(\frac{\det(\Omega/i)}{k^g}\right)^{1/2}
 \Theta_{1/k}(\Omega).
\]
This is the multivariable theta transformation in the scalar
normalization. It is a statement about theta sections and the
Fourier transform on the zero-mode lattice, not a statement that
the Heisenberg vertex algebra changes its OPE level. For integral
positive level one should read this as the usual finite Fourier
action on the vector of level-$k$ theta functions.

Koszul duality is a separate operation. It sends
$\mathcal{H}_k$ to the curved commutative algebra
$\mathcal{H}_k^!=(\mathrm{Sym}^{\mathrm{ch}}(V^*),m_0=-k\omega)$
(Theorem~\ref{thm:heisenberg-curved-structure}). It does not
identify $\mathcal{H}_k$ with $\mathcal{H}_{-k}$ or
$\mathcal{H}_{1/k}$.

\emph{Step~5: Geometric interpretation.}
Taking cohomology of the bar complex and applying the
isomorphisms of Steps~2--3:
\[
H^*(\bar{B}_{g,0}^{\mathrm{ch}}(\mathcal{H}_k))
\cong H^*(\mathrm{Jac}(\Sigma_g), \mathcal{L}_k).
\]
This is the cohomology of the level-$k$ line bundle
of theta functions on the Jacobian, which is the statement of
part~(3) of the theorem.

\medskip
\emph{Genus-1 illustration.}
For $g = 1$, the Jacobian is the elliptic curve
$T^2 = \mathbb{C}/(\mathbb{Z} + \tau\mathbb{Z})$ with
$\tau \in \mathbb{H}$.
The zero-mode lattice is $\mathfrak{h}_1 = \mathbb{C}^2$
with symplectic basis $\{A, B\}$ and $\omega(A,B) = 1$.
The CE complex is $\Lambda^*(\mathbb{C}^2)$ with
$d_{\mathrm{CE}}\xi = k\, dA \wedge dB$, and the local
system $\mathcal{L}_k$ has quasi-periodicity
$f(z+1) = f(z)$,
$f(z + \tau) = e^{-i\pi k\tau - 2\pi i k z}\, f(z)$.
An element of $\bar{B}_1^{\mathrm{geom}}(\mathcal{H}_k)$ has
the form
$\alpha(z_1) \otimes \alpha(z_2) \otimes f(z_1,z_2)\,\eta_{12}$,
and the Fourier expansion
$f(z_1,z_2) = \sum_{n,m \in \mathbb{Z}}
 c_{n,m}\, e^{2\pi i(nz_1/\tau + mz_2)}$
gives the bar differential
$d(\alpha \otimes \alpha \otimes f\,\eta_{12})
= k \sum_n c_{n,n}$.
The constraint $\sum_n c_{n,n} = 0$ defines the kernel, and the
quotient by the image identifies with sections of $\mathcal{L}_k$
on $T^2 = \mathrm{Jac}(\Sigma_1)$.
Under $\tau \to -1/\tau$ the $A$- and $B$-cycles exchange,
and the Jacobi identity
$\theta_3(-1/\tau) = (-i\tau)^{1/2}\,\theta_3(\tau)$
gives the modular transformation of part~(1).
\end{proof}

\begin{example}[Explicit genus-1 computation]\label{ex:heisenberg-genus1-explicit}
For the torus $T^2$ with modulus $\tau$, we compute the bar complex explicitly in low degrees:

\emph{Degree~0.}
\[\bar{B}_0 = \mathbb{C} \quad \text{(vacuum)}\]

\emph{Degree~1.}
\[\bar{B}_1 = \bigoplus_{(n,m) \in \mathbb{Z}^2} \mathbb{C} \cdot a_{n,m}\]
where $a_{n,m}$ represents the mode $a(z) e^{2\pi i(n\text{Re}(z) + m\text{Im}(z))}$.

\emph{Degree~2.}
\[\bar{B}_2 = \bigoplus_{(n_1,m_1,n_2,m_2)} \mathbb{C} \cdot (a_{n_1,m_1} \otimes a_{n_2,m_2}) \cdot \eta_{12}\]

The differential is:
\[d(a_{n_1,m_1} \otimes a_{n_2,m_2} \cdot \eta_{12}) = k\delta_{n_1+n_2,0} \delta_{m_1+m_2,0}\]

\emph{Cohomology.}
\[H^1(\bar{B}^{\text{geom}}_1(\mathcal{H}_k)) = \frac{\ker(d: \bar{B}_1 \to \bar{B}_2)}{\text{im}(d: \bar{B}_0 \to \bar{B}_1)}\]

This computes the homology of the Jacobian variety $\text{Jac}(T^2) \cong T^2$ with the level-$k$ structure.

\emph{Modular transformation.} The S-transformation $\tau \to -1/\tau$ acts on modes:
\[a_{n,m} \to a_{m,-n}\]

This exchanges the roles of $n$ and $m$, swapping A-cycles and
B-cycles. In the bar complex, this induces the modular pullback
\[
\bar{B}^{\text{geom}}_1(\mathcal{H}_k,\tau)
\xrightarrow{S}
\bar{B}^{\text{geom}}_1(\mathcal{H}_k,-1/\tau)
\]
with the same OPE level~$k$. The complementary curvature
$-k$ belongs to the Koszul dual algebra
$\cH_k^! = \mathrm{Sym}^{\mathrm{ch}}(V^*)$, not to the
modular transform of $\cH_k$.
\end{example}

\begin{remark}[Physical interpretation: quantum complementarity]\label{rem:heisenberg-physics}
At genus~$1$, complementarity is electromagnetic duality on $T^2$:
A-cycle winding exchanges with B-cycle winding, and the
$S$-transformation $\Omega \to -1/\Omega$ acts as Fourier transform
on $\mathrm{Jac}(T^2)$. The OPE level $k$ is the curvature
coefficient of the Heisenberg algebra. Its sign reversal occurs on
the Verdier-linear Koszul-dual side, not under the modular
$S$-action on the same Heisenberg algebra.
\end{remark}

\begin{remark}[Koszul duality as higher-dimensional Fourier transform]
\label{rem:koszul-fourier}
Chiral Koszul duality is a higher-dimensional Fourier transform: the bar construction
\[
\bar{B}(\cA) \;=\; \bigoplus_n
 \int_{C_n(X)} \cA^{\boxtimes n}
 \otimes \Omega^*_{\log}(\overline{C}_n(X))
\]
with kernel $\eta_{ij} = d\log(z_i - z_j)$ produces the Koszul
coalgebra of $\cA$. Verdier duality gives the homotopy Koszul-dual
factorization algebra
$\mathbb{D}_{\operatorname{Ran}}\,\bar{B}(\cA)\simeq\cA^!_\infty$
(Theorem~\ref{thm:bar-cobar-isomorphism-main}), while the separate
cobar counit $\Omega\bar{B}(\cA)\simeq\cA$ is bar-cobar inversion.
For commutative algebras, the $n$-point contributions factor pairwise,
recovering a linear integral transform.

For $\cH_k$ on $\Sigma_g$, this reduction is literal on the
abelian zero modes: bar cohomology computes
$H^*(\mathrm{Jac}(\Sigma_g), \mathcal{L}_k)$, and Poisson summation
on $\mathbb{Z}^g$ gives Fourier inversion for theta sections. The
Verdier-linear Koszul dual is the separate curved algebra
$\cH_k^!=(\mathrm{Sym}^{\mathrm{ch}}(V^*),m_0=-k\omega)$; the
modular formula $\Omega\mapsto-\Omega^{-1}$ acts on the Jacobian
variables, not by replacing $\cH_k$ with $\cH_{-k}$.
\end{remark}

\subsection{From Fourier duality to chiral Koszul duality}
\label{subsec:fourier-to-koszul}
\index{Fourier duality!chiral generalization}
\index{Koszul duality!as Fourier transform}

Remark~\ref{rem:koszul-fourier} identified the chiral bar construction
as a nonlinear generalization of the classical Fourier transform. We
now develop this identification in three stages of increasing
nonlinearity: the classical abelian transform, its higher-dimensional
extension to Jacobians of curves, and the genuinely non-abelian
generalization to configuration space integrals with non-factorizing
kernel. The passage between stages is governed by a single datum:
the order of the OPE pole.

\subsubsection{Stage~I: pairwise factorization and the abelian kernel}

\begin{proposition}[Abelian factorization of the bar differential; \ClaimStatusProvedHere]
\label{prop:abelian-bar-factorization}
\index{bar complex!abelian factorization}
Let\/ $\cA$ be a chiral algebra on a smooth curve $X$ whose OPE
contains no simple pole:
\[
\phi^a(z)\phi^b(w)
= \frac{C^{ab}}{(z-w)^2} + \mathrm{reg}.
\]
Then the bar differential on
$\cA^{\boxtimes n} \otimes \Omega^p_{\log}(\overline{C}_n(X))$
decomposes as a sum of commuting pairwise operators:
\[
d_{\mathrm{bar}} = \sum_{1 \leq i < j \leq n} d_{ij},
\qquad
[d_{ij},\, d_{kl}] = 0
\;\;\text{for}\;\;
\{i,j\} \cap \{k,l\} = \emptyset.
\]
Each $d_{ij}$ extracts the second-order residue at the collision
divisor $D_{ij} \subset \overline{C}_n(X)$, acting only on the
tensor factors $\cA^{\boxtimes\{i,j\}}$ and the logarithmic
form~$\omega$.
\end{proposition}

\begin{proof}
On an element
$\alpha = \phi^{a_1}(z_1) \otimes \cdots \otimes \phi^{a_n}(z_n)
\otimes \omega$,
the bar differential extracts OPE residues at each collision
divisor. When the OPE has no simple pole, the only contribution
from $D_{ij}$ is the double-pole coefficient:
\[
d_{ij}(\alpha)
= C^{a_i a_j}
 \cdot \bigl(\text{contraction at } z_i = z_j\bigr)
 \cdot \mathrm{Res}_{D_{ij}}(\omega).
\]
This operator depends only on the tensor factors at positions $i$
and $j$ and on the restriction of $\omega$ to $D_{ij}$. For
disjoint index pairs $\{i,j\} \cap \{k,l\} = \emptyset$, the
operators $d_{ij}$ and $d_{kl}$ act on independent tensor factors
and extract residues at transverse divisors in
$\overline{C}_n(X)$, so they commute.
\end{proof}

\begin{remark}[From pairwise decomposition to classical Fourier analysis]
\label{rem:pairwise-to-fourier}
The pairwise decomposition (Proposition~\ref{prop:abelian-bar-factorization}) is why abelian chiral algebras produce linear integral transforms: the multivariate Fourier kernel factors as $e^{i\langle x,\xi\rangle} = \prod_j e^{ix_j\xi_j}$, and the chiral bar differential on $\cH_k$ reduces to iterated pairwise residue extraction, which on $\Sigma_g$ becomes Fourier coefficient extraction on $\mathrm{Jac}(\Sigma_g)$.
\end{remark}

\begin{remark}[Convolution--multiplication exchange]
\label{rem:convolution-exchange}
\index{convolution!chiral bar construction}
The bar functor exchanges the chiral product $\mu\colon \cA \boxtimes \cA \to \Delta_* \cA$ with the coproduct
\[
\Delta\colon \bar{B}^{\mathrm{ch}}_n(\cA)
\;\longrightarrow\;
\bigoplus_{p+q=n} \bar{B}^{\mathrm{ch}}_p(\cA)
 \otimes \bar{B}^{\mathrm{ch}}_q(\cA),
\]
defined by restriction to the boundary strata $D_{S,T}$ of $\overline{C}_n(X)$, the chiral analogue of $\widehat{f * g} = \hat{f} \cdot \hat{g}$. For $\cH_k$ on $\Sigma_g$, this identifies the chiral product with convolution on the Jacobian (Theorem~\ref{thm:heisenberg-genus-g}).
\end{remark}

The parallel between classical Fourier duality and chiral bar-cobar
duality is summarized in the following dictionary.

\medskip
\begin{center}
\small
\renewcommand{\arraystretch}{1.35}
\begin{tabular}{@{}lll@{}}
\hline
\textbf{Classical Fourier} & \textbf{Chiral bar-cobar}
 & \textbf{Abelian specialization} \\
\hline
Dual group $\hat{G}$
 & Koszul dual $\cA^!$
 & $\cH_k^! = \mathrm{Sym}^{\mathrm{ch}}(V^*)$ \\[2pt]
Kernel $e^{i\langle x,\xi\rangle}$
 & $\eta_{ij} = d\log(z_i - z_j)$
 & genus $0$: $dz/(z-w)$ \\[2pt]
$2$-point integral
 & $n$-point configuration integral
 & $n$-pt $= \tbinom{n}{2}$ pairwise \\[2pt]
$\widehat{f * g} = \hat{f}\hat{g}$
 & $\bar{B}(\mu) = \Delta_{\bar{B}}$
 & convolution $\leftrightarrow$ product \\[2pt]
Fourier inversion
 & $\mathbb{D}_{\mathrm{Ran}}\, \bar{B}(\cA)
 \simeq \cA^!_\infty$
 & $\Omega \mapsto -\Omega^{-1}$ \\[2pt]
Plancherel measure
 & spectral discriminant $\Delta_\cA(x)$
 & $\Delta_{\cH} = 1$ (flat) \\[2pt]
Poisson summation
 & theta function identity
 & Siegel modular involution \\[2pt]
$L^1(G)$ (convolution)
 & $U^{\mathrm{ch}}(\cA)$ (chiral envelope)
 & lattice algebra $\C[\Z^g]$ \\
\hline
\end{tabular}
\end{center}
\medskip

\subsubsection{Stage~II: the non-abelian departure}

The pairwise factorization of
Proposition~\ref{prop:abelian-bar-factorization} fails when the OPE
contains a simple pole. This failure is the structural origin of
the genuinely nonlinear phenomena in the bar complex of non-abelian
chiral algebras.

\begin{proposition}[Non-abelian kernel non-factorization; \ClaimStatusProvedHere]
\label{prop:nonabelian-kernel-nonfactorization}
\index{bar complex!non-abelian kernel}
Let\/ $\cA$ be a chiral algebra with OPE
\[
\phi^a(z)\phi^b(w)
= \frac{C^{ab}}{(z-w)^2}
 + \frac{f^{ab}_{\;\;c}\, \phi^c(w)}{z-w}
 + \mathrm{reg},
\qquad f^{ab}_{\;\;c} \neq 0.
\]
Then the bar differential decomposes as
\begin{equation}\label{eq:bar-diff-pair-bracket}
d_{\mathrm{bar}} = d_{\mathrm{pair}} + d_{\mathrm{bracket}},
\end{equation}
where $d_{\mathrm{pair}} = \sum_{i<j} d_{ij}^{(2)}$ extracts
double-pole residues and
$d_{\mathrm{bracket}} = \sum_{i<j} d_{ij}^{(1)}$ extracts
simple-pole residues. Then:
\begin{enumerate}[label=\textup{(\roman*)}]
\item For $n \geq 3$, the operator $d_{\mathrm{bracket}}$ on
 $n$-point elements is irreducibly $n$-body: the output
 $d_{ij}^{(1)}(\alpha) \in \bar{B}^{\mathrm{ch}}_{n-1}(\cA)$
 depends on all fields, including the pair
 $(\phi^{a_i}, \phi^{a_j})$.
\item $d_{\mathrm{bracket}}^2 \neq 0$ in general. The full
 identity $d_{\mathrm{bar}}^2 = 0$ requires the cross-terms
 $d_{\mathrm{pair}} \circ d_{\mathrm{bracket}}
 + d_{\mathrm{bracket}} \circ d_{\mathrm{pair}}
 = -d_{\mathrm{bracket}}^2$
 and relies on the Borcherds identity mixing pairwise and
 bracket contributions.
\end{enumerate}

In the language of homotopy chiral algebras
(Theorem~\ref{thm:cech-hca}), the identity
$d_{\mathrm{bracket}}^2 \neq 0$ is the obstruction to strict
chirality, and the cross-terms
$d_{\mathrm{pair}} \circ d_{\mathrm{bracket}} +
d_{\mathrm{bracket}} \circ d_{\mathrm{pair}}$ are the
secondary Borcherds operations $j'_{(p,q,r)}$
(Definition~\ref{def:secondary-borcherds}) that restore
$d_{\mathrm{bar}}^2 = 0$ at the Ch$_\infty$ level.
\end{proposition}

\begin{proof}
The simple-pole extraction at $D_{ij}$ produces
\[
d_{ij}^{(1)}(\phi^{a_1} \otimes \cdots \otimes \phi^{a_n}
 \otimes \omega)
= f^{a_i a_j}_{\;\;\;\;c}\;
 \bigl(\phi^{a_1} \otimes \cdots \otimes
 \underset{\text{pos.\ }i}{\phi^c} \otimes \cdots
 \otimes \widehat{\phi^{a_j}} \otimes \cdots
 \otimes \phi^{a_n}\bigr)
 \otimes \mathrm{Res}_{D_{ij}}(\omega).
\]
The contracted field $\phi^c$ retains position~$z_i$ and depends on
both $a_i$ and $a_j$ through $f^{a_i a_j}_{\;\;\;\;c}$, while the
remaining fields persist at their original positions. The output is
an $(n{-}1)$-point element whose dependence on all remaining fields
is irreducible: $d_{ij}^{(1)}$ cannot be expressed as an operator on
the pair $(\phi^{a_i}, \phi^{a_j})$ alone, since the contracted
generator $\phi^c$ interacts with $\phi^{a_k}$ at subsequent bar
degrees. This establishes~(i).

For~(ii), consider $n = 3$ and compute:
\begin{align*}
d_{23}^{(1)} \circ d_{12}^{(1)}
 (\phi^a \otimes \phi^b \otimes \phi^c
 \otimes \eta_{12} \wedge \eta_{23})
&= f^{ab}_{\;\;e}\, f^{ec}_{\;\;g}\, \phi^g, \\
d_{12}^{(1)} \circ d_{23}^{(1)}
 (\phi^a \otimes \phi^b \otimes \phi^c
 \otimes \eta_{12} \wedge \eta_{23})
&= f^{bc}_{\;\;e}\, f^{ae}_{\;\;g}\, \phi^g.
\end{align*}
The sum $d_{\mathrm{bracket}}^2
= f^{ab}_{\;\;e}\, f^{ec}_{\;\;g}
+ f^{bc}_{\;\;e}\, f^{ae}_{\;\;g}
+ f^{ca}_{\;\;e}\, f^{be}_{\;\;g}$
vanishes by the Jacobi identity, but the cancellation requires the
additional contribution from $d_{13}^{(1)}$ and the double-pole
cross-terms. In isolation, $d_{\mathrm{bracket}}^2 \neq 0$: the
identity $d_{\mathrm{bar}}^2 = 0$ is a \emph{coupled} identity
between the Lie bracket and the level, not a consequence of either
alone.
\end{proof}

The non-abelian departure is sharpest in the following side-by-side
comparison.

\begin{example}[Heisenberg versus Kac--Moody at bar degree~\texorpdfstring{$2$}{2}]
\label{ex:heisenberg-vs-km}
\index{bar complex!Heisenberg vs.\ Kac--Moody}

\emph{Heisenberg.}
The OPE $\alpha(z)\alpha(w) = k/(z{-}w)^2 + \mathrm{reg}$ has no
simple pole. By Proposition~\ref{prop:abelian-bar-factorization},
the differential on $3$-point elements is purely pairwise:
\[
d(\alpha_1 \otimes \alpha_2 \otimes \alpha_3
 \otimes \eta_{12} \wedge \eta_{23})
= k\, \alpha_3 \otimes \eta_{23}
 \;-\; k\, \alpha_1 \otimes \eta_{13}.
\]
Each term involves a single residue extraction and a single
surviving field. The two terms cancel under $d^2$ by telescoping
(Theorem~\ref{thm:heisenberg-bar-complete}).

\emph{Kac--Moody $\AffKM{g}_k$.}
The OPE has both poles, so the differential acquires a non-abelian
correction:
\begin{align*}
d(J^a \otimes J^b \otimes J^c
 &\otimes \eta_{12} \wedge \eta_{23}) \\
&= \underbrace{k\delta^{ab}\, J^c \otimes \eta_{23}
 \;-\; k\delta^{bc}\, J^a
 \otimes \eta_{12}}_{\displaystyle
 d_{\mathrm{pair}}\;\text{(level, pairwise, ``Fourier'')}} \\
&\quad+\; \underbrace{f^{ab}_{\;\;e}\, J^e \otimes J^c
 \otimes \eta_{23}
 \;+\; f^{bc}_{\;\;e}\, J^a \otimes J^e
 \otimes \eta_{12}}_{\displaystyle
 d_{\mathrm{bracket}}\;\text{(structure
 constants, $3$-body, non-abelian)}}.
\end{align*}
The pairwise terms yield $1$-point elements (scalar times a single
current), as in the Heisenberg case. The bracket terms yield
\emph{$2$-point} elements carrying the Lie structure constants
$f^{abc}$; these have no abelian counterpart and encode the failure
of the bar kernel to factorize.

The identity $d^2 = 0$ on the bracket terms encodes the Jacobi
identity
$f^{ab}_{\;\;e}\, f^{ec}_{\;\;g}
+ f^{bc}_{\;\;e}\, f^{ae}_{\;\;g}
+ f^{ca}_{\;\;e}\, f^{be}_{\;\;g} = 0$:
the three cyclic contributions from $d_{\mathrm{bracket}}^2$ cancel
by antisymmetry. This is the sense in which the Jacobi identity is
the \emph{non-abelian Arnold relation}: it is the algebraic
constraint imposed by $d^2 = 0$ on the configuration space
$\overline{C}_3(X)$.
\end{example}

\begin{example}[Explicit bar differential matrix for
 \texorpdfstring{$\AffKM{sl}_{2,k}$}{sl_2,k}]
\label{ex:sl2-bar-matrix}
\index{bar complex!$\AffKM{sl}_2$ explicit matrix}

We make Example~\ref{ex:heisenberg-vs-km} fully computational for
$\AffKM{sl}_{2,k}$. The Lie algebra $\mathfrak{sl}_2$ has basis
$\{H, E, F\}$ with $[H, E] = 2E$, $[H, F] = -2F$, $[E, F] = H$
and normalized trace $(H, H) = 2$, $(E, F) = 1$. The OPE is
\begin{align*}
J^H(z)\, J^H(w) &= \frac{2k}{(z-w)^2}, &
J^H(z)\, J^E(w) &= \frac{2J^E(w)}{z-w}, \\
J^H(z)\, J^F(w) &= \frac{-2J^F(w)}{z-w}, &
J^E(z)\, J^F(w) &= \frac{k}{(z-w)^2} + \frac{J^H(w)}{z-w}.
\end{align*}

\emph{Degree-$1$ bar group} (two-point elements).
The space $\barBgeom^1$ is spanned by
$J^a \otimes J^b \otimes \eta_{12}$ with
$a, b \in \{H, E, F\}$: nine generators. The bar differential
extracts the full OPE:
\[
d(J^a \otimes J^b \otimes \eta_{12})
= \underbrace{k(a,b)}_{\text{double pole}} \cdot \mathbf{1}
 \;+\; \underbrace{[a,b]}_{\text{simple pole}}.
\]
Explicitly, splitting $d = d_{\mathrm{pair}} + d_{\mathrm{bracket}}$:

\medskip
\begin{center}
\small
\renewcommand{\arraystretch}{1.2}
\begin{tabular}{@{}l@{\;\;\;}c@{\;\;\;}c@{}}
\hline
Element $\otimes\, \eta_{12}$ & $d_{\mathrm{pair}}$ (level)
 & $d_{\mathrm{bracket}}$ (structure) \\
\hline
$J^H \otimes J^H$ & $2k \cdot \mathbf{1}$ & $0$ \\
$J^H \otimes J^E$ & $0$ & $2J^E$ \\
$J^H \otimes J^F$ & $0$ & $-2J^F$ \\
$J^E \otimes J^H$ & $0$ & $-2J^E$ \\
$J^E \otimes J^E$ & $0$ & $0$ \\
$J^E \otimes J^F$ & $k \cdot \mathbf{1}$ & $J^H$ \\
$J^F \otimes J^H$ & $0$ & $2J^F$ \\
$J^F \otimes J^E$ & $k \cdot \mathbf{1}$ & $-J^H$ \\
$J^F \otimes J^F$ & $0$ & $0$ \\
\hline
\end{tabular}
\end{center}
\medskip

In matrix form, with source basis
$(J^{HH},\, J^{HE},\, J^{HF},\, J^{EH},\, J^{EE},\, J^{EF},\,
J^{FH},\, J^{FE},\, J^{FF}) \otimes \eta_{12}$
and target basis
$(\mathbf{1},\; J^H,\; J^E,\; J^F)$:
\begin{equation}\label{eq:sl2-bar-matrix-fourier}
d_{\mathrm{bar}}
= \begin{pmatrix}
2k & 0 & 0 & 0 & 0 & k & 0 & k & 0 \\
0 & 0 & 0 & 0 & 0 & 1 & 0 & -1 & 0 \\
0 & 2 & 0 & -2 & 0 & 0 & 0 & 0 & 0 \\
0 & 0 & -2 & 0 & 0 & 0 & 2 & 0 & 0
\end{pmatrix}.
\end{equation}

\emph{Rank and kernel analysis.}
\begin{enumerate}[label=\textup{(\roman*)}]
\item \emph{Generic level} ($k \neq 0$):
 $\operatorname{rank}(d_{\mathrm{bar}}) = 4$
 (the four rows are independent: row~$1$ is the unique row
 hitting column~$1$). The kernel is
 $5$-dimensional, spanned by:
 \begin{gather*}
 J^E \otimes J^E, \quad
 J^F \otimes J^F, \quad
 J^H \otimes J^E + J^E \otimes J^H, \\
 J^H \otimes J^F + J^F \otimes J^H, \quad
 -J^H \otimes J^H + J^E \otimes J^F + J^F \otimes J^E.
 \end{gather*}
 The first four are symmetric tensors in $\operatorname{Sym}^2(\mathfrak{sl}_2)$; the
 fifth is the negative Casimir element $-C_2 \in
 (\mathfrak{sl}_2 \otimes \mathfrak{sl}_2)^{\mathfrak{sl}_2}$.
 These five elements are bar cocycles, the ``harmonic modes''
 that survive the non-abelian Fourier transform.

\item \emph{Abelian limit} ($k = 0$):
 $\operatorname{rank}(d_{\mathrm{bar}}) = 3$ (row~$1$ vanishes).
 The kernel gains one dimension: the Casimir class
 $\kappa = J^H \otimes J^H / 2 + J^E \otimes J^F + J^F \otimes J^E$
 becomes a cocycle, since its image under $d$ is
 $3k \cdot \mathbf{1} = 0$. This extra mode is the
 zero-frequency Fourier coefficient of the Killing form.
\end{enumerate}
The Heisenberg algebra at the same degree has a $1 \times 1$ matrix
$(k)$, with $0$-dimensional kernel for $k \neq 0$. The passage from
a $1 \times 1$ scalar to the $4 \times 9$
matrix~\eqref{eq:sl2-bar-matrix} is the passage from abelian Fourier
(one frequency, one coefficient) to non-abelian Fourier ($\dim(\mathfrak{g})^2 = 9$ frequency channels, $\dim(\mathfrak{g}) + 1 = 4$
coefficient slots, and a $5$-dimensional kernel of harmonic modes
governed by the representation theory of~$\mathfrak{g}$).
\end{example}

\subsubsection{Stage~III: non-abelian Fourier symmetry and spectral
 measure}

\begin{remark}[The Weyl group as non-abelian Poisson summation]
\label{rem:weyl-nonabelian-fourier}
\index{Weyl group!as Fourier symmetry}
\index{Poisson summation!non-abelian analogue}
The Weyl--Kac character formula
\textup{(}Theorem~\textup{\ref{thm:weyl-kac-geometric}}\textup{)}
replaces the dual-lattice sum in classical Poisson summation with an
alternating sum over $W_{\mathrm{aff}} = W \ltimes Q^\vee$. The bar
spectral sequence realizes the full generic/admissible passage:
$E_1$ is indexed by Verma modules $\mathcal{M}(w \cdot \lambda)$,
collapsing at generic level by
Theorem~\ref{thm:weyl-kac-geometric} and reducing to a finite
admissible Weyl set on the worked $\widehat{\mathfrak{sl}}_2$
surface, while in general rank the corresponding finite-sum
character formula is the external
Kac--Wakimoto theorem
\textup{(}Theorem~\ref{thm:kw-bar-general-rank}\textup{)}.
\end{remark}

\begin{remark}[Spectral discriminant as Plancherel density]
\label{rem:spectral-plancherel}
\index{spectral discriminant!Plancherel density}
\index{Plancherel measure!chiral analogue}
The spectral discriminant $\Delta_\cA(x) = \prod_i (1 - \lambda_i x)$ (Theorem~\ref{thm:ds-bar-gf-discriminant}) is the chiral Plancherel density: $\Delta_\cH = 1$ (uniform, partition growth), $\Delta_{\beta\gamma} = (1 - 3x)(1 + x)$ (oscillatory, DS-stable), and zeros of $\Delta_{\AffKM{sl}_2}(x)$ detect singular vectors.
\end{remark}

\subsubsection{Stage~IV: the homotopy-theoretic uplift}

Stages~I--III operate at the M-level: explicit chain complexes,
logarithmic propagators, matrix differentials. The full picture
lives at the H-level, in the $\infty$-category of $\En$-algebras,
where the Fourier duality becomes Poincar\'e--Koszul duality and the
bar-cobar adjunction is a derived equivalence.

\begin{proposition}[\texorpdfstring{$\En$}{En} Fourier hierarchy; \ClaimStatusProvedHere]
\label{prop:en-fourier-hierarchy}
\index{Fourier duality!$\En$ hierarchy}
\index{propagator!$\En$ hierarchy}
The bar construction admits a hierarchy of integral transforms
indexed by the real dimension~$n$ of the underlying manifold.
At each level, the propagator, the relations ensuring $d^2 = 0$,
and the duality involution generalize the classical Fourier data:

\medskip
\begin{center}
\small
\renewcommand{\arraystretch}{1.3}
\begin{tabular}{@{}cllll@{}}
\hline
$n$ & Propagator & Relations & Duality & Physical setting \\
\hline
$1$ & $\delta(x_i - x_j) \in H^0$
 & associativity & linear duality
 & quantum mechanics \\[2pt]
$2$ & $\eta_{ij} = d\log(z_i - z_j) \in \Omega^1$
 & Arnold & Verdier on $\overline{C}_k(X)$
 & CFT / chiral algebras \\[2pt]
$3$ & $G = \varepsilon_{abc} r^a\, dr^b \wedge dr^c / 4\pi|r|^3
 \in \Omega^2$
 & Totaro & linking number
 & Chern--Simons \\[2pt]
$n$ & $(n{-}1)$-form on $\Conf_2(\R^n)$
 & Totaro & $\En$-Poincar\'e--Koszul
 & $\En$-TFT \\
\hline
\end{tabular}
\end{center}
\medskip

\noindent
In each case, the bar differential is
\[
d_{\En}
= d_{\mathrm{dR}}
+ \sum_{\substack{S \subset [k] \\ |S| \geq 2}}
\mathrm{Res}_{\partial_S}
\Bigl(\bigwedge_{i < j \in S} G_{ij} \wedge \mu_S\Bigr),
\]
where $G_{ij}$ is the $\En$ propagator
\textup{(}an $(n{-}1)$-form on $\Conf_2(\R^n)$,
Definition~\textup{\ref{def:en-propagator})},
$\mathrm{Res}_{\partial_S}$ is the linking-sphere residue
\textup{(}Proposition~\textup{\ref{prop:linking-sphere-residue})},
and $\mu_S$ is the $\En$-multiplication. The identity $d^2 = 0$
follows from the Totaro relations at dimension~$n$
\textup{(}Theorem~\textup{\ref{thm:totaro-presentation})},
which specialize to Arnold relations at $n = 2$
\textup{(}Theorem~\textup{\ref{thm:e2-d-squared})}.
\end{proposition}

\begin{proof}
The $\En$ propagator is a closed $(n{-}1)$-form representing the
Poincar\'e dual of the diagonal
(Definition~\ref{def:en-propagator}), with singularity
$\sim |r|^{1-n}\, \mathrm{vol}_{S^{n-1}}$ near collision. The
residue operation
$\mathrm{Res}_{\partial_S}\bigl(\bigwedge G_{ij} \wedge \mu_S\bigr)$
integrates $G$ over the linking sphere
$S^{n-1}_\varepsilon$ around the collision locus, extracting the
leading singularity of $\mu_S$
(Proposition~\ref{prop:linking-sphere-residue}).

At $n = 2$, the linking sphere $S^1$ is the contour
$|z_i - z_j| = \varepsilon$, the propagator
$G = d\log(z_i - z_j)$ is the holomorphic $1$-form, and the
residue is the classical OPE extraction. At $n = 3$, the linking
sphere $S^2$ computes linking numbers, with Green's two-form
$\frac{1}{4\pi}\epsilon_{abc}r^a\,dr^b\wedge dr^c/|r|^3$, and the
bar differential encodes the perturbative Chern--Simons invariants.

The identity $d^2 = 0$ at general~$n$ follows from
Theorem~\ref{thm:e2-d-squared} (proved there for $n = 2$ via
Arnold, extending to general~$n$ via Totaro): the
codimension-$2$ boundary strata of $\overline{\Conf}_k(\R^n)$
cancel in pairs by the Totaro three-term relation
$\alpha_{ij} \cdot \alpha_{jk} + \alpha_{jk} \cdot \alpha_{ki}
+ \alpha_{ki} \cdot \alpha_{ij} = 0$ in
$H^{2(n-1)}(\Conf_3(\R^n))$.
\end{proof}

The $\En$ Fourier hierarchy admits a fully derived formulation.

\begin{theorem}[Derived Fourier duality via Poincar\'e--Koszul
 duality; \ClaimStatusConditional]
\label{thm:derived-fourier}
\index{Poincar\'e--Koszul duality!as derived Fourier}
\index{Fourier duality!derived}
Let $A$ be an $\En$-algebra in chain complexes.
\begin{enumerate}[label=\textup{(\roman*)}]
\item \textup{(Lurie~\cite{HA}, \S5.2)}\quad
 The bar-cobar adjunction
 \[
 \Omega_{\En} \;\dashv\; \barB_{\En}
 \colon
 \En\text{-}\mathrm{CoAlg}_\infty
 \;\rightleftarrows\;
 \En\text{-}\mathrm{Alg}_\infty
 \]
 is an adjunction of\/ $\infty$-categories. On the
 $\En$-Koszul locus (where the bar spectral sequence
 collapses), the counit
 $\Omega_{\En}(\barB_{\En}(A)) \xrightarrow{\sim} A$
 is an equivalence
 \textup{(}Theorem~\textup{\ref{thm:en-koszul-duality})}.

\item \textup{(Ayala--Francis~\cite{AF15}, Theorem~7.8)}\quad
 For a closed oriented $n$-manifold~$M$,
 Poincar\'e--Koszul duality gives
 \begin{equation}\label{eq:pk-duality-fourier}
 \int_M A
 \;\simeq\;
 \int_M A^{\vee}[-n],
 \end{equation}
 where $\int_M$ denotes $\En$-factorization homology and
 $A^\vee$ is the $\En$-Koszul dual coalgebra
 \textup{(}Theorem~\textup{\ref{thm:af-pkd})}. This is the
 derived Fourier inversion formula: it exchanges the algebra~$A$
 (``position space'') with its Koszul dual~$A^\vee$ (``momentum
 space'') via factorization homology over~$M$ (``integration'').

\item \textup{(This monograph, $n = 2$)}\quad
 At $n = 2$, the equivalence~\eqref{eq:pk-duality-fourier}
 specializes to the Verdier identification
 $\mathbb{D}_{\mathrm{Ran}}\, \bar{B}(\cA) \simeq \cA^!_\infty$
 of Theorem~\textup{\ref{thm:bar-cobar-isomorphism-main}},
 realized at the chain level by propagator integrals
 \textup{(}Proposition~\textup{\ref{prop:refines-af})}.
 The holomorphic propagator $\eta_{ij} = d\log(z_i - z_j)$
 is a specific representative of the $\Etwo$ propagator class
 $[\alpha_{ij}] \in H^1(\Conf_2(X))$
 \textup{(}Corollary~\textup{\ref{cor:n2-recovery})}.
\end{enumerate}
\end{theorem}

\begin{remark}[Nonlinear Fourier interpretation]
\label{rem:fh-fourier}\label{rem:htt-fourier}\label{rem:curved-fourier}
\index{factorization homology!as Fourier integral}\index{homotopy transfer!Fourier interpretation}\index{coderived category!Fourier interpretation}
The derived Fourier transform (Theorem~\ref{thm:derived-fourier}) replaces $\int_G f\, e^{i\langle\cdot,\cdot\rangle}$ with factorization homology $\int_M A$, the $\En$ propagator serving as kernel. Homotopy transfer (Theorem~\ref{thm:htt}) yields an $\Ainf$-structure via the tree formula (Theorem~\ref{thm:tree-formula}); for Koszul algebras $\tilde{m}_k = 0$ for $k \geq 3$ (Theorem~\ref{thm:chiral-htt}). At genus $g \geq 1$, curvature forces passage to $D^{\mathrm{co}}(\barB_g(\cA))$, where the derived inversion~\eqref{eq:pk-duality-fourier} holds. Complementarity supplies the separate scalar statement
\[
\mathrm{obs}_g(\cA)+\mathrm{obs}_g(\cA^!)
=\bigl(\kappa(\cA)+\kappa(\cA^!)\bigr)\lambda_g,
\]
with the value of $\kappa+\kappa^!$ determined family by family.
\end{remark}

\subsection{Explicit bar complex calculation}

The bar complex $\bar{B}^{\text{geom}}_*(\mathcal{H}_k)$ through degree 5:

\begin{theorem}[Heisenberg geometric bar differential; \ClaimStatusProvedHere]\label{thm:heisenberg-bar-complete}
For the Heisenberg algebra $\mathcal{H}_k$ on a curve $X$, the bar complex has chain groups
$\bar{B}_0 = \mathbb{C}$,
$\bar{B}_1 = \mathcal{H}_k$,
$\bar{B}_n = \mathcal{H}_k^{\otimes n} \otimes \Omega^*_{\log}(\overline{C}_n(X))$
for $n \geq 2$, with differential
$d = d_{\text{int}} + d_{\text{res}} + d_{dR}$
satisfying $d^2 = 0$.
\end{theorem}

\begin{proof}
\emph{Degree-by-degree structure.}

\begin{align}
\bar{B}_0 &= \mathbb{C} \quad \text{(vacuum)} \\
\bar{B}_1 &= \mathcal{H}_k \quad \text{(the algebra itself)} \\
\bar{B}_2 &= \mathcal{H}_k \otimes \mathcal{H}_k \otimes \Omega^1_{\log}(\overline{C}_2(X)) \\
\bar{B}_3 &= \mathcal{H}_k^{\otimes 3} \otimes \Omega^*_{\log}(\overline{C}_3(X)) \\
&\vdots
\end{align}

\emph{Differential structure.}
The differential has three components:
\[d = d_{\text{int}} + d_{\text{res}} + d_{dR}\]

For genus 0:
\begin{itemize}
\item $d_{\text{int}} = 0$ (Heisenberg has no internal operations beyond the bilinear bracket)
\item $d_{\text{res}}$: Extracts residues at collision divisors using the OPE coefficient $k$
\item $d_{dR}$: de Rham differential on logarithmic forms
\end{itemize}

\emph{Explicit formulas through degree~3.}

\[d_{\text{res}}(a(z_1) \otimes a(z_2) \otimes \eta_{12}) = k \cdot 1\]

\[d_{\text{res}}(a(z_1) \otimes a(z_2) \otimes a(z_3) \otimes \eta_{12} \wedge \eta_{23}) = k \cdot a(z_3) \otimes \eta_{23} - k \cdot a(z_1) \otimes \eta_{13}\]

We verify $d^2 = 0$ by direct computation. The two terms in $d(a \otimes a \otimes a \otimes \eta_{12} \wedge \eta_{23})$ yield, upon a second application of $d$:
\[d(k \cdot a(z_3) \otimes \eta_{23}) = k^2, \qquad d(-k \cdot a(z_1) \otimes \eta_{13}) = -k^2\]
so $d^2 = k^2 - k^2 = 0$ by telescoping cancellation of the two residue terms.
\end{proof}

\begin{remark}[Comparison with literature]\label{rem:heisenberg-literature}
This calculation agrees with:
\begin{itemize}
\item \emph{Gui--Li--Zeng \cite{GLZ22}.} Their Theorem 4.2 for Heisenberg specializes to the formulas above.
\item \emph{Beilinson--Drinfeld \cite{BD04}.} Section 4.7 on chiral homology, specialized to Heisenberg.
\item \emph{Costello--Gwilliam \cite{CG17}.} Volume 2, Chapter 5 on factorization algebras for Heisenberg.
\end{itemize}

The agreement provides non-trivial verification of the geometric approach via configuration spaces.
\end{remark}

\begin{lemma}[Bar dimensions as partition numbers; \ClaimStatusProvedHere]%
\label{lem:bar-dims-partitions}
\index{partition function!bar complex dimensions}
The reduced weight-window chain dimensions for the free-field bar
complexes are governed by integer partitions. These are chain-level
counts before taking the bar differential; they are not the
genus-$0$ bar cohomology groups. For the free fermion $\psi$ (one
fermionic generator of conformal weight~$\frac{1}{2}$):
\begin{equation}\label{eq:fermion-bar-partitions}
\dim \bar{B}^n(\psi) \;=\; p(n-1), \qquad n \geq 1,
\end{equation}
where $p(m)$ is the number of partitions of~$m$
($p(0) = 1$, $p(1) = 1$, $p(2) = 2$, $p(3) = 3$, $p(4) = 5$).
This gives the master-census sequence $1, 1, 2, 3, 5$ at the
chain-counting level.
The identity~\eqref{eq:fermion-bar-partitions} follows from
$\bar{B}^n(\psi) \cong \Lambda^n_{\mathrm{co}}(\psi)$
(cofree cocommutative coalgebra on the odd mode system):
the Arnold relations on $\overline{C}_{n+1}(\mathbb{P}^1)$ act
trivially on this exterior coalgebra, and enumerating the basis
elements by conformal weight decomposition yields $p(n-1)$
chain generators at bar degree~$n$. The cohomology of the same
genus-$0$ complex is the collapse stated in
Theorem~\ref{thm:fermion-bar-complex-genus-0}.

For the Heisenberg algebra $\mathcal{H}_k$ (one bosonic generator
of conformal weight~$1$):
\begin{equation}\label{eq:heisenberg-bar-partitions}
\dim \bar{B}^n(\mathcal{H}_k) \;=\;
\begin{cases}
1 & n = 1, \\
p(n-2) & n \geq 2,
\end{cases}
\end{equation}
giving the Master Table sequence $1, 1, 1, 2, 3$.
The shift from $p(n-1)$ to $p(n-2)$ reflects the higher
conformal weight of the Heisenberg generator ($h = 1$ versus
$h = \frac{1}{2}$): the desuspension $s^{-1}\bar{\mathcal{H}}$
has generators starting at weight~$0$, and the Arnold quotient
eliminates one degree of freedom per tensor factor compared to
the fermionic case. Equivalently,
$\bar{B}^n(\mathcal{H}) \cong \mathrm{Sym}^n_{\mathrm{co}}(J)$
(cofree cocommutative coalgebra on one even generator), and the
graded dimension of the $n$-th symmetric power is $p(n-2)$
for $n \geq 2$.

The generating functions:
\[
\sum_{n \geq 1} \dim \bar{B}^n(\psi)\, t^n
= \frac{t}{\prod_{m=1}^\infty (1-t^m)},
\qquad
\sum_{n \geq 1} \dim \bar{B}^n(\mathcal{H})\, t^n
= t + \frac{t^2}{\prod_{m=1}^\infty (1-t^m)}.
\]
\end{lemma}

\begin{remark}[Partition bridge to Yangian prefundamentals]
The same partition sequence governs the weight multiplicities
\(\dim(L^-_a)_{-2k} = p(k)\) of the Yangian negative
prefundamental modules
\textup{(}Theorem~\textup{\ref{thm:shifted-prefundamental-generation}},
Step~2\textup{)}. This gives a partition-theoretic bridge between the
free-field bar complex and the type-$A$ Yangian prefundamental sector
appearing in the shifted-prefundamental-generation theorem. By
itself, this partition match does not prove the full compact-completed
identification. Thick generation on the evaluation-generated core is
proved for all simple types by Corollary~\ref{cor:mc3-all-types};
extension beyond evaluation modules is a separate theorem.
\end{remark}

\subsection{\texorpdfstring{Curved duality under level inversion $k \mapsto -k$}{Curved duality under level inversion k -k}}

\begin{theorem}[Heisenberg level inversion: curved duality; \ClaimStatusProvedHere]\label{thm:heisenberg-level-inversion}
The Heisenberg algebras at levels $k$ and $-k$ form a \emph{curved/filtered dual pair} (distinct from standard Koszul duality) with:
\begin{enumerate}
\item Curvature terms:
$m_0^{(k)}=k\,\omega_{\mathrm{cen}}$.
\item Modified pairing: $\langle J \otimes J, J \otimes J \rangle_k = k \cdot \delta^{(2)}(z-w)$
\item Curved cobar complexes related by:
$\Omega^{\mathrm{curv}}\bar{B}(\mathcal{H}_k)\cong
\Omega^{\mathrm{curv}}\bar{B}(\mathcal{H}_{-k})$ as underlying
filtered vector spaces with opposite curvature.
\end{enumerate}

This is \emph{not} the same as the standard Koszul duality $\mathcal{H}^! = \mathrm{Sym}^{\mathrm{ch}}(V^*)$ established above; it is an additional duality structure that exists in the curved/filtered category.
\end{theorem}

\begin{proof}
The Heisenberg OPE has no simple-pole bracket term, but its
double-pole coefficient is exactly the trace-form collision datum:
\[
r_{\cH_k}^{\mathrm{coll}}(z)=\frac{k}{z}.
\]
Equivalently, the Borcherds identity extracts the second-order pole
coefficient as the boundary pairing
\[
\langle J \otimes J, J \otimes J \rangle_k = k
\]
(the level appears as the second-order pole coefficient, and after
the $d\log$ convention as the simple rational residue $k/z$).

Under $k \mapsto -k$, this pairing changes sign, establishing curved
level inversion. The reduced bar coalgebras have the same underlying
conilpotent coalgebra; the curvature on the curved cobar side changes
sign:
\begin{itemize}
\item $B(\mathcal{H}_k)=T^c(s^{-1}\bar{\mathcal{H}}_k)$ with
the same deconcatenation coalgebra for $k$ and $-k$;
\item the simple-pole bracket differential is zero in both cases;
\item the curved cobar term is
$m_0^{(k)}=k\,\omega_{\mathrm{cen}}$ and
$m_0^{(-k)}=-k\,\omega_{\mathrm{cen}}$.
\end{itemize}
Thus the ordinary genus-$0$ bar differential still squares to zero,
while the curved cobar differential satisfies the Positselski relation
$m_1^2=[m_0,-]$.
\end{proof}

For the full Koszul duality table (including all entries formerly listed here), see \S\ref{sec:kd-computations-free-fields}.


\section{Free fermion genus expansion and spin structures}
\label{sec:free-fermion-genus-expansion}
\index{free fermion!genus expansion|textbf}
\index{spin structure!genus expansion}

The free fermion $\cF$ ($c = 1/2$, $\kappa(\cF)=1/4$) is
an example where spin structures enter the genus
expansion. The Heisenberg genus expansion
(Chapter~\ref{ch:heisenberg-frame}) involves determinants of
the $\bar\partial$ operator on the trivial line bundle; the free
fermion replaces determinants with Pfaffians and theta functions
with theta constants. The shadow obstruction tower prediction
$F_g(\cF)=F_g^{\mathrm{sc}}(\cF)
=\kappa(\cF) \cdot \lambda_g^{\mathrm{FP}}
= \frac{1}{4}\lambda_g^{\mathrm{FP}}$ remains valid
(Theorem~\ref{thm:genus-universality}), but the partition
function from which $F_g$ is extracted has a richer structure
than the Heisenberg case: it depends on a choice of spin
structure $\sigma$ on $\Sigma_g$.

\begin{center}
\renewcommand{\arraystretch}{1.3}
\begin{tabular}{ll}
\textbf{Algebra} & Free fermion~$\cF$,
OPE $\psi(z)\,\psi(w) \sim 1/(z{-}w)$ \\
\textbf{Central charge} & $c = 1/2$ \\
\textbf{Shadow archetype} & G (Gaussian), $r_{\max} = 2$ \\
\textbf{Modular characteristic} & $\kappa(\cF) = c/2 = 1/4$ \\
\textbf{Koszul dual} & $\cF^! = \mathrm{Sym}^{\mathrm{ch}}(\gamma)$
(single bosonic generator) \\
\textbf{Complementarity} & $\kappa + \kappa' = 1/4 + (-1/4) = 0$ \textup{(}free-field pair\textup{)} \\
\textbf{$r$-matrix} & $r(z) = 0$ (simple pole absorbed by $d\log$ extraction) \\
\textbf{Genus expansion} & $F_g = \tfrac{1}{4}\,\lambda_g^{\mathrm{FP}}$,
GF $= \tfrac{1}{4\hbar^2}(\hat{A}(i\hbar) - 1)$
\end{tabular}
\end{center}

\begin{remark}[$r$-matrix vanishing]
\label{rem:fermion-r-matrix-vanishing}
\index{free fermion!r-matrix@$r$-matrix}
The free fermion OPE has a \emph{simple} pole
$\psi(z)\psi(w) \sim 1/(z{-}w)$. The bar construction
extracts the collision residue along $d\log(z{-}w)$,
absorbing one power of $(z{-}w)$
(Remark~\ref{rem:heisenberg-bar-absorbs-pole}). The resulting
collision term is scalar, not field-valued; after projection to the
trace-form rational $r$-matrix it contributes no pole:
$r(z)=0$. This is consistent with the abelian nature of the fermion
algebra: the chiral bracket $\{-,-\}$ is the $(z{-}w)^{-1}$ residue
of the OPE, which for the free fermion is $1$ (a scalar central
term, not a generator-valued expression), so the corresponding
field-valued Lie bracket is trivial.
\end{remark}

\subsection{Spin structures on Riemann surfaces}
\label{subsec:spin-structures}
\index{spin structure|textbf}

A spin structure on a closed Riemann surface $\Sigma_g$ is a
square root $L$ of the canonical bundle: $L^{\otimes 2} \cong
\omega_{\Sigma_g}$. The conformal weight $h = 1/2$ of the
free fermion means that $\psi$ is a section of $L$, so the
definition of the theory on $\Sigma_g$ requires a choice
of spin structure.

\begin{proposition}[Spin structure count;
\ClaimStatusProvedElsewhere]
\label{prop:spin-structure-count}
\index{spin structure!count}
A genus-$g$ surface has exactly $2^{2g}$ spin structures,
classified by the mod-$2$ cohomology
$H^1(\Sigma_g, \mathbb{Z}/2)$. They decompose by parity
\textup{(}the Atiyah invariant, equal to
$\dim H^0(\Sigma_g, L) \bmod 2$\textup{)} into:
\begin{align}
\text{even:} &\quad 2^{g-1}(2^g + 1), \\
\text{odd:} &\quad 2^{g-1}(2^g - 1).
\end{align}
At genus~$1$: $4$ spin structures, $3$ even and $1$ odd.
At genus~$2$: $16$ spin structures, $10$ even and $6$ odd.
\end{proposition}

\begin{proof}[References]
The count is classical; see Atiyah~\cite{Atiyah71} for the
parity invariant and Mumford~\cite{Mum71} for the
relationship to theta characteristics. A spin structure on
$\Sigma_g$ is equivalent to a theta characteristic, i.e.\ a
divisor class $[D]$ with $2[D] \sim K_{\Sigma_g}$. The number
of theta characteristics is $|H^1(\Sigma_g,
\mathbb{Z}/2)| = 2^{2g}$, and the parity decomposition
follows from counting the zeros of the Riemann theta function
with half-integer characteristic.
\end{proof}

\subsection{Genus-1 partition function: theta functions and spin
structures}
\label{subsec:fermion-genus1}
\index{free fermion!genus-1 partition function}

At genus~$1$, the four spin structures are labeled by
half-integer characteristics
$[\alpha, \beta] \in (\mathbb{Z}/2)^2$:
\[
\begin{array}{lcl}
(\alpha,\beta) = (1,0) & \longleftrightarrow &
 \text{NS, antiperiodic--periodic
 (even, } \vartheta_2\text{)} \\[2pt]
(\alpha,\beta) = (0,0) & \longleftrightarrow &
 \text{NS, periodic--periodic
 (even, } \vartheta_3\text{)} \\[2pt]
(\alpha,\beta) = (0,1) & \longleftrightarrow &
 \text{NS, periodic--antiperiodic
 (even, } \vartheta_4\text{)} \\[2pt]
(\alpha,\beta) = (1,1) & \longleftrightarrow &
 \text{R, antiperiodic--antiperiodic
 (odd, } \vartheta_1 = 0\text{)}
\end{array}
\]

\begin{theorem}[Free fermion genus-1 partition functions;
\ClaimStatusProvedHere]
\label{thm:fermion-genus1-partition}
For the free fermion on the elliptic curve
$E_\tau = \mathbb{C}/(\mathbb{Z} + \tau\mathbb{Z})$,
the partition function in spin structure $\sigma$ is:
\begin{equation}\label{eq:fermion-Z1-spin}
Z_1(\cF, \sigma)
= \left(\frac{\vartheta_\sigma(0|\tau)}{\eta(\tau)}\right)^{1/2},
\end{equation}
where $\vartheta_\sigma$ is the Jacobi theta function with
characteristic $\sigma$. The unsquared ratio
$\vartheta_\sigma/\eta$ is the complex fermion, equivalently two
real fermions. For the single real fermion of this chapter:
\begin{align}
\label{eq:fermion-Z1-theta2}
Z_1(\cF, \sigma_2) &=
\left(\frac{\vartheta_2(0|\tau)}{\eta(\tau)}\right)^{1/2}, \\[4pt]
\label{eq:fermion-Z1-theta3}
Z_1(\cF, \sigma_3) &=
\left(\frac{\vartheta_3(0|\tau)}{\eta(\tau)}\right)^{1/2}, \\[4pt]
\label{eq:fermion-Z1-theta4}
Z_1(\cF, \sigma_4) &=
\left(\frac{\vartheta_4(0|\tau)}{\eta(\tau)}\right)^{1/2}, \\[4pt]
\label{eq:fermion-Z1-odd}
Z_1(\cF, \sigma_1) &=
\left(\frac{\vartheta_1(0|\tau)}{\eta(\tau)}\right)^{1/2} = 0
\quad \text{\textup{(}odd spin structure\textup{)}}.
\end{align}
The ratio $\vartheta_\sigma(0|\tau)/\eta(\tau)$ is the determinant
of the complex fermion. Its square root is the Pfaffian line of the
single real fermion, with spin structure~$\sigma$.
\end{theorem}

\begin{proof}
The fermion partition function on $E_\tau$ with boundary
conditions $(\alpha, \beta)$ is
$\mathrm{Tr}_{H_{\alpha}}\bigl((-1)^{\beta F}\, q^{L_0 - c/24}\bigr)$,
where $H_\alpha$ is the Fock space with periodicity $\alpha$
and $(-1)^{\beta F}$ inserts a sign for the $\beta$-cycle.
For a single real fermion of weight $h = 1/2$ with
$c = 1/2$:

\emph{Step~1.} In the NS sector ($\alpha = 0$), the mode
expansion $\psi(z) = \sum_{r \in \mathbb{Z} + 1/2}
\psi_r z^{-r-1/2}$ has half-integer modes. The Fock space
trace gives:
\[
\mathrm{Tr}_{\mathrm{NS}}\bigl(q^{L_0 - 1/48}\bigr)
= q^{-1/48} \prod_{n=1}^\infty (1 + q^{n-1/2})
= \sqrt{\frac{\vartheta_3(0|\tau)}{\eta(\tau)}}.
\]
Taking the single real fermion trace gives
\eqref{eq:fermion-Z1-theta3} and
\eqref{eq:fermion-Z1-theta4} (the latter with
$(-1)^F$ insertion). Squaring these formulas gives the complex
fermion ratio $\vartheta_\sigma/\eta$.

\emph{Step~2.} In the Ramond sector ($\alpha = 1$), the mode
expansion has integer modes, and the zero mode $\psi_0$
satisfies $\psi_0^2 = 1/2$. The Ramond Pfaffian is the
chosen square root of the corresponding complex determinant:
\[
\mathrm{Tr}_{\mathrm{R}}\bigl(q^{L_0 - 1/48}\bigr)
= \sqrt{\frac{\vartheta_2(0|\tau)}{\eta(\tau)}},
\]
and the odd spin structure
$\mathrm{Tr}_{\mathrm{R}}\bigl((-1)^F q^{L_0 - 1/48}\bigr) = 0$
because the odd theta characteristic has a zero mode, so the
Pfaffian vanishes.
\end{proof}

\begin{computation}[GSO-summed partition function]
\label{comp:fermion-gso}
\index{GSO projection!free fermion}
\index{free fermion!GSO projection}
The GSO (Gliozzi--Scherk--Olive) projection sums over spin
structures with equal weight, discarding the odd spin
structure (which vanishes):
\begin{equation}\label{eq:fermion-gso}
Z_1^{\mathrm{GSO}}(\cF)
= \frac{1}{2}\left(
 \frac{\vartheta_2(0|\tau)^2}{\eta(\tau)^2}
+ \frac{\vartheta_3(0|\tau)^2}{\eta(\tau)^2}
+ \frac{\vartheta_4(0|\tau)^2}{\eta(\tau)^2}
\right).
\end{equation}
This is the squared, even-spin GSO block appropriate to a complex
fermion or to a paired real-fermion sector. For a single real fermion
one must choose Pfaffian square roots before summing; the scalar
obstruction coefficient remains $\kappa(\cF)=1/4$.
The Jacobi identity
$\vartheta_3^4 = \vartheta_2^4 + \vartheta_4^4$
controls the even-spin vector, but it does not collapse the sum of
theta squares to a single theta square. The expression above is
therefore best read as the even-spin block for a paired
real-fermion sector; full modular invariance requires the standard
spin-structure vector and its multiplier system, not a scalar
identity among the three squares. This is in contrast to the Heisenberg
partition function $Z_1(\cH_\kappa) = 1/\eta(\tau)$,
which transforms under the full $\mathrm{SL}_2(\mathbb{Z})$
with a multiplier system
\textup{(}Computation~\textup{\ref{comp:heisenberg-partition-g1})}.
\end{computation}

\begin{remark}[Heisenberg comparison: $\eta^{-1}$ vs
$\vartheta/\eta$]
\label{rem:fermion-heisenberg-genus1-comparison}
\index{free fermion!vs Heisenberg at genus 1}
The structural difference between the two Gaussian archetypes
at genus~$1$ is:
\[
\underbrace{Z_1(\cH_\kappa)
= \eta(\tau)^{-1}}_{\text{determinant}}
\qquad\text{vs}\qquad
\underbrace{Z_1(\cF, \sigma)
= (\vartheta_\sigma / \eta)^{1/2}}_{\text{Pfaffian}}.
\]
The Heisenberg partition function involves the determinant of
the scalar Laplacian (an even functional of the metric), so
no spin structure is needed. The fermion partition function
involves the Pfaffian of the Dirac operator (an odd
functional), which requires and depends on a spin structure.
Both have $\eta(\tau)$ in the denominator (the universal
conformal anomaly), but the numerator is $1$ for the boson
and $\vartheta_\sigma$ for the fermion.
\end{remark}

\subsection{Free energy extraction and the shadow obstruction tower
prediction}
\label{subsec:fermion-free-energy}
\index{free fermion!free energy}
\index{shadow tower!free fermion verification}

\begin{theorem}[Free fermion genus-1 free energy;
\ClaimStatusConditional]
\label{thm:fermion-F1-shadow}
The genus-$1$ free energy of the free fermion is:
\begin{equation}\label{eq:fermion-F1-shadow-verification}
F_1(\cF) = \frac{\kappa(\cF)}{24}
= \frac{1/4}{24} = \frac{1}{96}.
\end{equation}
This matches the shadow obstruction tower prediction
$F_1(\cF) = \kappa(\cF) \cdot \lambda_1^{\mathrm{FP}}
= \kappa(\cF)/24$
\textup{(}Theorem~\textup{\ref{thm:genus-universality})}.
\end{theorem}

\begin{proof}
The free energy $F_1$ is extracted from the partition function
by $F_1 = -\log Z_1|_{q^0}$ (the constant term in the
$q$-expansion of $\log Z_1$). For each even spin structure
$\sigma$, the ratio $\vartheta_\sigma/\eta$ has
$q$-expansion beginning with $q^{h_\sigma - 1/48 + 1/24}$
where $h_\sigma$ is the ground-state energy. The regularized
free energy is extracted from the effective central charge:
\[
c_{\mathrm{eff}} = c - 24 h_{\min} = \frac{1}{2},
\]
giving
$F_1(\cF)=c_{\mathrm{eff}}/48=1/96=\kappa(\cF)/24$.
Alternatively, the result follows from the genus universality
theorem applied to $\kappa(\cF) = c/2 = 1/4$:
\[
F_1(\cF) = \kappa(\cF) \cdot \frac{1}{24}
= \frac{1}{4} \cdot \frac{1}{24} = \frac{1}{96}. \qedhere
\]
\end{proof}

\begin{computation}[Free fermion genus expansion]
\textup{(}UNIFORM-WEIGHT\textup{})
\label{comp:fermion-genus-expansion}
\index{free fermion!genus expansion table}
By Theorem~\textup{\ref{thm:genus-universality}} with
$\kappa(\cF) = 1/4$, the scalar/full free energy at genus $g \geq 1$
is:
\begin{equation}\label{eq:fermion-Fg-formula}
F_g(\cF)=F_g^{\mathrm{sc}}(\cF)
= \frac{1}{4} \cdot
\frac{2^{2g-1}-1}{2^{2g-1}} \cdot \frac{|B_{2g}|}{(2g)!}.
\end{equation}
Explicit values:
\begin{center}
\renewcommand{\arraystretch}{1.3}
\begin{tabular}{ccc}
\toprule
$g$ & $F_g(\cF)$ \text{(exact)} & $F_g(\cF)$ \text{(decimal)} \\
\midrule
$1$ & $1/96$ & $1.04 \times 10^{-2}$ \\[2pt]
$2$ & $7/23040$ & $3.04 \times 10^{-4}$ \\[2pt]
$3$ & $31/3870720$ & $8.01 \times 10^{-6}$ \\[2pt]
$4$ & $127/619315200$ & $2.05 \times 10^{-7}$ \\[2pt]
$5$ & $511/98099527680$ & $5.21 \times 10^{-9}$ \\
\bottomrule
\end{tabular}
\end{center}
The ratio to the Heisenberg values
\textup{(}Table~\textup{\ref{tab:free-energy-values})} is
$F_g(\cF)/F_g(\cH_1)
=F_g^{\mathrm{sc}}(\cF)/F_g^{\mathrm{sc}}(\cH_1)
=1/4$ at every genus, confirming
$\kappa(\cF)/\kappa(\cH_1) = (1/4)/1 = 1/4$.
The generating function is:
\begin{equation}\label{eq:fermion-GF}
\sum_{g=1}^\infty F_g(\cF)\,\hbar^{2g}
=\sum_{g=1}^\infty F_g^{\mathrm{sc}}(\cF)\,\hbar^{2g}
= \frac{1}{4\hbar^2}\bigl(\hat{A}(i\hbar) - 1\bigr)
= \frac{1}{4}\Bigl(\frac{1}{24}
+ \frac{7}{5760}\hbar^2
+ \frac{31}{967680}\hbar^4 + \cdots\Bigr).
\end{equation}
\end{computation}

\subsection{Higher genus: Pfaffians and theta constants}
\label{subsec:fermion-higher-genus}
\index{free fermion!higher genus}
\index{Pfaffian!genus-$g$ partition function}

\begin{theorem}[Free fermion at genus $g$;
\ClaimStatusConditional]
\label{thm:fermion-genus-g}
On a genus-$g$ surface $\Sigma_g$ with spin structure $\sigma$,
the free fermion partition function is:
\begin{equation}\label{eq:fermion-Zg}
Z_g(\cF, \sigma)
= \bigl(\det \bar\partial_\sigma\bigr)^{1/2}
= \mathrm{Pf}(\bar\partial_\sigma),
\end{equation}
the Pfaffian of the Dirac operator twisted by $\sigma$.
In terms of the period matrix
$\Omega \in \mathfrak{H}_g$ and the theta function with
characteristic $\delta = [\alpha, \beta]$
\textup{(}$\alpha, \beta \in (\mathbb{Z}/2)^g$\textup{)}:
\begin{equation}\label{eq:fermion-theta-char}
Z_g(\cF, \delta)
= \frac{\vartheta[\delta](0|\Omega)^{1/2}}
 {\bigl(\det \mathrm{Im}\,\Omega\bigr)^{1/4}},
\end{equation}
where the genus-$g$ theta function with characteristic is:
\begin{equation}\label{eq:genus-g-theta-char}
\vartheta[\delta](z|\Omega)
= \sum_{n \in \mathbb{Z}^g}
 \exp\bigl(\pi i(n+\alpha)^T \Omega\, (n+\alpha)
 + 2\pi i(n+\alpha)^T(z+\beta)\bigr).
\end{equation}
The obstruction class is
$\mathrm{obs}_g(\cF) = \tfrac{1}{4}\,\lambda_g$
\textup{(}UNIFORM-WEIGHT\textup{}), independent
of the spin structure $\sigma$, because $\kappa(\cF) = 1/4$ is
a scalar invariant of the algebra that does not depend on global
choices.
\end{theorem}

\begin{proof}
\emph{Step~1: Pfaffian structure.}
The free fermion path integral on $\Sigma_g$ with spin
structure $\sigma$ is
$\int \mathcal{D}\psi\, e^{-S[\psi]}$
where $S[\psi] = \frac{1}{2}\int \psi\, \bar\partial_\sigma\,\psi$.
Since $\bar\partial_\sigma$ is a skew-symmetric operator on
sections of $L_\sigma$ (the square root of $\omega_{\Sigma_g}$
determined by $\sigma$), the Gaussian integral evaluates to
$\mathrm{Pf}(\bar\partial_\sigma) = (\det \bar\partial_\sigma)^{1/2}$.
This is the fundamental distinction from the Heisenberg case,
where $\bar\partial$ acts on the trivial bundle and the integral
gives $(\det \bar\partial)^{-1}$: the boson contributes
an \emph{inverse} determinant, the fermion a
\emph{half-determinant}.

\emph{Step~2: Theta function representation.}
The bosonization formula (Proposition~\ref{prop:bosonization}
below) identifies
$\mathrm{Pf}(\bar\partial_\sigma)$ with the genus-$g$
theta constant $\vartheta[\delta](0|\Omega)$ up to a universal
normalization factor. The denominator
$(\det \mathrm{Im}\,\Omega)^{1/4}$ absorbs the conformal
anomaly ($c = 1/2$ contributes half of the $c = 1$
normalization).

\emph{Step~3: Obstruction class.}
The genus-$g$ obstruction class in the modular deformation is computed
from the $B$-cycle monodromies of the spin-$1/2$ propagator. On
$\Sigma_g$,
the genus-$g$ fermion propagator is the Szeg\H{o} kernel
$S_\sigma(z,w) = \vartheta[\delta](z{-}w|\Omega) /
(E(z,w)\, \vartheta[\delta](0|\Omega))$,
where $E(z,w)$ is the prime form. The propagator depends on
$\sigma$, but the obstruction differential extracts only the
universal part: the $B$-cycle monodromy defect of
$d\log E(z,w)$ (the bar propagator, which has weight $1$ and
is spin-structure-independent; see
Remark~\ref{rem:fermion-propagator-universality}).
Therefore $\mathrm{obs}_g(\cF) = \kappa(\cF) \cdot \lambda_g
= \frac{1}{4}\lambda_g$, exactly as for the Heisenberg
(Theorem~\ref{thm:heisenberg-genus2-obstruction}).
\end{proof}

\begin{remark}[Spin structure independence of the
obstruction class]
\label{rem:fermion-propagator-universality}
\index{free fermion!spin structure independence of
$\kappa$}
The spin structure $\sigma$ enters the genus-$g$ partition
function $Z_g(\cF, \sigma)$ through the Szeg\H{o} kernel, but
the modular characteristic $\kappa(\cF) = 1/4$ and the
obstruction class
\(\mathrm{obs}_g(\cF)=\kappa(\cF)\lambda_g=\frac14\lambda_g\)
\textup{(}UNIFORM-WEIGHT\textup{}) are
spin-structure-independent. The reason: $\kappa(\cF)$ is determined
by the \emph{local} OPE data (the quartic-pole coefficient of
$T(z)T(w)$, which is $c/2$), while spin structures are
\emph{global} data on $\Sigma_g$. In the language of the
shadow obstruction tower, $\kappa(\cF)$ is a genus-$0$ invariant
(Definition~\ref{def:modular-characteristic-package}) that determines
the genus-$g$ obstruction class via the universal formula
\(\mathrm{obs}_g(\cF)=\kappa(\cF)\lambda_g\)
\textup{(}UNIFORM-WEIGHT\textup{)}.
\end{remark}

\begin{remark}[Genus-$2$ theta constants]
\label{rem:fermion-genus2}
\index{free fermion!genus-2 partition function}
At genus~$2$, the $10$ even spin structures give $10$
genus-$2$ theta constants
$\vartheta[\delta](0|\Omega)$ for the $10$ even
characteristics $\delta$. The $6$ odd characteristics
give $\vartheta[\delta](0|\Omega) = 0$ (the odd theta
constants vanish at $z = 0$). The paired real-fermion GSO block is
\[
Z_2^{\mathrm{GSO}}(\cF)
= \frac{1}{2^4}\sum_{\delta\,\mathrm{even}}
 \vartheta[\delta](0|\Omega)^2,
\]
where the sum runs over the $10$ even characteristics.
For a single real fermion the Pfaffian square roots must be chosen
before summing over spin structures.
This expression is a Siegel modular form of weight $1/2$
for a subgroup of $\mathrm{Sp}_4(\mathbb{Z})$, in
contrast to the Heisenberg genus-$2$ partition function
$Z_2(\cH_\kappa) = (\det \mathrm{Im}\,\Omega)^{-1/2}$
\textup{(}Computation~\textup{\ref{comp:partition-genus-two})},
which transforms under the full $\mathrm{Sp}_4(\mathbb{Z})$.
\end{remark}

\subsection{The Ising model identification}
\label{subsec:ising-identification}
\index{Ising model!free fermion identification|textbf}
\index{minimal model!$M(3,4)$}

\begin{proposition}[Ising $=$ free fermion;
\ClaimStatusProvedElsewhere]
\label{prop:ising-fermion}
The free fermion $\cF$ at $c = 1/2$ is isomorphic to the
minimal model $M(3,4)$, the critical Ising model. The primary
fields and their conformal weights are:
\begin{center}
\renewcommand{\arraystretch}{1.3}
\begin{tabular}{lcl}
\toprule
\textbf{Primary field} & $h$ & \textbf{Physical interpretation} \\
\midrule
$\mathbf{1}$ & $0$ & Identity (vacuum) \\
$\sigma$ & $1/16$ & Spin field (order parameter) \\
$\psi$ & $1/2$ & Free fermion (energy operator) \\
\bottomrule
\end{tabular}
\end{center}
The fusion rules are:
$\psi \times \psi = \mathbf{1}$,\quad
$\psi \times \sigma = \sigma$,\quad
$\sigma \times \sigma = \mathbf{1} + \psi$.
\end{proposition}

\begin{proof}[References]
The identification is classical: the $c = 1/2$
minimal model has Kac table with $p = 3$, $p' = 4$,
central charge
$c = 1 - 6(p-p')^2/(pp') = 1 - 6/12 = 1/2$,
and three primary fields with weights
$h_{r,s} = ((4r - 3s)^2 - 1)/48$
for $(r,s) \in \{(1,1), (1,2), (1,3)\}$, giving
$h = 0, 1/16, 1/2$. The fermion $\psi$ with
$h = 1/2$ is the energy operator; the spin field $\sigma$
with $h = 1/16$ is the disorder operator that intertwines
the NS and Ramond sectors. See
\cite{DiFrancesco1997}, Chapter~7.
\end{proof}

\begin{remark}[Koszul duality and the Ising model]
\label{rem:ising-koszul}
\index{Ising model!Koszul dual}
The Koszul dual $\cF^!
= \mathrm{Sym}^{\mathrm{ch}}(\gamma)$ is a single free boson
of weight~$1/2$
\textup{(}Theorem~\textup{\ref{thm:single-fermion-boson-duality})}.
This is \emph{not} the Ising model at a different temperature:
Koszul duality exchanges statistics
(fermionic $\leftrightarrow$ bosonic) and changes the structure
of the bar complex, but the Ising model interpretation applies
only to the original $\cF$. The Koszul dual has
$\kappa(\cF^!) = -1/4$ and $c(\cF^!) = -1/2$
(non-unitary), placing it outside the minimal model
classification.
\end{remark}

\subsection{Boson-fermion correspondence}
\label{subsec:boson-fermion-correspondence}
\index{bosonization!free fermion|textbf}
\index{boson-fermion correspondence}

\begin{proposition}[Bosonization formula;
\ClaimStatusProvedElsewhere]
\label{prop:bosonization}
Let $\phi(z)$ be a free boson compactified on a circle of
radius $R = 1$ \textup{(}i.e.\ the Heisenberg algebra
$\cH_1$ at level $\kappa = 1$ with the identification
$\phi \sim \phi + 2\pi$\textup{)}. Then the vertex operators:
\begin{equation}\label{eq:bosonization}
\psi(z) = e^{i\phi(z)}, \qquad
\bar\psi(z) = e^{-i\phi(z)}
\end{equation}
satisfy the OPE of a \textup{complex} free fermion:
$\psi(z)\bar\psi(w) \sim 1/(z{-}w)$. The composite
$\normord{\psi\bar\psi} = i\partial\phi = iJ$ reproduces
the Heisenberg current.
\end{proposition}

\begin{proof}[References]
Standard; see \cite{FMS86}, Chapter~14, or
\cite{Pol98}, \S10.3.
\end{proof}

\begin{remark}[Bosonization is not Koszul duality]
\label{rem:bosonization-not-koszul}
\index{bosonization!vs Koszul duality}
The boson-fermion correspondence
\eqref{eq:bosonization} relates a \emph{complex}
fermion ($c = 1$) to a compactified boson ($c = 1$):
both sides have the same central charge. Koszul
duality, by contrast, maps the \emph{real} fermion
$\cF$ ($c = 1/2$) to a bosonic algebra $\cF^!$
($c = -1/2$): the central charges have opposite signs.
The two operations are structurally distinct:
\begin{center}
\renewcommand{\arraystretch}{1.3}
\begin{tabular}{lll}
\toprule
& \textbf{Bosonization} & \textbf{Koszul duality} \\
\midrule
Input & complex fermion ($c = 1$) & real fermion ($c = 1/2$) \\
Output & compactified boson ($c = 1$) & free boson ($c = -1/2$) \\
Mechanism & vertex operator & bar construction + Verdier \\
$c$-relation & $c_{\mathrm{in}} = c_{\mathrm{out}}$
& $c_{\mathrm{in}} + c_{\mathrm{out}} = 0$ \\
$\kappa$-relation & N/A & $\kappa + \kappa' = 0$ \textup{(}free-field pair\textup{)} \\
\bottomrule
\end{tabular}
\end{center}
Bosonization is an isomorphism of conformal field theories;
Koszul duality applies the bar construction to $\cA$ and then
Verdier-dualizes the resulting coalgebra, producing the homotopy
Koszul-dual factorization algebra
$\mathbb{D}_{\mathrm{Ran}}B(\cA) \simeq \cA^!_\infty$
\textup{(}Theorem~\textup{\ref{thm:bar-cobar-isomorphism-main})}.
\end{remark}

\begin{remark}[The real free fermion and the Heisenberg
current]
\label{rem:real-fermion-heisenberg}
\index{free fermion!relation to Heisenberg current}
The single real fermion $\cF$ does not contain a Heisenberg current:
$\normord{\psi\psi}=0$ by fermionic antisymmetry. A Heisenberg
current appears only after taking two real fermions, equivalently one
complex fermion with fields $\psi,\bar\psi$, where
\[
J(z)=\normord{\psi(z)\bar\psi(z)}
\]
has OPE $J(z)J(w)\sim (z-w)^{-2}$. Thus
$\cH_1$ embeds in the complex-fermion algebra
$\cF\otimes\cF$, not in a single copy of~$\cF$. This is another
manifestation of the difference between $c=1/2$ and $c=1$.
\end{remark}

\subsection{Comparison: free fermion vs Heisenberg}
\label{subsec:fermion-heisenberg-comparison}
\index{free fermion!vs Heisenberg}

\begin{table}[ht]
\centering
\caption{Gaussian archetypes: free fermion vs Heisenberg}
\label{tab:fermion-heisenberg-comparison}
\begin{tabular}{lll}
\toprule
& \textbf{Free fermion} $\cF$ & \textbf{Heisenberg} $\cH_\kappa$ \\
\midrule
Central charge & $c = 1/2$ (fixed) & $c = 1$ (fixed, $\kappa$ free) \\
$\kappa$ & $1/4$ (fixed) & $\kappa$ (tunable) \\
Shadow class & G, $r_{\max} = 2$ & G, $r_{\max} = 2$ \\
OPE pole & simple: $1/(z{-}w)$ & double: $\kappa/(z{-}w)^2$ \\
$r$-matrix & $r(z) = 0$ & $r(z) = \kappa/z$ \\
Statistics & fermionic & bosonic \\
Spin structure & required & not required \\
$Z_1$ & $(\vartheta_\sigma/\eta)^{1/2}$ & $\eta^{-1}$ \\
Partition function type & Pfaffian & determinant \\
$F_1$ & $1/96$ & $\kappa/24$ \\
$F_g/F_g(\cH_1)$ & $1/4$ (all $g$) & $\kappa$ (all $g$) \\
Koszul dual & $\mathrm{Sym}^{\mathrm{ch}}(\gamma)$ & $\operatorname{Sym}^{\mathrm{ch}}(V^*)$ \\
$\kappa + \kappa'$ & $0$ & $0$ \\
Sewing mechanism & Szeg\H{o} kernel & Bergman kernel \\
Bosonization & $\psi = e^{i\phi}$ (complex) & N/A \\
Ising model & $M(3,4)$ at $c = 1/2$ & N/A \\
\bottomrule
\end{tabular}
\end{table}

\begin{remark}[Universality of the $\hat{A}$-genus]
\label{rem:fermion-ahat-universality}
\index{free fermion!$\hat{A}$-genus}
Both Gaussian archetypes (the free fermion and the Heisenberg)
produce the same $\hat{A}$-genus generating function,
differing only in the overall scalar modular characteristic.
By the genus universality theorem
(Theorem~\ref{thm:genus-universality}), class~G
is characterized by $r_{\max} = 2$: all higher
shadow invariants vanish and the genus expansion is
controlled by the scalar modular characteristic alone.
The free fermion confirms this at the non-trivial
value $\kappa(\cF)=1/4$ (non-integer, unlike the Heisenberg
where $\kappa$ is typically an integer). The
$\hat{A}$-genus structure is insensitive to statistics
(bosonic vs fermionic), spin structure dependence, and the
distinction between determinants and Pfaffians: these are
features of the partition function $Z_g$, not of the
free energy
$F_g(\cF)=F_g^{\mathrm{sc}}(\cF)
=\kappa(\cF) \cdot \lambda_g^{\mathrm{FP}}$
\textup{(}UNIFORM-WEIGHT, class~G exact lane\textup{)}.
\end{remark}


\section{Virasoro, strings, and moduli}\label{sec:virasoro-strings-moduli}

The Virasoro bar complex is analyzed at general central charge. On the
critical locus $c = 26$, the genus-$0$ equivariant anomaly vanishes,
giving the proved critical BRST/bar comparison on the total
matter+ghost string algebra. Stronger claims about full moduli-space
cohomology, physical amplitudes, or modular invariance require
additional input and are fenced below.

\subsection{Virasoro at critical central charge}

The free-field dualities above have finite shadow depth. The
Heisenberg scalar lane is governed by the level:
$\kappa(\cH_k)=k$ and
$r_{\cH_k}^{\mathrm{coll}}(z)=k/z$, although
$c(\cH_k)=1$. The Virasoro scalar lane is different. Its generator is
the stress tensor itself, so
$\kappa(\mathrm{Vir}_c)=c/2$ and
\[
r_{\mathrm{Vir}_c}^{\mathrm{coll}}(z)
=\frac{c/2}{z^3}+\frac{2T}{z}.
\]
The autonomous stress-tensor OPE has an infinite shadow tower on the
non-singular surface $c(5c+22)\neq0$; this is class~$\mathsf M$, not a
Gaussian free-field lane.

\subsubsection{Setup}

\begin{definition}[Virasoro algebra]
The Virasoro algebra $\text{Vir}_c$ has stress-energy tensor $T(z)$ of weight 2 with OPE:
\[
T(z)T(w) = \frac{c/2}{(z-w)^4} + \frac{2T(w)}{(z-w)^2} + \frac{\partial T(w)}{z-w} + \text{regular}
\]
At $c = 26$ (critical dimension), special cancellations occur.
\end{definition}
\index{Virasoro algebra|textbf}

\subsubsection{Bar complex and critical descent}\label{sec:virasoro-bar-moduli-main}

The fourth-order pole in $T(z)T(w)$ prevents naive Koszulness, but at
the critical central charge $c = 26$ the conformal anomaly cancels the
ghost contribution and the genus-$0$ equivariant bar complex becomes an
honest cochain complex on the quotient configuration space. This is
the proved algebraic input behind the BRST/bar dictionary. A
stronger identification with the cohomology of
$\overline{\mathcal{M}}_{0,n+3}$ is a separate moduli-space theorem.

\begin{theorem}[Critical Virasoro descent at $c = 26$; \ClaimStatusProvedHere]\label{thm:virasoro-moduli}
For $\text{Vir}_c$ on $\mathbb{P}^1$, the genus-$0$ equivariant bar
complex satisfies
\[
d_{\mathrm{equiv}}^2 = (c/2 - 13)\,\mu,
\qquad
\mu = \kappa_1 \in H^2(\overline{\mathcal{M}}_{0,n+3}).
\]
Hence at $c = 26$ the equivariant bar complex descends to a genuine
cochain complex on
$\overline{C}_{n+3}(\mathbb{P}^1)/\mathrm{PSL}_2(\mathbb{C})
\cong \overline{\mathcal{M}}_{0,n+3}$.
\end{theorem}

\begin{proof}
At $c = 26$, the conformal anomaly of the
$bc$ ghost system cancels that of the matter sector, so the BRST
differential squares to zero without correction terms. In bar-complex
language: the relevant equivariant obstruction vanishes, and the
quotient bar complex is an honest cochain complex.

\emph{Step~1.}
The Virasoro algebra contains $\mathfrak{sl}_2 = \langle L_{-1}, L_0, L_1 \rangle$ acting on $\mathbb{P}^1$ by Möbius transformations. The $\mathrm{PSL}_2(\mathbb{C})$ action on $C_{n+3}(\mathbb{P}^1)$ is free, and the quotient is $C_{n+3}(\mathbb{P}^1)/\mathrm{PSL}_2 \cong C_n(\mathbb{C}^*)$, the configuration of $n$ points on the thrice-punctured sphere. The Fulton--MacPherson compactification is equivariant, giving $\overline{C}_{n+3}(\mathbb{P}^1)/\mathrm{PSL}_2 \cong \overline{\mathcal{M}}_{0,n+3}$ (Keel's theorem \cite{Kee92}).

\emph{Step~2.}
The bar complex $\bar{B}^n_{\mathrm{geom}}(\mathrm{Vir}_c)$ consists of logarithmic forms on $\overline{C}_{n+1}(\mathbb{P}^1)$ with coefficients in $\mathrm{Vir}_c^{\boxtimes(n+1)}$. The Virasoro Ward identities (annihilation by $L_{-1}, L_0, L_1$) precisely implement the $\mathrm{PSL}_2$-equivariance condition: a form $\omega \in \bar{B}^n$ satisfying $\iota_{L_k}\omega = 0$ for $k = -1, 0, 1$ descends to $\overline{\mathcal{M}}_{0,n+3}$.

\emph{Step~3.}
For general $c$, the descent involves an anomalous term: the Lie
derivative $\mathcal{L}_{L_k}$ acting on bar complex elements picks up
a central extension contribution proportional to $c$. The equivariant
bar complex differential satisfies
\[
d_{\mathrm{equiv}}^2=(c/2-13)\,\mu,
\]
where $\mu$ is the Mumford class
$\kappa_1 \in H^2(\overline{\mathcal{M}}_{0,n+3})$. The coefficient is
$\kappa(\mathrm{Vir}_c)-\kappa(\mathrm{Vir}_{26})=c/2-13$, measuring
the departure from the critical Koszul dual. At $c=26$, this anomaly
vanishes: $d_{\mathrm{equiv}}^2=0$, so the equivariant bar complex is a
genuine cochain complex.
\end{proof}

\begin{remark}[Ghost system as Koszul dual]
\label{rem:ghost-koszul-identification}
\index{ghost system!Koszul dual identification}
\index{BRST operator!bar differential}
The $bc$ ghost system of central charge $c_{\mathrm{ghost}} = -26$
that Polyakov introduces for gauge-fixing the worldsheet
diffeomorphism group supplies the Koszul-dual ghost sector in the
critical total-string comparison of
Theorem~\ref{thm:brst-bar-genus0}. On that proved genus-$0$ locus, the
relative BRST complex of the total algebra
$\cA_{\mathrm{tot}} = \mathrm{Vir}_{26} \otimes bc$ is
quasi-isomorphic to the genus-$0$ bar complex of the same total
algebra. Thus physical BRST classes are compared to bar-cohomology
classes through that chain map. The manuscript does \emph{not} prove a
direct matter-only identity
$Q_{\mathrm{BRST}} = d_{\barB}$ or
$H^*_{\mathrm{BRST}}(\mathrm{String}) = H^*(\barB(\mathrm{Vir}_{26}))$.
\par
The ghost number grading is the cohomological degree on the BRST side;
the bar complex provides the algebraic model on the proved comparison
locus.
\end{remark}

\begin{remark}[Polyakov action as degree-$2$ shadow]
\label{rem:polyakov-action-degree-two}
\index{Polyakov action!shadow tower interpretation}
\index{shadow tower!Polyakov action}
The Polyakov action
$S_{\mathrm{Pol}} = \frac{1}{4\pi\alpha'}\int \sqrt{g}\,
g^{ab}\,\partial_a X^\mu\,\partial_b X_\mu$
is quadratic in the fields; it controls only the degree-$2$ sector
of the worldsheet theory. The degree-$2$ shadow
$\Theta_\cA^{\leq 2} = \kappa \cdot \eta \otimes \Lambda \in
\mathrm{Def}_{\mathrm{cyc}}^{\mathrm{mod}}(\cA)$ is the algebraic
counterpart: its Chern--Weil evaluation against a conformal
variation reproduces the Polyakov formula with coefficient
$\kappa/(6\pi)$ (Proposition~\ref{prop:polyakov-chern-weil}).
The higher-degree shadows (cubic~$\mathfrak{C}$ at $r = 3$,
quartic~$\mathfrak{Q}^{\mathrm{contact}}$ at $r = 4$, and the
full tower $\Theta_\cA^{\leq r}$ for $r \geq 5$) encode the
nonlinear OPE corrections that the quadratic Polyakov action
cannot detect.
\end{remark}

\begin{remark}[Critical descent and the moduli comparison]
\label{rem:virasoro-bar-moduli-polyakov}
\index{Virasoro algebra!moduli correspondence!Polyakov significance}
Theorem~\ref{thm:virasoro-moduli} proves only the critical descent
statement: the genus-$0$ equivariant anomaly vanishes at $c = 26$, so
the quotient bar complex is an honest cochain complex on the
$\overline{\mathcal{M}}_{0,n+3}$ stage. The stronger claim that this
complex computes the full cohomology
$H^*(\overline{\mathcal{M}}_{0,n+3})$, or that the bosonic string path
integral literally localises there, requires the separate
moduli-integration comparison
\textup{(}Corollary~\ref{cor:string-amplitude-genus0}\textup{)}, and
the higher-genus amplitude identification is
\textup{(}Conjecture~\ref{conj:string-amplitude-higher-genus}\textup{)}.
Away from $c = 26$, the curvature
$d_{\mathrm{equiv}}^2 = (c/2 - 13)\,\kappa_1$ obstructs the critical
descent picture.
\end{remark}

\subsubsection{A residue model on moduli space}

\begin{proposition}[Boundary-residue differential on moduli forms; \ClaimStatusProvedHere]
\label{prop:moduli-degeneration}
On logarithmic forms on $\overline{\mathcal{M}}_{0,n+3}$, the natural
boundary-residue operator
\[
d_{\partial}\colon \Omega^n_{\log}(\overline{\mathcal{M}}_{0,n+3})
\longrightarrow
\bigoplus_I \Omega^{n-1}_{\log}(D_I)
\]
is given by
\[
d_{\partial}\omega = \sum_{\text{nodes}} \operatorname{Res}_{\text{node}} \omega,
\]
where the sum is over all possible nodal degenerations.
\end{proposition}

\begin{proof}
The boundary $\partial\overline{\mathcal{M}}_{0,n+3}$ is a union of divisors $D_I$, one for each partition $[n+3] = I \sqcup J$ with $|I|, |J| \geq 2$, parametrizing curves where the $I$-labeled points bubble off onto a separate component. Each $D_I \cong \overline{\mathcal{M}}_{0,|I|+1} \times \overline{\mathcal{M}}_{0,|J|+1}$, and the collision of points $z_i \to z_j$ in $\overline{C}_{n+3}(X)$ maps to the boundary stratum $D_{\{i,j\}}$ in moduli space. A logarithmic form $\omega \in \Omega^n_{\log}(\overline{\mathcal{M}}_{0,n+3})$ has at most a simple pole along $D_I$; in a local coordinate $t$ with $D_I = \{t = 0\}$, write $\omega = f \cdot d\log t + g$ with $g$ regular. Then $\mathrm{Res}_{D_I}\omega = f|_{t=0} \in \Omega^{n-1}(D_I)$.
\[
d_{\partial}\omega = \sum_I \mathrm{Res}_{D_I}\omega.
\]
This is the standard boundary-residue operator on the logarithmic
de~Rham model of $\overline{\mathcal{M}}_{0,n+3}$. It matches the
nodal-degeneration pattern of the critical descent construction. The
chain-level identification between $d_{\partial}$ and the Virasoro bar
differential is a separate comparison map, not part of this
proposition.
\end{proof}

\subsubsection{Explicit low-degree computation}

\begin{example}[Low degrees for the moduli-space model]
\begin{itemize}
\item Degree 0: the geometric model has $H^0(\overline{\mathcal{M}}_{0,3}) = \mathbb{C}$.
\item Degree 1: $H^1(\overline{\mathcal{M}}_{0,4}) = 0$ since $\overline{\mathcal{M}}_{0,4} \cong \mathbb{P}^1$.
\item Degree 2: $H^2(\overline{\mathcal{M}}_{0,5}) \cong \mathbb{C}^5$ since $\overline{\mathcal{M}}_{0,5} \cong \mathrm{Bl}_4\mathbb{P}^2$ is the del~Pezzo surface of degree~$5$.
\item Higher degrees require the full geometry of the boundary-residue complex on $\overline{\mathcal{M}}_{0,n+3}$.
\end{itemize}
These values describe the moduli-space model only. They do not prove
that
$H^*(\bar{B}_{\mathrm{geom}}(\mathrm{Vir}_{26}))$ is identified with
the cohomology of $\overline{\mathcal{M}}_{0,n+3}$.
\end{example}


\subsection{String vertex algebra}

\subsubsection{Setup}

\begin{definition}[String vertex algebra]
The bosonic string vertex algebra combines:
\begin{itemize}
\item Matter: 26 free bosons $X^\mu$ with $T_{\text{matter}} = -\frac{1}{2}\partial X^\mu \partial X_\mu$
\item Ghosts: $(b,c)$ with weights $(2,-1)$ and $T_{\text{ghost}} = -2b\partial c - (\partial b)c$
\item BRST charge: $Q = \oint \left(c T_{\text{matter}} + bc\partial c + \frac{3}{2}\partial^2 c\right)$
\end{itemize}
The anomaly condition is
\[
c_{\mathrm{tot}}=c_{\mathrm{matter}}+c_{bc}=c_{\mathrm{matter}}-26=0.
\]
With the standard intercept, this is exactly the condition
$Q^2=0$.
\end{definition}


\subsection{String theory: Virasoro--BRST duality}

The Virasoro bar complex is computed in
\S\ref{sec:virasoro-bar-moduli-main} above and the WZW bar complex in
Chapter~\ref{chap:kac-moody-koszul}. The BRST comparison uses the
total matter-plus-ghost vertex algebra, not the matter Virasoro
algebra alone.

\subsubsection{Physical states}

\begin{theorem}[Critical bosonic BRST complex \cite{Pol98};
\ClaimStatusProvedElsewhere]
\label{thm:brst-cohomology}
Let $\mathcal V_{\mathrm m}$ be the free-boson matter vertex algebra
with $c_{\mathrm m}=26$, and let $bc_{(2,-1)}$ be the reparametrization
ghost system with $c_{bc}=-26$. The BRST charge
$Q_{\mathrm B}$ on
\[
\mathcal V_{\mathrm{tot}}
=\mathcal V_{\mathrm m}\otimes bc_{(2,-1)}
\]
is square-zero. The physical state space is the relative cohomology
\[
H^\bullet_{\mathrm{rel}}(\mathcal V_{\mathrm{tot}},Q_{\mathrm B})
=
H^\bullet\bigl(\ker b_0\cap\ker L_0^{\mathrm{tot}},\,Q_{\mathrm B}\bigr).
\]
For $c_{\mathrm m}\neq26$ the anomaly term makes
$Q_{\mathrm B}^2\neq0$, so this relative BRST complex is not defined
without an additional anomaly-cancelling sector.
\end{theorem}

\begin{proof}
The OPE of the BRST current with itself has anomaly coefficient
$c_{\mathrm m}-26$. Hence the contour charge satisfies
$Q_{\mathrm B}^2=0$ exactly on the critical surface
$c_{\mathrm m}=26$. The relative complex imposes the constraints
$b_0=0$ and $L_0^{\mathrm{tot}}=0$ before taking cohomology; quotienting
by $Q_{\mathrm B}$-exact vectors gives the physical state space. The
no-ghost and mass-shell descriptions are additional computations in
this relative cohomology, not the definition of the complex.
\end{proof}

\subsubsection{Verifying duality}

\begin{conjecture}[Conditional Virasoro/string comparison; \ClaimStatusConjectured]\label{conj:virasoro-string}
At the critical point, the genus-$0$ BRST/bar comparison relates the
critical bosonic string to the bar complex of $\mathrm{Vir}_{26}$.
The proved statement is the
quasi-isomorphism between the relative BRST complex of the total
critical string vertex algebra and the genus-$0$ bar complex
\textup{(}Theorem~\ref{thm:brst-bar-genus0}\textup{)}. A direct
matter-only identification
\[
H^*(\bar{B}_{\mathrm{geom}}(\mathrm{Vir}_{26})) \cong
H^*_{\mathrm{BRST}}(\mathrm{String})
\]
is a conditional shorthand for that total-string comparison, not a
theorem proved in this chapter.
\end{conjecture}

\begin{remark}[Comparison input]
The proved input is Theorem~\ref{thm:brst-bar-genus0}, which compares
the relative BRST complex of the critical string with the genus-$0$ bar
complex. The additional steps needed to isolate the matter Virasoro
bar complex and treat the ghost sector as an explicit Koszul factor are
not established in the manuscript, so the displayed matter-only
isomorphism remains conditional.
\end{remark}


\subsection{String theory and holographic dualities}

\subsubsection{The algebraic string theory dictionary}

\begin{theorem}[Algebraic bar/BRST genus dictionary;
\ClaimStatusConditional]\label{thm:algebraic-string-dictionary}
\index{string theory!algebraic dictionary}
The algebraic package behind
perturbative string-theory language has the following bar-complex
shadows on the proved comparison surfaces:
\begin{enumerate}[label=\textup{(\roman*)}]
\item \emph{Bar complex / BRST comparison on the proved loci.}
 For Kac--Moody algebras at $k \neq -h^\vee$,
 $H^*(\barB(\cA)) \cong H^{\infty/2+*}(\cA)$
 \textup{(}Theorem~\textup{\ref{thm:bar-semi-infinite-km})}.
 For $c = 26$, the genus-$0$ bar complex is
 quasi-isomorphic to the relative BRST complex
 \textup{(}Theorem~\textup{\ref{thm:brst-bar-genus0})}.
\item \emph{Curvature and anomaly.}
 The obstruction coefficient $\kappa(\cA)$ is the scalar Hodge
 curvature in the modular bar package. On the BRST string surface, the
 nilpotence condition is the separate total-central-charge equation
 $c_{\mathrm{tot}}=0$. For the bosonic free-field sector this is
 $c_{\mathrm m}+c_{bc}=26-26=0$, equivalently
 $\kappa_{\mathrm m}+\kappa_{bc}=13-13=0$ on the Virasoro scalar lane.
	 Free-field/Kac--Moody anomaly duality
	 $\kappa + \kappa' = 0$ holds for affine Kac--Moody and free-field
	 Koszul pairs
	 \textup{(}Corollary~\textup{\ref{cor:anomaly-duality-km-vol1}}\textup{)}.
	 For $\mathcal{W}$-algebras the sum $\kappa + \kappa'$ is
	 nonzero in general, see Remark~\textup{\ref{rem:w3-kappa-sums}}.
\item \emph{Genus grading = loop-order filtration.}
 The genus-graded bar complex
 $\barB^{(g)}(\cA)$ organizes the algebraic genus decomposition and
 the higher-genus correction package
 \textup{(}Theorem~\textup{\ref{thm:genus-universality})}.
\item \emph{Boundary factorization = clutching law.}
 The genus-graded bar classes satisfy the modular-operadic clutching
 identities along boundary divisors of
 $\overline{\mathcal{M}}_{g,n}$, so their separating and
 non-separating residues are expressed in terms of lower-genus bar
 classes
 \textup{(}Theorem~\textup{\ref{thm:genus-induction-strict})}.
\end{enumerate}
\end{theorem}

\begin{proof}
Each clause references its proof:
(i)~Theorems~\ref{thm:bar-semi-infinite-km}
and~\ref{thm:brst-bar-genus0};
(ii)~Corollaries~\ref{cor:anomaly-duality-km-vol1}
and~\ref{cor:anomaly-duality-w},
Theorem~\ref{thm:anomaly-koszul};
(iii)~Theorem~\ref{thm:genus-universality};
(iv)~Theorem~\ref{thm:genus-induction-strict}.
\end{proof}

\begin{remark}[Analytic boundary of the string dictionary]
The theorem records the algebraic dictionary only. The remaining
analytic or physical inputs are the moduli-space integration pairings,
the identification with literal string amplitudes, convergence of the
genus series, and the open/closed bulk-comparison interpretation.
\end{remark}

\begin{remark}[Conditional DS-reduced \texorpdfstring{$\mathcal{W}$}{W} extension]
The principal $\mathcal{W}$-algebra analogue of item~\textup{(i)}
is conditional on Theorem~\ref{thm:bar-semi-infinite-w}. On the
current record, the Kac--Moody and $c=26$ BRST cases are proved,
and the DS-reduced case is precisely the conditional transport clause.
\end{remark}

\begin{corollary}[Conditional genus-\texorpdfstring{$0$}{0} amplitude pairing from the bar complex;
\ClaimStatusConditional]\label{cor:string-amplitude-genus0}
\index{string amplitude!genus zero}
Let $\cA$ be a conformal vertex algebra with $c = 26$, and assume
that the quasi-isomorphism of
Theorem~\ref{thm:brst-bar-genus0} intertwines the physical top-degree
integration functional on $\overline{\mathcal{M}}_{0,n}$ with the
regularized top-degree pairing on the genus-$0$ bar complex. Then the
tree-level amplitude can be computed from the bar complex:
\[
\mathcal{A}_{0,n}^{\mathrm{string}}(V_1, \ldots, V_n)
= \int_{\overline{\mathcal{M}}_{0,n}}
\langle \barBgeom^{(0)}_n
(V_1 \otimes \cdots \otimes V_n) \rangle_{\mathrm{reg}}.
\]
\end{corollary}

\begin{proof}
Tree-level string amplitudes are computed by integration of
BRST-closed vertex operator insertions over
$\overline{\mathcal{M}}_{0,n} \cong
\overline{C}_{n}(\mathbb{P}^1)/\mathrm{PSL}_2$
\cite[Ch.~5]{Pol98}:
$\mathcal{A}_{0,n}^{\mathrm{string}}
= \int_{\overline{\mathcal{M}}_{0,n}}
\omega_{\mathrm{BRST}}(V_1, \ldots, V_n)$.
By Theorem~\ref{thm:brst-bar-genus0}, the chain map
$\Phi\colon (\cA_{\mathrm{tot}}^{\mathrm{rel}}, Q_{\mathrm{BRST}})
\xrightarrow{\sim}
(\barB^{\mathrm{ch}}_0(\cA_{\mathrm{tot}}), d_{\mathrm{bar}})$
is a quasi-isomorphism (for $c = 26$, ensuring
$Q_{\mathrm{BRST}}^2 = 0$). Under the additional pairing-compatibility
hypothesis, the BRST and bar representatives determine the same
top-degree cohomology class against
$[\overline{\mathcal{M}}_{0,n}]$, so the integrals agree.
\end{proof}

\subsubsection{Worldsheet perspective: higher genus}

The genus-$0$ algebraic bridge above has the following genus-$g$
continuation; the physical amplitude identification remains a separate
analytic comparison.

\begin{theorem}[Genus-\texorpdfstring{$g$}{g} chiral homology from bar complex;
\ClaimStatusProvedHere]
\label{thm:genus-g-chiral-homology}
\index{chiral homology!genus g@genus $g$}
For a chiral algebra $\cA$ on a smooth projective
curve~$X$ of genus~$g$, the genus-$g$ bar complex
computes chiral homology:
\[
H_\bullet^{\mathrm{ch}}(X,\cA)
\cong H^\bullet(\barB^{(g)}(X,\cA)).
\]
For a rational $C_2$-cofinite vertex algebra on a stable pointed
curve, the zeroth coinvariant space is the dual of the conformal-block
fiber. These fibers form the usual finite-rank conformal-block sheaf on
the stable locus, with extension across the boundary supplied by the
factorization package.
\end{theorem}

\begin{proof}
Chiral homology is computed via the bar complex by
Theorem~\ref{thm:bar-computes-chiral-homology}.
For $g \geq 1$, the bar differential incorporates
genus-$g$ corrections via the modular operad
(Convention~\ref{conv:higher-genus-differentials})
(Theorem~\ref{thm:genus-induction-strict}); the
resulting complex computes chiral homology on the
genus-$g$ curve by functoriality of factorization homology. For
rational $C_2$-cofinite vertex algebras, Zhu finiteness and
factorization identify the duals of the zeroth coinvariants with the
fibers of the conformal-block sheaf. Coinvariants are not themselves
the conformal blocks; the duality is part of the statement.
\end{proof}

\begin{conjecture}[String amplitude correspondence,
\texorpdfstring{$g \geq 1$}{g 1}; \ClaimStatusConjectured]
\label{conj:string-amplitude-higher-genus}
The chiral homology class of
Theorem~\ref{thm:genus-g-chiral-homology}, integrated
over $\overline{\mathcal{M}}_{g,n}$, computes the
$g$-loop string amplitude:
\[
\mathcal{A}_{g,n}^{\mathrm{string}}
= \int_{\overline{\mathcal{M}}_{g,n}}
\langle \barB^{(g)}_n
(\mathcal{V}_1 \otimes \cdots
\otimes \mathcal{V}_n) \rangle.
\]
(Contributing to Conjecture~\ref{conj:master-bv-brst}.)
\end{conjecture}

\begin{remark}[Higher-genus amplitude comparison]
Theorem~\ref{thm:genus-g-chiral-homology} is proved.
Conjecture~\ref{conj:master-bv-brst} identifies chiral homology
with physical string amplitudes at $g \geq 1$ under the higher-genus
BRST chain map and convergence hypotheses.
At $g = 0$ the bar/BRST comparison is proved, while the
moduli-integration/amplitude identification is only conditional
\textup{(}Corollary~\ref{cor:string-amplitude-genus0}\textup{)}.

The conjecture consists of two chain-level assertions: a compatible
BRST-to-bar morphism over the Deligne--Mumford boundary strata, and
equality of the resulting top-degree trace with the physical
polyform measure. The first assertion is algebraic. The second is an
analytic statement about convergence and regularization of the
string measure.
\end{remark}

\subsubsection{Holographic object firewall}

\begin{conjecture}[Holographic Koszul package; \ClaimStatusConjectured]
\label{conj:bulk-boundary-correspondence}
Let $\cA$ be a boundary chiral algebra with finite-dimensional weight
pieces, convergent completed bar coalgebra, and nondegenerate
Verdier/BPZ pairing. A holographic enhancement of $\cA$ must specify
the distinct objects
\[
\cA,\qquad
B(\cA),\qquad
\cA^{\mathrm i}=H^\bullet B(\cA),\qquad
\cA^!=(\cA^{\mathrm i})^\vee,\qquad
\cZ^{\mathrm{der}}_{\mathrm{ch}}(\cA).
\]
Here $B(\cA)$ is the boundary bar coalgebra, $\cA^!$ is the
Verdier-linear line/defect algebra, and
$\cZ^{\mathrm{der}}_{\mathrm{ch}}(\cA)$ is the chiral Hochschild
cochain-level closed-sector algebra. Physical bulk observables are
identified with it only after a named open/closed comparison
quasi-isomorphism. The cobar counit $\Omega B(\cA)\simeq\cA$ is
inversion; it neither constructs the bulk nor identifies $\cA^!$ with
the derived center. A radial interpretation of the bar filtration is an
extra filtered trace comparison, not a consequence of bar-cobar alone.
(Contributing to Conjecture~\ref{conj:master-bv-brst}.)
\end{conjecture}


\subsection{Segal extension criteria}

\begin{theorem}[Finite BRST-Segal extension criterion \cite{Pol98,Seg04,Zhu96};
\ClaimStatusConditional]\label{thm:classification-extendable}
Let $\mathcal{V}_{\mathrm{m}}$ be a conformal vertex algebra on
$\mathbb{P}^1$ and let $\mathcal{V}_{\mathrm{gh}}$ be the appropriate
ghost system, bosonic or superconformal. Assume that
\[
\mathcal{V}_{\mathrm{tot}}
=\mathcal{V}_{\mathrm{m}}\otimes\mathcal{V}_{\mathrm{gh}}
\]
has a BRST operator $Q$ and finite-dimensional relative cohomology in
each conformal weight. If the following conditions hold, the relative
BRST cohomology defines a finite Segal modular functor and hence
formal perturbative amplitudes satisfying sewing:
\begin{enumerate}
\item \emph{Anomaly cancellation.} The total central charge satisfies
$c_{\mathrm{tot}}=0$; equivalently
$c_{\mathrm{m}}=26$ in the bosonic case and $c_{\mathrm{m}}=15$ in the
superstring case after including the appropriate ghost systems.
\item \emph{Nilpotent BRST charge.} The charge satisfies
$Q^2=0$ on $\mathcal{V}_{\mathrm{tot}}$, so physical states are
$H^\bullet_{\mathrm{rel}}(\mathcal{V}_{\mathrm{tot}},Q)$.
\item \emph{Modular descent.} The vacuum partition function and
inserted conformal blocks descend under the relevant mapping-class
group action.
\item \emph{Finite sewing.} The conformal blocks form finite-rank
local systems with factorization along separating and
non-separating boundary divisors; for rational vertex algebras this is
supplied by Zhu finiteness and the Verlinde package.
\item \emph{Trace-class amplitudes.} The sewing traces converge after
the chosen regularization and are compatible with the boundary
factorization pairings.
\end{enumerate}
Conversely, within the finite-rank BRST-Segal framework these
conditions are necessary. They are not a classification of all chiral
algebras that can appear in string backgrounds.
\end{theorem}

\begin{proof}
\emph{Necessity.}
The Belavin--Knizhnik anomaly is the obstruction to descending the
determinant line over moduli; its coefficient is the total central
charge, hence descent forces $c_{\mathrm{tot}}=0$. Gauge independence
of amplitudes is precisely BRST exactness of the variation with
respect to Beltrami differentials, so $Q^2=0$ is forced. Genus one
descent gives modular covariance of the torus trace; a modular
invariant amplitude requires an invariant vector or trace in that
representation. Boundary factorization forces finite sewing for the
conformal-block sheaves in the finite-rank regime, and the trace-class
condition is the analytic input needed to turn formal sewing into
amplitudes.

\emph{Sufficiency.}
Assumptions (1)--(5) give a finite modular functor. Segal sewing
constructs the genus-$g$ amplitudes by gluing the relative BRST
cohomology along the boundary strata of
$\overline{\mathcal{M}}_{g,n}$. The nilpotence of $Q$ makes the
amplitude independent of gauge slice, anomaly cancellation descends
the determinant line, and modular descent identifies equivalent
sewing presentations. Finite sewing gives the convergence and
factorization input required by
Theorem~\ref{thm:convergence-criterion-spectral}.
\end{proof}

\begin{remark}[Provenance]
Anomaly cancellation is the Belavin--Knizhnik determinant-line
calculation \cite{BK86}. Modular descent of the genus-one trace is the
Friedan--Martinec--Shenker condition. Finiteness and the Verlinde
factorization statement for rational vertex algebras follow from Zhu
\cite{Zhu96} and its modular-functor refinements. BRST nilpotence is
the Kato--Ogawa and Borcherds condition \cite{Bor92}. Segal sewing is
the construction mechanism \cite{Seg04}.
\end{remark}


\subsection{Quantum corrections and deformed Koszul duality}

\begin{theorem}[Genus expansion as curved deformation;
\ClaimStatusConditional]
\label{thm:genus-deformation-exact}
\index{genus expansion!curved deformation}
Let $\cA$ be a complete filtered chiral algebra whose modular
curvature $\kappa(\cA)$ acts centrally, and assume the completion and
finite-type Verdier hypotheses of
Theorem~\ref{thm:curved-koszul-pairs}. The fiberwise
genus-expanded bar operator
\[
d_{\hbar} = \dzero + \sum_{g=1}^\infty \hbar^g\, d_g
\]
on $\barB(\cA)[[\hbar]]$ is a curved coderivation:
\[
d_\hbar^2 = h_\hbar\ast(-),\qquad
h_\hbar=\hbar^2\mcurv{1}+O(\hbar^3),
\]
where $\mcurv{1}=\kappa(\cA)\lambda_1$ is the scalar Hodge curvature
of Convention~\textup{\ref{conv:higher-genus-differentials}}.
The total corrected operator
$\D_\hbar=d_\hbar+d_{\mathrm{per}}+\nabla_{\mathrm{KS}}^{\mathrm{GM}}$
is flat. After passing to the bar-dual coalgebra
$\cA^{\mathrm i}_\hbar=H^\bullet(\barB(\cA)[[\hbar]],\D_\hbar)$ and
then applying Verdier-linear duality, one obtains a completed curved
$A_\infty$ algebra $(\cA^!_\hbar,m_0^\hbar,m_1^\hbar,\ldots)$ whose
curvature term is dual to $h_\hbar$.
\end{theorem}

\begin{proof}
Convention~\ref{conv:higher-genus-differentials} separates the
fiberwise curved operator from the total corrected operator. For the
fiberwise operator, the order-$\hbar^0$ identity is
$\dzero^2=0$. The order-$\hbar$ term is
$\dzero d_1+d_1\dzero=0$, since the genus-one correction is compatible
with collision residues. At order $\hbar^2$ the Fay identity leaves
the central Hodge curvature
$\mcurv{1}=\kappa(\cA)\lambda_1$, giving
$d_\hbar^2=h_\hbar\ast(-)$ in the completed curved convolution
algebra. Higher orders give the remaining terms of $h_\hbar$ by the
same modular-operadic boundary calculation.

The period correction and the Gauss--Manin term cancel this curvature
on the total complex; this is exactly
Theorem~\ref{thm:quantum-diff-squares-zero}. Hence
$\D_\hbar^2=0$, and the cohomology coalgebra
$\cA^{\mathrm i}_\hbar$ is defined. The finite-type Verdier
hypothesis then identifies the continuous dual coalgebra with the
completed curved $A_\infty$ algebra $\cA^!_\hbar$. This is the curved
compatibility statement of
Theorem~\ref{thm:curved-koszul-pairs}, applied to the genus expansion.
\end{proof}

\begin{conjecture}[Loop corrections as string coupling;
\ClaimStatusConjectured]
\label{conj:loop-string-coupling}
In the holographic context, $\hbar = g_s = 1/N$
identifies the algebraic genus expansion parameter with the string
coupling only after a trace functor comparing the bar genus filtration
with the physical loop expansion is supplied. Under such a comparison,
the completed curved Verdier-Koszul dual
$(\cA^!_{g_s},m_0^{g_s},m_1^{g_s},\ldots)$ is conjectured to model
the boundary-side algebraic shadow of loop corrections.
(Contributing to Conjecture~\ref{conj:master-bv-brst}.)
\end{conjecture}

\begin{example}[One-loop correction in \texorpdfstring{AdS$_3$}{AdS3}]
The one-loop correction to the boundary two-point
function:
\[
\langle \mathcal{O}(z) \mathcal{O}(w)
\rangle_{1\text{-loop}} = \frac{1}{N}
\int_{\mathrm{AdS}_3} G(z,w;z')
K(\mathcal{O}^{\vee}, \mathcal{O}^{\vee}, \Phi_{\mathrm{grav}})
\]
where $G$ is the bulk-to-boundary propagator and
$\Phi_{\mathrm{grav}}$ is the graviton field. This is a physical
model for Conjecture~\ref{conj:loop-string-coupling}. The bar-side
curved term would have the form
\[
m_0^{(1)}=\frac{\alpha_{\mathrm{1-loop}}}{N}\lambda_1,
\]
where $\alpha_{\mathrm{1-loop}}$ is fixed by the one-loop determinant
and trace normalization. This coefficient is not computed by Koszul
duality alone.
\end{example}

\begin{remark}[Trace normalization]
Theorem~\ref{thm:genus-deformation-exact} is the algebraic curved
deformation statement under its stated completion hypotheses. The
identification $\hbar=g_s=1/N$ requires the holographic trace
comparison and the one-loop determinant normalization.
\end{remark}


\subsection{String amplitudes via bar complex}

\begin{theorem}[Bar classes on moduli and boundary factorization;
\ClaimStatusProvedHere]
\label{thm:bar-string-integrand}
\index{bar complex!moduli class}
For vertex operators $V_1,\ldots,V_n$, the genus-$g$, $n$-point bar
class
$\barBgeom^{(g)}_n(V_1 \otimes \cdots \otimes V_n)$
satisfies the following properties:
\begin{enumerate}[label=\textup{(\roman*)}]
\item Along a separating boundary divisor
 $D_\Gamma \cong \overline{\mathcal{M}}_{g_1,n_1+1}
 \times \overline{\mathcal{M}}_{g_2,n_2+1}$, its residue is the
 contracted tensor product of the corresponding lower-genus bar
 classes.
\item Along a non-separating boundary divisor
 $D_\Gamma \cong \overline{\mathcal{M}}_{g-1,n+2}$, its residue is
 the self-gluing contraction of the $(g-1,n+2)$ bar class.
\item For the Heisenberg algebra $\mathcal{H}_\kappa$, pairing the bar
 class with the Gaussian measure at all genera computes free boson
 correlation functions by Wick contraction and the determinant formula
 for the oscillator modes.
\end{enumerate}
\end{theorem}

\begin{proof}
Items~(i) and~(ii) are the separating and non-separating boundary
components of the modular operad structure on the genus-graded bar
package: the bar differential on stable curves is assembled from
boundary residues, and Theorem~\ref{thm:genus-induction-strict}
identifies those residues with the corresponding clutching
contractions of lower-genus bar classes. Item~(iii) is the Gaussian
calculation: Wick's theorem contracts the bar tensor factors pairwise,
and the oscillator determinant gives the Dedekind $\eta$-factor.
\end{proof}

\begin{conjecture}[String amplitude formula;
\ClaimStatusConjectured]\label{conj:string-amplitude}
For a general interacting chiral algebra $\cA$ at
$g \geq 1$:
\[
\mathcal{A}_{g,n}^{\mathrm{string}}(V_1, \ldots, V_n)
= \int_{\overline{\mathcal{M}}_{g,n}}
\langle \barBgeom^{(g)}_n
(V_1 \otimes \cdots \otimes V_n)
\rangle_{\mathrm{reg}}
\]
where $\langle \cdot \rangle_{\mathrm{reg}}$ denotes
Costello regularization of the string measure
$\det(\mathrm{Im}\,\Omega)^{-c/2}$.
(Contributing to Conjecture~\ref{conj:master-bv-brst}.)
\end{conjecture}


\begin{theorem}[String factorization and bar residues \cite{Pol98};
\ClaimStatusConditional]\label{thm:amplitude-factorization}
Physical string amplitudes factorize at a stable boundary divisor
$D_\Gamma$ by summing over intermediate BRST states:
\[
\operatorname{Res}_{D_\Gamma}\mathcal{A}_{g,n}^{\mathrm{string}}
=
\sum_i
\mathcal{A}_{g_1,n_1+1}^{\mathrm{string}}(V_I,V_i)\,
\mathcal{A}_{g_2,n_2+1}^{\mathrm{string}}(V_i^*,V_J),
\]
with the analogous self-gluing formula on non-separating divisors.
If a higher-genus BRST/bar comparison map intertwines the BPZ
propagator with the bar residue pairing, then the bar class satisfies
\[
\Phi\!\left(
\operatorname{Res}_{D_\Gamma}[\barBgeom^{(g)}_n]
\right)
=
\operatorname{Res}_{D_\Gamma}\mathcal{A}_{g,n}^{\mathrm{string}}.
\]
The first statement is the physical factorization theorem; the second
is the conditional bar comparison.
\end{theorem}

\begin{proof}
By Theorem~\ref{thm:boundary-higher-genus}, the boundary of
$\overline{\mathcal{M}}_{g,n}$ decomposes into divisors $D_\Gamma$
indexed by stable graphs $\Gamma$. A separating node gives
\[
D_\Gamma \cong \overline{\mathcal{M}}_{g_1, n_1+1} \times \overline{\mathcal{M}}_{g_2, n_2+1}
\]
where $g_1 + g_2 = g$, $n_1 + n_2 = n$, and the extra marked points on each factor correspond to the node. Non-separating nodes give $D_\Gamma \cong \overline{\mathcal{M}}_{g-1, n+2}$.

The string amplitude is $\mathcal{A}_{g,n} = \int_{\overline{\mathcal{M}}_{g,n}} \Omega_{g,n}$ where $\Omega_{g,n}$ is the top form constructed from vertex operators and the string measure. Near a boundary divisor $D_\Gamma$ with local parameter $t$ (the plumbing fixture), the integrand has a pole:
\[
\Omega_{g,n} \sim \frac{dt}{t} \wedge \Omega_{g_1, n_1+1} \wedge \Omega_{g_2, n_2+1} \cdot \sum_i V_i \otimes V_i^*
\]
where the sum is over a complete set of intermediate states $\{V_i\}$ (the propagator), and $V_i^*$ denotes the dual state under the BPZ pairing \cite{BPZ84}.

The geometric bar complex has its own boundary residue:
\[
\operatorname{Res}_{D_\Gamma}[\barBgeom^{(g)}_n] = \barBgeom^{(g_1)}_{n_1} \otimes_{\mathrm{Prop}} \barBgeom^{(g_2)}_{n_2}
\]
where $\otimes_{\mathrm{Prop}}$ is the bar pairing along the node.
Identifying this algebraic pairing with the BPZ propagator is exactly
the extra BRST/bar compatibility hypothesis.
\index{state-field correspondence}
\end{proof}

\begin{example}[Tree-level four-point amplitude]
For the tree-level four-point amplitude in closed string theory, the bar complex is
\[\barBgeom^{(0)}_4 = \text{span}\{V_1 \otimes V_2 \otimes V_3 \otimes V_4 \otimes \eta_{12} \wedge \eta_{23} \wedge \eta_{34}\}\]
and the amplitude is
\[\mathcal{A}_{0,4} = \int_{\overline{C}_4(\mathbb{P}^1)} \frac{dz_1 \wedge dz_2 \wedge dz_3 \wedge dz_4}{(z_1-z_2)(z_2-z_3)(z_3-z_4)(z_4-z_1)} \prod_{i=1}^4 V_i(z_i)\]

This gives a Veneziano-type amplitude:
\[\mathcal{A}_{0,4} \sim B(-\alpha' s, -\alpha' t) = \frac{\Gamma(-\alpha' s)\,\Gamma(-\alpha' t)}{\Gamma(-\alpha' s - \alpha' t)}\]
where $s, t$ are the Mandelstam variables and $\alpha'$ is the Regge slope (the precise form depends on the vertex operator insertions).
\end{example}

\begin{example}[One-loop two-point amplitude]
For the one-loop two-point amplitude, the bar complex is
\[\barBgeom^{(1)}_2 = \text{span}\{V_1 \otimes V_2 \otimes \eta_{12} \otimes \omega_{\text{moduli}}\}\]
where $\omega_{\text{moduli}} = d\tau \wedge d\bar{\tau}/(\text{Im}\tau)^2$ is the K\"ahler form on $\mathcal{M}_1$.
The amplitude is
\[\mathcal{A}_{1,2} = \int_{\mathcal{M}_1} \frac{d\tau \wedge d\bar{\tau}}{(\text{Im}\tau)^2} \int_{\mathbb{T}_\tau} \frac{dz_1 \wedge dz_2}{(z_1-z_2)^2} V_1(z_1) V_2(z_2)\]

This gives the one-loop correction with modular invariance.
\end{example}

\begin{theorem}[Belavin--Knizhnik anomaly condition;
\ClaimStatusProvedElsewhere]\label{thm:modular-anomaly}
For a matter-plus-ghost worldsheet theory, the determinant-line
anomaly at genus~$1$ has scalar class
\[
\frac{c_{\mathrm{tot}}}{24}\lambda_1,
\qquad
c_{\mathrm{tot}}=c_{\mathrm{matter}}+c_{\mathrm{ghost}}.
\]
For the bosonic string, $c_{\mathrm{ghost}}=-26$, so the anomaly
vanishes exactly when $c_{\mathrm{matter}}=26$. For the superstring,
the corresponding matter value is $c_{\mathrm{matter}}=15$ after the
superconformal ghost sector is included. This vanishing is necessary
for Weyl and modular descent of amplitudes; a scalar modular-invariant
partition function still requires the representation-theoretic descent
data of the modular classification theorem.
\end{theorem}

\begin{proof}
The Belavin--Knizhnik calculation identifies the curvature of the
determinant line with the total central charge divided by~$24$. The
reparametrization ghost system has $c_{bc}=-26$, giving
$c_{\mathrm{tot}}=c_{\mathrm{matter}}-26$ in the bosonic case. The
superconformal ghost system gives the analogous $15$ in the
superstring case. This proves the anomaly cancellation condition. It
does not by itself produce a modular-invariant scalar partition
function; that is a separate question about the
$\mathrm{SL}_2(\mathbb Z)$ representation carried by the conformal
blocks.
\end{proof}


\subsection{\texorpdfstring{Genus-one modular covariance}{Genus-one modular covariance}}

\begin{theorem}[Genus-one bar character covariance;
\ClaimStatusConditional]\label{thm:modular-invariance}
Let $\mathcal A$ be a finite-type rational $C_2$-cofinite vertex
algebra with simple modules $M_i$. The genus-$1$ bar traces
\[
Z_i(\tau)=\operatorname{Tr}_{M_i}\bigl(q^{L_0-c/24}\bigr),
\qquad q=e^{2\pi i\tau},
\]
span a finite-dimensional representation of
$\mathrm{SL}_2(\mathbb Z)$. Equivalently,
\[
Z_i(\gamma\tau)=\sum_j \rho(\gamma)_{ij}Z_j(\tau).
\]
The factor $q^{-c/24}$ is the determinant-line normalization. It is
not a scalar modular weight for the bar complex, and a scalar
partition function is invariant only after choosing an invariant
linear combination of the character vector.
\end{theorem}

\begin{proof}[Proof via genus-1 bar complex]
The genus-$1$ bar trace is the ordinary VOA character trace with the
determinant-line normalization $q^{-c/24}$. Zhu's theorem gives the
finite-dimensional modular representation for rational $C_2$-cofinite
VOAs. The geometric bar representative is built from the Weierstrass
$\zeta$-function:
\[
G_\tau(z) = \zeta_\tau(z) - \frac{\pi^2 E_2(\tau)}{3}\,z
\]
where $E_2(\tau)$ is the quasi-modular Eisenstein series. Under $\tau \mapsto \gamma\cdot\tau = (a\tau+b)/(c\tau+d)$:
\[
G_{\gamma\cdot\tau}\!\left(\frac{z}{c\tau+d}\right) = (c\tau+d)\,G_\tau(z) + 2\pi i c\,z
\]
The quasi-modular anomaly of $E_2$ is absorbed by the flat
genus-$1$ connection and produces the same cocycle as the character
representation. Hence the bar trace is covariant as a vector-valued
character. Scalar invariance is an additional invariant-vector
condition.
\end{proof}

\begin{theorem}[Modular anomaly for affine Kac--Moody algebras;
\ClaimStatusConditional]
\label{thm:modular-anomaly-km-w}
\index{modular anomaly!KM and W-algebras}
Let $\cA=V_k(\fg)$ be an affine Kac--Moody vertex algebra with
Sugawara central charge $c=k\dim\fg/(k+h^\vee)$ on an integrable
positive-level WZW surface. Then:
\begin{enumerate}[label=\textup{(\roman*)}]
\item The genus-$1$ character vector transforms under the WZW
 $\mathrm{SL}_2(\bZ)$ representation
 \textup{(}Theorem~\textup{\ref{thm:modular-invariance})}.
\item BRST\slash semi-infinite
 cohomology is computed via the bar complex:
 $H^*(\barB^{\mathrm{ch}}(\cA)) \cong
 H^*_{\mathrm{si}}(\cA)$
 \textup{(}Theorem~\textup{\ref{thm:bar-semi-infinite-km})}.
\item The anomaly coefficient satisfies $\kappa + \kappa' = 0$ for affine Kac--Moody and free-field Koszul pairs
 \textup{(}Proposition~\textup{\ref{prop:ff-channel-shear})};
 in particular, the affine Kac--Moody duality is antisymmetric.
\end{enumerate}
For affine Kac--Moody, the intrinsic vanishing
$\kappa(\cA) = 0$ occurs at the critical level $k = -h^\vee$.
The physical string anomaly condition is separate: after coupling to
the bosonic $bc$ ghosts it is $c_{\mathrm{matter}}-26=0$.
\end{theorem}

\begin{proof}
Combine the three established results.
By~(i), the genus-$1$ bar characters transform as the WZW
character vector. By~(ii), this complex computes the relevant
semi-infinite\slash BRST cohomology on the affine Kac--Moody side.
By~(iii), the anomaly
coefficient $\kappa = (k+h^\vee)\dim(\fg)/(2h^\vee)$ for
$\widehat{\fg}_k$ satisfies $\kappa + \kappa' = 0$ for affine Kac--Moody
under $k \mapsto -k - 2h^\vee$. The critical level condition
$k=-h^\vee$ is therefore an intrinsic Koszul-curvature condition.
The bosonic string condition instead uses the ghost contribution:
the $bc$ system at weights $(2,-1)$ has $c_{\mathrm{gh}}=-26$, so
total anomaly cancellation requires $c_{\mathrm{matter}}=26$.
\end{proof}

\begin{remark}[Conditional principal \texorpdfstring{$\mathcal{W}$}{W} extension]
The principal $\mathcal{W}$-algebra analogue of
Theorem~\ref{thm:modular-anomaly-km-w} is conditional on the
DS/bar transport statement
\textup{(}Theorem~\textup{\ref{thm:bar-semi-infinite-w}}\textup{)}.
What is already proved for principal $\mathcal{W}$-pairs on the
current record is the curvature-complementarity package of
Corollary~\ref{cor:anomaly-duality-w}, not the full
semi-infinite/BRST identification.
\end{remark}

\begin{conjecture}[Determinant-line anomaly for general chiral
algebras; \ClaimStatusConjectured]
\label{conj:modular-anomaly-brst}
For a general chiral algebra $\cA$ with a finite-type relative BRST
complex and trace-class genus-$1$ amplitudes, the determinant-line
part of the modular anomaly is
\[
\frac{c_{\mathrm{tot}}}{24}\lambda_1.
\]
The remaining anomaly data are representation-valued modular
cocycles of the conformal-block sheaf. No universal formula depending
only on $c$ and a regularized dimension of BRST cohomology is asserted.
(Contributing to Conjecture~\ref{conj:master-bv-brst}.)
\end{conjecture}


\begin{example}[Virasoro algebra modular invariance]
The Virasoro bar complex is computed in
\S\ref{sec:virasoro-bar-moduli-main}. For holomorphic self-dual CFTs,
a scalar vacuum partition function requires the usual
$c\equiv0\pmod{24}$ condition, as in the Monster module at $c=24$
\cite{FLM88}. The string-critical values $c_{\mathrm m}=26$ and
$c_{\mathrm m}=15$ are matter-plus-ghost anomaly conditions, not
standalone holomorphic genus-$1$ modular-invariance conditions.
\end{example}

\begin{example}[WZW model modular invariance]
The WZW/Kac--Moody bar complex is computed in Chapter~\ref{chap:kac-moody-koszul}; the central charge is $c = k\dim\mathfrak{g}/(k + h^\vee)$.
For integrable positive integer levels, the WZW characters form the
standard finite-dimensional modular representation; diagonal or
ADE-type modular invariants are additional invariant tensors in that
representation. Coupling such a matter sector to the bosonic string
ghosts imposes the separate equation $c-26=0$.
\end{example}

\begin{theorem}[Genus-one modular descent criteria \cite{Zhu96, MooreSeiberg89};
\ClaimStatusConditional]\label{thm:modular-classification}
The following mechanisms give genus-$1$ modular descent in their
stated finite-type settings:

\begin{enumerate}
\item \emph{Holomorphic self-dual theories.} A scalar vacuum
 character descends when the determinant-line multiplier is trivial;
 in the standard holomorphic case this forces $c\in24\mathbb Z$.
\item \emph{Rational $C_2$-cofinite theories.} The characters form a
 finite-dimensional $\mathrm{SL}_2(\mathbb Z)$ representation, and
 modular-invariant partition functions are invariant tensors in that
 representation.
\item \emph{Critical BRST strings.} Matter plus ghosts descend when
 $c_{\mathrm{tot}}=0$ and the relative BRST cohomology is finite and
 sewing-compatible.
\item \emph{Orbifolds.} A finite orbifold descends when the group
 action is anomaly-free and the twisted-sector modular data close.
\end{enumerate}
These criteria are not mutually exclusive and not exhaustive; they
replace no classification theorem.
\end{theorem}

\begin{proof}[Proof via representation theory]
For holomorphic self-dual theories, the only character is the vacuum
character, so descent requires cancellation of the determinant-line
multiplier; the standard normalization gives $c\in24\mathbb Z$. For
rational $C_2$-cofinite VOAs, Zhu's theorem gives the finite modular
representation, and Moore--Seiberg compatibility identifies the
invariant tensors with modular-invariant partition functions. For
critical BRST strings, Theorem~\ref{thm:modular-anomaly} supplies the
determinant-line cancellation and the finite BRST-Segal criterion
supplies sewing. The orbifold statement is the corresponding
equivariant modular-functor condition. None of these arguments proves
a bidirectional classification of chiral algebras.
\end{proof}



\section{W-algebras and Wakimoto representations}\label{sec:w-algebras-wakimoto}

The $\mathcal{W}$-algebra construction via Drinfeld--Sokolov reduction,
its Slodowy slice geometry, the nilpotent orbit poset, and the full
bar complex computation are developed in
Chapter~\ref{chap:w-algebras}. The connection between
the Wakimoto free-field realization and the bar complex is as follows.

\begin{theorem}[Filtered Wakimoto bar model;
\ClaimStatusConditional]\label{thm:wakimoto-bar}
Let $\mathcal{M}_{\mathrm{Wak}}$ be the Wakimoto free-field module of
$\widehat{\mathfrak{g}}_k$. In the associated free-field bar model one
has
\begin{itemize}
\item Degree 0: $H^0 = \mathbb{C}[\phi_1, \ldots, \phi_r]$ in the
 Cartan polynomial sector.
\item Degree 1: the free-field generators lie in
 $\bigoplus_{\alpha \in \Delta_+} \mathbb{C}\beta_\alpha
 \oplus \bigoplus_{i=1}^r \mathbb{C}\partial\phi_i$ before imposing
 the Drinfeld--Sokolov differential.
\item If the Drinfeld--Sokolov BRST differential is compatible with
 the bar filtration and the resulting spectral sequence converges,
 the reduced cohomology gives the corresponding filtered
 $\mathcal{W}$-model.
\end{itemize}
The last clause is the reduction hypothesis; it is not a statement
that the unreduced Wakimoto bar complex is already
$\mathcal{W}^{-h^\vee}(\mathfrak g)$.
\end{theorem}

\begin{proof}
The Wakimoto module decomposes as
$\mathcal{M}_{\mathrm{Wak}} =
\bigotimes_{\alpha \in \Delta_+}(\beta\gamma)_\alpha\otimes
\mathcal{H}_{\mathfrak h}$. The free-field bar cohomology is therefore
the tensor product of the $\beta\gamma$ and Heisenberg computations
from~\S\S\ref{sec:bar-complexes-free-fields}--\ref{sec:heisenberg-bar-complex-sec2}.
The passage to a $\mathcal W$-algebra requires the
Drinfeld--Sokolov BRST differential. If $Q_{\mathrm{DS}}$ preserves the
bar filtration and the induced spectral sequence converges, its
$E_1$ page is the displayed free-field bar cohomology and its abutment
is the reduced bar model. This proves the conditional statement and
isolates the comparison still supplied by Chapter~\ref{chap:w-algebras}.
\end{proof}

\begin{proposition}[Graphical Wakimoto model;
\ClaimStatusConditional]\label{prop:wakimoto-graph}
After choosing the Wakimoto normal ordering and the compatible
bar-filtration data of Theorem~\ref{thm:wakimoto-bar}, the free-field
bar model is represented by decorated graphs: vertices carry Wakimoto
generators, edges carry log-form propagators
$\eta_{ij}=d\log(z_i-z_j)$, and the differential is Wick contraction
together with the induced Drinfeld--Sokolov term.
\end{proposition}

\begin{theorem}[Transferred \texorpdfstring{$A_\infty$}{A-infinity} model for
\texorpdfstring{$\mathcal{W}$}{W}-algebras; \ClaimStatusConditional]
\label{thm:w-algebra-ainfty}
\index{W-algebra!$A_\infty$ structure}
Assume that the Drinfeld--Sokolov bar complex of
$\mathcal{W}^k(\mathfrak g)$ admits a filtered contraction onto a
chosen minimal model. Homotopy transfer gives operations
$\{m_n\}_{n\ge1}$ with $m_1=0$ on that model; $m_2$ is induced by the
OPE product, and $m_n$ for $n\ge3$ is the contribution of the
$n$-fold collision strata and propagator homotopies. The
$A_\infty$ relations hold on the transferred model because the
source differential squares to zero.
\end{theorem}

\begin{conjecture}[W-algebra \texorpdfstring{$A_\infty$}{A-infinity} and quantum
cohomology]
\label{conj:w-algebra-quantum-cohomology}
\ClaimStatusConjectured
The conditional $A_\infty$ operations of
Theorem~\ref{thm:w-algebra-ainfty} are conjectured to encode the quantum
cohomology of the flag variety: $m_k$ corresponds to
contributions from Schubert cells in~$G/B$. For
$\fg = \mathfrak{sl}_n$, this recovers $QH^*(G/B)$
with quantum parameter $q = e^{2\pi i \tau}$.
\end{conjecture}

\begin{remark}[Quantum cohomology comparison]
The transferred $A_\infty$ model of
Theorem~\ref{thm:w-algebra-ainfty} is conditional on the filtered
contraction. Conjecture~\ref{conj:w-algebra-quantum-cohomology}
identifies it with $QH^*(G/B)$ under the additional matching of
Schubert cell combinatorics with collision strata.
\end{remark}

\begin{theorem}[Classical \texorpdfstring{$\mathcal{W}$}{W}-algebra integrability; \ClaimStatusProvedElsewhere]\label{thm:w-classical-integrability}
The classical W-algebra $\mathcal{W}^{\mathrm{cl}}(\mathfrak{g})$ encodes integrability of the generalized Toda lattice via the Drinfeld--Sokolov construction \cite{DS}: the $m_2$ product gives the Poisson bracket, and the integrals of motion Poisson-commute.
For $\mathfrak{g} = \mathfrak{sl}_n$, this recovers the $\mathfrak{sl}_n$ Toda hierarchy.
\end{theorem}

\begin{theorem}[Higher \texorpdfstring{$A_\infty$}{A-infinity} corrections in quantum \texorpdfstring{$\mathcal{W}$}{W}-algebras; \ClaimStatusConditional]\label{thm:w-integrability}
Under the hypotheses of Theorem~\ref{thm:w-algebra-ainfty}, the
transferred quantum $\mathcal{W}$-algebra model provides a
filtered deformation of the classical binary Drinfeld--Sokolov
package: in the classical limit the binary operation recovers the
Poisson bracket of
Theorem~\ref{thm:w-classical-integrability}, while in the quantum
algebra the higher operations $m_k$ \textup{(}$k \geq 3$\textup{)}
record the non-quadratic correction terms detected by the Li/PBW
filtration. This is the algebraic deformation package underlying
possible quantum-integrability comparisons, but no identification
with a specific quantum integrable hierarchy is proved here.
\end{theorem}

\begin{proof}
Under the filtered contraction hypothesis, the existence of the higher
operations $m_k$ is
Theorem~\ref{thm:w-algebra-ainfty}. The classical binary
Drinfeld--Sokolov integrable structure is the content of
Theorem~\ref{thm:w-classical-integrability}. Passing from the
classical Poisson algebra to the quantum $\mathcal{W}$-algebra
introduces the higher multilinear operations supplied by the
bar-side $A_\infty$ model; these higher terms are the non-binary
correction data measured by the Li/PBW filtration in that model.
\end{proof}



\section{Comparison and filtered structures}\label{sec:comparison-filtered}

Free-field examples isolate the Koszul residue/BPZ comparison surface,
the interaction between filtrations and bar complexes, and the
five-theorem projections of the universal MC element.

\subsection{Comparison of free-field examples}

Geometric complexity correlates inversely with algebraic simplicity:
the free fermion has no positive-degree genus-$0$ bar cohomology; the
$\beta\gamma$ system has bar chain sectors of rank
$2 \cdot 3^{n-1}$ and bar cohomology governed by
$\sqrt{(1+x)/(1-3x)}$; the Heisenberg gains a level central class;
the Virasoro connects to $\overline{\mathcal{M}}_{0,n}$; W-algebras
link to flag varieties.

\begin{remark}[Five projections of $\Theta_\cA$: free-field archetypes]
\label{rem:free-field-five-theorems}
\index{main theorems!free field verification}
\index{Chriss--Ginzburg structure principle!free fields}
Each free-field archetype verifies all five main theorems as projections
of the universal MC element
$\Theta_\cA \in \MC(\gAmod)$
(Theorem~\ref{thm:mc2-bar-intrinsic}). The verification tables below
organize the proved data per family; the shadow archetype column
refers to the classification of
Table~\ref{tab:shadow-tower-census}.
\end{remark}

\subsubsection{Heisenberg: five-theorem verification}

\begin{computation}[Five projections of $\Theta_{\cH_k}$;
\ClaimStatusConditional]
\label{comp:heisenberg-five-theorems}
\index{Heisenberg algebra!five-theorem verification}
Shadow archetype: G (Gaussian, $r_{\max} = 2$).
$\kappa(\cH_k) = k$.
Central charge $c = 1$.
\begin{center}
\small
\renewcommand{\arraystretch}{1.3}
\begin{tabular}{lll}
\toprule
\emph{Theorem} & \emph{Projection of $\Theta_{\cH_k}$}
 & \emph{Content} \\
\midrule
A (bar-cobar) &
 $\Theta_{\cH_k}\big|_{\hbar=0} = \tau$ &
 $\cH_k^! = \mathrm{Sym}^{\mathrm{ch}}(V^*)$ (curved, $m_0 = -k\omega$) \\
B (inversion) &
 $\Omega(\barBgeom(\cH_k)) \simeq \cH_k$ &
 cobar counit; dual algebra computed in Thm~\ref{thm:heisenberg-koszul-dual-early} \\
C (complementarity) &
 $\kappa(\cH_k) + \kappa(\cH_k^!) = 0$ \textup{(}free-field pair\textup{)} &
 scalar obstruction cancellation; $B(\cH_k)$, $\cH_k^!$, and
 $Z_{\mathrm{ch}}^{\mathrm{der}}(\cH_k)$ are distinct objects \\
D (modular char.) &
 $\Theta_{\cH_k}\big|_{g=1,n=0} = k \cdot \lambda_1$ &
$\kappa(\cH_k) = k$;\; $F_g = k \cdot \lambda_g^{\mathrm{FP}}$ \textup{(}UNIFORM-WEIGHT\textup{}) \\
H (chiral Hochschild) &
 $\ChirHoch^{>2}(\cH_k)=0$ &
 $P_{\cH_k}(t)=1+t+t^2$ \\
\bottomrule
\end{tabular}
\end{center}
For Theorem~H the Heisenberg computation is explicit:
\[
\ChirHoch^0(\cH_k)=\mathbb{C},\qquad
\ChirHoch^1(\cH_k)=\mathbb{C}\cdot[D],\qquad
\ChirHoch^2(\cH_k)=\mathbb{C}\cdot[\partial_k],\qquad
\ChirHoch^n(\cH_k)=0 \textup{ for } n \ge 3,
\]
where $D(J)=\mathbf{1}$ is the outer derivation and $[\partial_k]$
is the infinitesimal level deformation. The bar resolution has one
cogenerator in bar degree~$1$ and its quadratic square in bar
degree~$2$, so applying
$\operatorname{Hom}_{\cH_k\textup{-}\cH_k}(-,\cH_k)$ produces a
cochain complex concentrated in degrees $0,1,2$. Hence
\[
P_{\cH_k}(t)=\sum_{n \ge 0}\dim \ChirHoch^n(\cH_k)\,t^n = 1+t+t^2,
\]
which is the Heisenberg specialization of
Theorem~\ref{thm:hochschild-polynomial-growth}; the one-dimensional
degree-$0$ and degree-$2$ terms match the shift-$[2]$ symmetry of
Theorem~\ref{thm:main-koszul-hoch}.
The tower terminates at $r = 2$: $\Theta_{\cH_k}^{\leq r}
= \Theta_{\cH_k}^{\leq 2} = \kappa \cdot \eta \otimes \Lambda$
for all $r \geq 2$. No cubic, quartic, or higher shadows.
All higher $L_\infty$ brackets vanish:
$\ell_k^{\mathrm{tr}} = 0$ for $k \geq 3$.
The modular convolution algebra $\gAmod$ is strictly formal.
\end{computation}

\subsubsection{Free fermion: five-theorem verification}

\begin{computation}[Five projections of $\Theta_{\mathcal{F}}$;
\ClaimStatusProvedHere]
\label{comp:fermion-five-theorems}
\index{free fermion!five-theorem verification}
Shadow archetype: G (Gaussian, $r_{\max} = 2$).
$\kappa(\mathcal{F}) = 1/4$.
Central charge $c = 1/2$.
\begin{center}
\small
\renewcommand{\arraystretch}{1.3}
\begin{tabular}{lll}
\toprule
\emph{Theorem} & \emph{Projection of $\Theta_{\mathcal{F}}$}
 & \emph{Content} \\
\midrule
A (bar-cobar) &
 $\Theta_\mathcal{F}\big|_{\hbar=0} = \tau$ &
 $\mathcal{F}^! = \mathrm{Sym}^{\mathrm{ch}}(\gamma)$ (Thm~\ref{thm:single-fermion-boson-duality}) \\
B (inversion) &
 $\Omega(\barBgeom(\mathcal{F})) \simeq \mathcal{F}$ &
 cobar counit; genus-$0$ cohomology collapses in Thm~\ref{thm:fermion-bar-complex-genus-0} \\
C (complementarity) &
 $\kappa(\mathcal{F}) + \kappa(\mathcal{F}^!) = 0$ \textup{(}free-field pair\textup{)} &
 scalar obstruction cancellation; $B(\mathcal F)$,
 $\mathcal F^{\mathrm i}$, $\mathcal F^!$, and
 $Z_{\mathrm{ch}}^{\mathrm{der}}(\mathcal F)$ are distinct \\
D (modular char.) &
 $\Theta_\mathcal{F}\big|_{g=1,n=0} = \tfrac{1}{4}\lambda_1$ &
 $\kappa(\mathcal{F}) = c/2 = 1/4$ \\
H (chiral Hochschild) &
 Degree-preserving sub-MC &
 $\ChirHoch^0(\mathcal{F})=\mathbb{C}$,\;
 $\ChirHoch^1(\mathcal{F})=0$,\;
 $\ChirHoch^2(\mathcal{F})=\mathbb{C}$ \\
\bottomrule
\end{tabular}
\end{center}
The tower terminates at $r = 2$.
Fermionic statistics kill all positive-degree genus-$0$ bar cohomology
(Theorem~\ref{thm:fermion-bar-complex-genus-0}):
the cyclic permutation on $\overline{C}_n(X)$ acts by
primitive roots of unity, leaving no invariant forms.
\end{computation}

\subsubsection{$\beta\gamma$ system: five-theorem verification}

\begin{computation}[Five projections of $\Theta_{\beta\gamma}$;
\ClaimStatusConditional]
\label{comp:betagamma-five-theorems}
\index{beta-gamma system@$\beta\gamma$ system!five-theorem verification}
Shadow archetype: C (contact/quartic, $r_{\max} = 4$).
$\kappa(\beta\gamma) = c/2 = 6\lambda^2 - 6\lambda + 1$.
Central charge $c_{\beta\gamma} = 2(6\lambda^2 - 6\lambda + 1)$.
\begin{center}
\small
\renewcommand{\arraystretch}{1.3}
\begin{tabular}{lll}
\toprule
\emph{Theorem} & \emph{Projection of $\Theta_{\beta\gamma}$}
 & \emph{Content} \\
\midrule
A (bar-cobar) &
 $\Theta_{\beta\gamma}\big|_{\hbar=0} = \tau$ &
 $(\beta\gamma)^! = bc$ (Thm~\ref{thm:betagamma-bc-koszul}) \\
B (inversion) &
 $\Omega(\barBgeom(\beta\gamma)) \simeq \beta\gamma$ &
 cobar counit; lowest-weight chain rank $2 \cdot 3^{n-1}$
 (Thm~\ref{thm:betagamma-bar-dim}) \\
C (complementarity) &
 $\kappa(\beta\gamma) + \kappa(bc) = 0$ \textup{(}free-field pair\textup{)} &
 scalar obstruction cancellation; the dual algebra is $bc$, not the
 derived center of $\beta\gamma$ \\
D (modular char.) &
 $\Theta_{\beta\gamma}\big|_{g=1,n=0}
 = \kappa \cdot \lambda_1$ &
 $\kappa = 6\lambda^2{-}6\lambda{+}1$;\; $= 1$ at $\lambda=0,1$;\; $= -\tfrac{1}{2}$ at $\lambda=\tfrac{1}{2}$ \\
H (chiral Hochschild) &
 Degree-preserving sub-MC &
 $\ChirHoch^0(\beta\gamma_\lambda)=\mathbb{C}$,\;
 $\ChirHoch^1(\beta\gamma_\lambda)=0$,\;
 $\ChirHoch^2(\beta\gamma_\lambda)=\mathbb{C}$ \\
\bottomrule
\end{tabular}
\end{center}
Shadow obstruction tower: the internal weight-changing line is
shadow-trivial, but the standard conformal-weight family
$\beta\gamma_\lambda$ satisfies
$S_2 = 6\lambda^2 - 6\lambda + 1$, $S_3 = 0$, $S_4 = -5/12$,
and $S_r = 0$ for $r \ge 5$
\textup{(}Chapter~\ref{chap:beta-gamma}; Proposition~\ref{prop:depth-gap-trichotomy}\textup{)}.
Hence the tower terminates at exactly $r = 4$.
The full shadow obstruction tower data (multi-channel structure,
special weight table, shadow metric, weight symmetry,
Mumford connection) is in
\S\ref{sec:betagamma-shadow-tower-free}.
The two-generator free-field systems ($\beta\gamma$ and $bc$)
are the free-field archetypes outside the Gaussian class.
Their conformal-weight line is a stress-tensor/shadow parameter, not
a curve-level $\ChirHoch^1$ class.
The infinite sequence in
Remark~\ref{rem:betagamma-bar-cohomology-gf} is bar cohomology, not
chiral Hochschild cohomology; Theorem~H concerns the latter.
\end{computation}

\subsubsection{$bc$ ghost system: five-theorem verification}

\begin{computation}[Five projections of $\Theta_{bc}$;
\ClaimStatusConditional]
\label{comp:bc-five-theorems}
\index{$bc$ system!five-theorem verification}
Shadow archetype: C (contact/quartic, $r_{\max} = 4$),
dual to $\beta\gamma$.
$\kappa(bc) = c_{bc}/2 = -(6\lambda^2 - 6\lambda + 1)$.
At the physical point $\lambda = 2$: $\kappa(bc) = -13$, $c_{bc} = -26$.
\begin{center}
\small
\renewcommand{\arraystretch}{1.3}
\begin{tabular}{lll}
\toprule
\emph{Theorem} & \emph{Projection of $\Theta_{bc}$}
 & \emph{Content} \\
\midrule
A (bar-cobar) &
 $\Theta_{bc}\big|_{\hbar=0} = \tau$ &
 $(bc)^! = \beta\gamma$ (Thm~\ref{thm:betagamma-bc-koszul}) \\
B (inversion) &
 $\Omega(\barBgeom(bc)) \simeq bc$ &
 cobar counit \\
C (complementarity) &
 $\kappa(bc) + \kappa(\beta\gamma) = 0$ \textup{(}free-field pair\textup{)} &
 scalar obstruction cancellation; the dual algebra is
 $\beta\gamma$, not the derived center of $bc$ \\
D (modular char.) &
 $\Theta_{bc}\big|_{g=1,n=0}
 = \kappa(bc) \cdot \lambda_1$ &
 $\kappa(bc) = -\kappa(\beta\gamma) = -(6\lambda^2{-}6\lambda{+}1)$ \\
H (chiral Hochschild) &
 Degree-preserving sub-MC &
 $\ChirHoch^0(bc_\lambda)=\mathbb{C}$,\;
 $\ChirHoch^1(bc_\lambda)=0$,\;
 $\ChirHoch^2(bc_\lambda)=\mathbb{C}$ \\
\bottomrule
\end{tabular}
\end{center}
By the complementarity theorem
(Theorem~\ref{thm:quantum-complementarity-main}),
the scalar genus tower of $bc_\lambda$ is determined by the
$\beta\gamma_\lambda$ tower through
$\kappa(bc_\lambda)=-\kappa(\beta\gamma_\lambda)$. The class-C
shadow depth is also preserved, but not because every invariant is
obtained by changing the sign of~$\kappa$: the standard family has
the same quartic coefficient $S_4=-5/12$ and the same vanishing
range $S_r=0$ for $r\geq 5$.
The ordered-to-symmetric averaging step for Theorem~H is the
Hochschild averaging of
Corollary~\ref{cor:hochschild-averaging-symmetric}; it does not
identify bar cohomology growth with Hochschild amplitude.
Ghost-number rescaling is zero-mode gauge, and the spin line changes
the conformal vector rather than the curve-level chiral product.
At the bosonic string point $\lambda = 2$:
$\kappa(\mathrm{matter}) + \kappa(\mathrm{ghost}) = 13 + (-13) = 0$ \textup{(}free-field string sector\textup{)},
the anomaly cancellation of the $c = 26$ string.
\end{computation}

\subsubsection{Lattice VOA: five-theorem verification}

\begin{computation}[Five projections of $\Theta_{V_\Lambda}$;
\ClaimStatusProvedHere]
\label{comp:lattice-five-theorems}
\index{lattice VOA!five-theorem verification}
Shadow archetype: G (Gaussian, $r_{\max} = 2$) for the scalar
shadow of every even lattice VOA. Root vectors carry lattice-graded
simple-pole channels, but they are not class-L scalar shadows.
$\kappa(V_\Lambda) = \operatorname{rank}(\Lambda)$.
Central charge $c = \operatorname{rank}(\Lambda)$.
\begin{center}
\small
\renewcommand{\arraystretch}{1.3}
\begin{tabular}{lll}
\toprule
\emph{Theorem} & \emph{Projection of $\Theta_{V_\Lambda}$}
 & \emph{Content} \\
\midrule
A (bar-cobar) &
 $\Theta_{V_\Lambda}\big|_{\hbar=0} = \tau$ &
 $V_\Lambda^! = V_{\Lambda^\vee}$ on finite lattice sectors; form duality \\
B (inversion) &
 $\Omega(\barBgeom(V_\Lambda)) \simeq V_\Lambda$ &
 Genus-$0$ qi \\
C (complementarity) &
 scalar curvature of the Verdier form-dual oscillator lane &
 $\operatorname{rank}(\Lambda)+(-\operatorname{rank}(\Lambda))=0$ \\
D (modular char.) &
 $\Theta_{V_\Lambda}\big|_{g=1,n=0}
 = \operatorname{rank}(\Lambda) \cdot \lambda_1$ &
 $\kappa(V_\Lambda) = \operatorname{rank}(\Lambda)$ \\
H (chiral Hochschild) &
 Degree-preserving sub-MC on Koszul finite-type sectors &
 $\ChirHoch^n(V_\Lambda)=0$ for $n\notin\{0,1,2\}$ \\
\bottomrule
\end{tabular}
\end{center}
For every even lattice, the Heisenberg oscillator factor controls the
scalar shadow and gives $r_{\max}=2$. Root vectors
$e^\alpha$ may contribute lattice-graded simple-pole residues and Weyl
combinatorics; they do not create a cubic scalar obstruction term.
Self-dual lattices $\Lambda=\Lambda^\vee$ (e.g.\ the $E_8$ lattice,
the Leech lattice) are form self-dual. The Verdier scalar curvature
still changes sign in the Koszul package, so form self-duality is not
bar-cobar inversion.
\end{computation}


\subsection{Witten diagrams and Koszul duality}

\begin{technique}[Koszul residue pairing, not a Witten diagram]
Assume that $\cA$ has finite-dimensional weight spaces, a
nondegenerate BPZ/Verdier pairing, and a completed bar coalgebra whose
ternary transferred operation $m_3$ is defined. The algebraic
three-point Koszul residue pairing is
\[
K_{\mathrm{bar}}(a^!,b^!,c^!)
=
\left\langle a^!\otimes b^!\otimes c^!,\,m_3\right\rangle_{\mathrm{Verdier}},
\]
where $a^!,b^!,c^!\in\cA^!$ are line/defect-sector classes. A
holographic Witten diagram is obtained from this expression only after
one supplies a trace functor from $\cA^!$-modules to bulk fields, a
bulk-to-boundary propagator, and an equality between the Verdier
pairing and the physical BPZ pairing. These are extra data; they are
not consequences of the bar-cobar adjunction.
\end{technique}

\begin{example}[Conditional three-point comparison in \texorpdfstring{AdS$_3$}{AdS3}]
For operators $\mathcal{O}_i$ of dimension $\Delta_i$ in the boundary CFT:

\[\langle \mathcal{O}_1(z_1) \mathcal{O}_2(z_2) \mathcal{O}_3(z_3) \rangle = \frac{C_{123}}{|z_{12}|^{\Delta_1+\Delta_2-\Delta_3}|z_{23}|^{\Delta_2+\Delta_3-\Delta_1}|z_{31}|^{\Delta_3+\Delta_1-\Delta_2}}\]

If the holographic trace functor sends
$\mathcal{O}_i^!\in\cA^!$ to the bulk fields dual to
$\mathcal{O}_i$, and if the trace identifies the Verdier pairing with
the physical BPZ pairing, then
\[
C_{123}
=
\left\langle
\mathcal{O}_1^!\otimes\mathcal{O}_2^!\otimes\mathcal{O}_3^!,
m_3
\right\rangle_{\mathrm{Verdier}}.
\]
Without those two hypotheses this is only an algebraic residue
pairing in the Koszul-dual line/defect sector.
\end{example}


\subsection{Filtered and graded structures: compatibility}

\begin{definition}[Compatible filtration]
A filtration $F_\bullet\mathcal{A}$ on a graded chiral algebra $\mathcal{A} = \bigoplus_{n \in \mathbb{Z}} \mathcal{A}_n$ is \emph{compatible} if:
\begin{enumerate}
\item $F_p\mathcal{A} = \bigoplus_{n} F_p\mathcal{A}_n$ (respects grading)
\item $\mu(F_p\mathcal{A} \otimes F_q\mathcal{A}) \subset F_{p+q}\mathcal{A}$ (respects multiplication)
\item $\text{Gr}_p\mathcal{A} = F_p\mathcal{A}/F_{p-1}\mathcal{A}$ is graded
\item The associated graded $\text{Gr}\mathcal{A} = \bigoplus_p \text{Gr}_p\mathcal{A}$ is a chiral algebra
\end{enumerate}
\end{definition}

\begin{theorem}[Filtered bar complex; \ClaimStatusProvedHere]
\label{thm:filtered-bar-complex}
For a filtered chiral algebra $(F_\bullet\mathcal{A}, d)$, the bar complex inherits a compatible filtration:
\[F_p\barBgeom(\mathcal{A}) = \sum_{i_0 + \cdots + i_n \leq p} \Omega^*(\ConfigSpace{n+1}) \otimes F_{i_0}\mathcal{A} \otimes \cdots \otimes F_{i_n}\mathcal{A}\]
with:
\[\text{Gr}\barBgeom(\mathcal{A}) \cong \barBgeom(\text{Gr}\mathcal{A})\]
\end{theorem}

\begin{proof}
The differential preserves filtration:
\[d(F_p\barBgeom) \subset F_p\barBgeom\]
because:
\begin{itemize}
\item $d_{\text{int}}$ preserves filtration degree
\item $d_{\text{fact}}$ via residues: $\text{Res}_{D_{ij}}(F_{i_1} \otimes \cdots \otimes F_{i_n}) \subset F_{i_1 + \cdots + i_n}$
\item $d_{\text{config}}$ does not change filtration
\end{itemize}

The isomorphism $\text{Gr}\barBgeom(\mathcal{A}) \cong \barBgeom(\text{Gr}\mathcal{A})$ follows from:
\[\text{Gr}_p(F_{i_0}\mathcal{A} \otimes \cdots \otimes F_{i_n}\mathcal{A}) = \bigoplus_{j_0 + \cdots + j_n = p} \text{Gr}_{j_0}\mathcal{A} \otimes \cdots \otimes \text{Gr}_{j_n}\mathcal{A}\]
\end{proof}

\begin{definition}[Curved filtered algebra]
A \emph{curved filtered chiral algebra} is $(F_\bullet\mathcal{A}, d, m_0)$ where:
\begin{itemize}
\item $d: F_p\mathcal{A} \to F_p\mathcal{A}[1]$ (preserves filtration)
\item $m_0 \in F_0\mathcal{A}[2]$ (curvature in filtration degree 0)
\item $d^2 = [m_0, \cdot]$ (curved differential equation)
\end{itemize}
\end{definition}

\begin{theorem}[Curved Koszul duality \cite{Positselski11}; \ClaimStatusProvedElsewhere]
\label{thm:curved-koszul-duality}
For complete augmented curved filtered chiral algebras with
conilpotent curved bar coalgebra and finite-type associated graded:
\begin{enumerate}
\item The curved bar construction is a curved coalgebra with curvature $\kappa_C = \bar{m}_0$
\item The cobar of a curved coalgebra is a curved algebra
\item If $\text{Gr}\mathcal{A}$ is Koszul, then:
\[\Omegach(\barBgeom(\mathcal{A})) \simeq \mathcal{A}\]
as curved filtered algebras.
\end{enumerate}
The displayed equivalence is bar-cobar inversion in the curved
filtered category. The Koszul-dual algebra is obtained only after
passing to the bar cohomology coalgebra and applying the continuous
Verdier dual.
\end{theorem}


\subsection{Explicit low-degree computations}
\label{sec:explicit-low-degree}

The explicit low-degree bar complex computations for the free fermion, Heisenberg, and $\beta\gamma$ system are carried out in Theorem~\ref{thm:fermion-bar-complex-genus-0} (free fermion, \S\ref{sec:free-fermion}), Theorem~\ref{thm:heisenberg-bar} and Theorem~\ref{thm:heisenberg-bar-complete} (Heisenberg, \S\ref{sec:heisenberg-bar-complex}), and Theorem~\ref{thm:betagamma-bc-koszul} with Proposition~\ref{prop:bc-betagamma-orthogonality} ($\beta\gamma \leftrightarrow bc$, \S\ref{sec:fermion-boson-koszul}). All results match the predictions of Theorem~\ref{thm:bar-cobar-isomorphism-main}.


\subsection{Mass deformation and contractible bar complexes}
\label{subsec:massive-chiral-contractible}
\index{mass deformation!contractible bar complex}
\index{bar complex!contractible}

Consider two free chirals $X$, $Y$ with superpotential $W = mXY$,
$m \neq 0$. The derived critical locus is
\[
\bigl\{ \partial W/\partial X = mY = 0,\;
 \partial W/\partial Y = mX = 0 \bigr\} = \{0\},
\]
a single point: all modes are massive.

\begin{proposition}[Massive chirals have contractible bar complexes;
\ClaimStatusProvedElsewhere]
\label{prop:massive-chiral-contractible}
For the chiral algebra $\mathcal{A}_{X,Y,m}$ associated to the
superpotential $W = mXY$:
\begin{enumerate}[label=\textup{(\roman*)}]
\item The boundary algebra is $H^*(\mathcal{A}_{X,Y,m}) \cong \mathbb{C}$
 in cohomology (trivial).
\item The mass term introduces a differential $m_1 = Q_W$ on the bar
 complex $\barBgeom(\mathcal{A}_{X,Y,m})$ which kills all cohomology:
 \[
 H^*\bigl(\barBgeom(\mathcal{A}_{X,Y,m}),\, d + Q_W\bigr) \cong \mathbb{C}.
 \]
\item In particular, $\barBgeom(\mathcal{A}_{X,Y,m})$ is contractible:
 the bar complex is acyclic in all positive degrees.
\end{enumerate}
\end{proposition}

This is the simplest instance of the superpotential truncation
principle: of the full $A_\infty$ structure $(m_0, m_1, m_2, \ldots)$,
only $m_1 = Q_W$ survives in the massive regime, and it suffices to
trivialize the bar complex.

\begin{remark}[Mass gap and bar acyclicity]
\label{rem:mass-gap-bar-acyclicity}
\index{mass gap!bar acyclicity}
The physical content is that massive theories have contractible bar
complexes. The mass gap (the property that all excitations above the
vacuum have strictly positive energy) is detected algebraically by
positive-degree acyclicity of the perturbed reduced bar complex
\((\barBgeom(\mathcal A),d+Q_W)\). Conversely, a nonzero
positive-degree class in this perturbed complex is a massless
boundary mode. This statement concerns the massive deformation
complex; it does not identify the Koszul-dual algebra
\(\mathcal A^!\) with the boundary algebra, nor does it assert a
gap criterion for arbitrary interacting chiral theories.
\end{remark}

\begin{remark}[Boundary algebra as quotient by massive modes]
\label{rem:massive-boundary-quotient}
\index{boundary algebra!massive quotient}
\index{mass deformation!boundary algebra quotient}
The algebraic content of
Proposition~\ref{prop:massive-chiral-contractible} is that the mass
term $W = mXY$ introduces a differential that kills all but the
vacuum sector:
\begin{equation}\label{eq:massive-boundary-quotient}
 \cA_\partial^{W=mXY}
 \;=\;
 H^0\!\bigl(\cA_{X} \otimes \cA_{Y},\; Q_{W}\bigr)
 \;\cong\;
 \mathbb{C},
\end{equation}
where $Q_W = m\oint(Y\partial_X + X\partial_Y)$. The differential
$Q_W$ pairs the $X$-modes with $Y$-modes bijectively (since $m
\neq 0$), leaving only the vacuum. This is the $d = 2$ case of the
superpotential truncation (Construction~\ref{constr:superpotential-ainfty}):
the quadratic superpotential forces $m_1$ to pair all modes, with no
room for higher $A_\infty$ operations. In this quadratic model the
perturbed bar complex detects the mass gap:
\[
H^*\bigl(\bar{B}(\cA),d+Q_W\bigr) \cong \mathbb{C}.
\]
The mass deformation adds a quadratic perturbation to
$\Theta_{\cA}$ that renders it gauge-trivial in
$\operatorname{MC}(\mathfrak{g}^{\mathrm{amb}}_{\cA})$: the
perturbed shadow algebra $\cA^{\mathrm{sh}}_{W}$
(Definition~\ref{def:shadow-algebra}) collapses to
$\mathbb{C}$. The vanishing is a property of the deformed infrared
theory; it is not a redefinition of the ultraviolet free-field
constants $\kappa$, $\Delta$, or the contact classes.
\end{remark}

% ================================================================
% ČECH RESOLUTION AND HCA FOR FREE FIELDS
% ================================================================

\subsection{\v{C}ech homotopy chiral algebras for free fields}
\label{subsec:cech-hca-free-fields}
\index{homotopy chiral algebra!free fields}
\index{Cech complex@\v{C}ech complex!free field examples}

Resolutions on the Ran space are modeled by the bar complex; resolutions
on the Zariski site are modeled by the \v{C}ech complex (Theorem~\ref{thm:cech-hca}). On $\mathbb{P}^1$ with the standard two-element
cover, a decisive simplification occurs: \emph{every} \v{C}ech HCA
is strict, regardless of the chiral algebra. The reason is
purely combinatorial: the Moore complex $M\check{C}^\bullet$ is
concentrated in degrees~$0$ and~$1$, so there is no room for
a nontrivial Jacobiator. Non-strictness in the free-field
examples lives on the \emph{bar} side, and the \v{C}ech--bar
comparison map $\Phi^*$ is what creates it.

By Proposition~\ref{prop:two-element-strict}, for the standard
two-element cover $\mathcal{U} = \{U_0, U_\infty\}$ of
$\mathbb{P}^1$ and \emph{any} chiral algebra $\cA$, the Moore
complex $M\check{C}^\bullet$ is concentrated in degrees~$0$
and~$1$, so $F_n = 0$ for all $n \geq 3$ and the
Maurer--Cartan element is
\[
\Phi \;=\; (d_M,\; F_2,\; 0,\; 0,\; \ldots\,).
\]

\noindent
\v{C}ech strictness is a property of the
\emph{cover}, not the algebra: on a two-element cover of
$\mathbb{P}^1$, the Heisenberg, free fermion, $\beta\gamma$ system,
and affine Kac--Moody algebras \emph{all} have strict \v{C}ech HCA.
Non-strictness manifests on the bar side, through integration over
FM compactifications.

\subsubsection{Heisenberg on $\mathbb{P}^1$: mode computation}
\label{subsec:cech-heisenberg-p1}
\index{Heisenberg!Cech@\v{C}ech complex on $\mathbb{P}^1$}

Consider $\cH_k$ on $X = \mathbb{P}^1$ with cover
$U_0 = \operatorname{Spec}k[z]$,
$U_\infty = \operatorname{Spec}k[z^{-1}]$, and
$U_{0\infty} = \operatorname{Spec}k[z,z^{-1}]$.
The current $J(z) = \sum_{n \in \mathbb{Z}} J_n z^{-n-1}\,dz$
has mode expansions:
\begin{align*}
\cH_k(U_0) &= \bigoplus_{n \geq 0} k \cdot J_n
 \oplus k \cdot \mathbf{1}, \\
\cH_k(U_\infty) &= \bigoplus_{n \leq 0} k \cdot J_n
 \oplus k \cdot \mathbf{1}, \\
\cH_k(U_{0\infty}) &= \bigoplus_{n \in \mathbb{Z}} k \cdot J_n
 \oplus k \cdot \mathbf{1}.
\end{align*}
The Moore complex of the \v{C}ech cosimplicial object is
concentrated in degrees~$0$ and~$1$:
\[
(M\check{C})^0 = \cH_k(U_0),
\qquad
(M\check{C})^1 = \cH_k(U_{0\infty}) / \cH_k(U_\infty)
 \cong \bigoplus_{n > 0} k \cdot J_n.
\]
The \v{C}ech differential maps
$J_n \in (M\check{C})^0$ to its class in
$(M\check{C})^1$, which is zero for $n \leq 0$ and $J_n$ itself
for $n > 0$. Thus $H^0 = k \cdot \mathbf{1}$ (global sections)
and $H^1 = \bigoplus_{n < 0} k \cdot J_n$ (negative modes), in
agreement with the chiral homology of the Heisenberg
on~$\mathbb{P}^1$.

The \v{C}ech HCA Maurer--Cartan element is
$\Phi = (d_M, F_2, 0, 0, \ldots)$ by
Proposition~\ref{prop:two-element-strict}. The binary bracket
$F_2$ acts by the Heisenberg commutator:
$F_2(J_m, J_n) = km\,\delta_{m+n,0} \cdot \mathbf{1}$.

\begin{remark}[Both sides strict for the Heisenberg]
\label{rem:cech-bar-strict-heisenberg}
The Heisenberg is the unique case among standard families where
\emph{both} the \v{C}ech and bar sides are strict. The bar
$A_\infty$ structure satisfies $m_n = 0$ for $n \geq 3$
(\S\ref{sec:explicit-low-degree}), and the \v{C}ech HCA has
$F_n = 0$ for $n \geq 3$. The comparison map $\Phi^*$ of
Conjecture~\ref{conj:cech-bar-intertwining} sends the strict
\v{C}ech data to strict bar data: no higher operations are created
or destroyed.
\end{remark}

\subsubsection{Free fermion on $\mathbb{P}^1$: strict \v{C}ech HCA
with nontrivial cohomology}
\label{subsec:cech-fermion-p1}

For the free fermion $\psi(z)\psi(w) \sim 1/(z-w)$ on
$\mathbb{P}^1$, the \v{C}ech complex has the same two-element
structure with generators $\psi_r$ in place of $J_n$. The
difference is that $\psi$ is a section of $K_X^{1/2}$ (NS sector),
with half-integer mode expansion:
$\psi(z) = \sum_{r \in \mathbb{Z}+1/2} \psi_r\,z^{-r-1/2}$.
The \v{C}ech cohomology is
$H^0(\mathbb{P}^1, K_X^{1/2}) = 0$ (no global sections of
$\mathcal{O}(-1)$) and $H^1(\mathbb{P}^1, K_X^{1/2}) = k$ (by
Serre duality).

By Proposition~\ref{prop:two-element-strict}, the \v{C}ech HCA is
strict: $\Phi = (d_M, F_2, 0, 0, \ldots)$. As with the
Heisenberg, the bar side is also strict ($m_3 = 0$,
Theorem~\ref{thm:fermion-bar-complex-genus-0}), so $\Phi^*$
again maps strict to strict.

\subsubsection{$\beta\gamma$ on $\mathbb{P}^1$: strict \v{C}ech,
non-strict bar}
\label{subsec:cech-betagamma-p1}

For the $\beta\gamma$ system on $\mathbb{P}^1$ with the two-element
cover, the \v{C}ech complex has two generators:
$\beta \in \Gamma(X, K_X^\lambda)$ and
$\gamma \in \Gamma(X, K_X^{1-\lambda})$. The binary bracket $F_2$ is
induced by $\beta(z)\gamma(w) \sim 1/(z-w)$; in modes,
$F_2(\beta_m, \gamma_n) = \delta_{m+n,0} \cdot \mathbf{1}$.

By Proposition~\ref{prop:two-element-strict}, the \v{C}ech HCA is
again strict:
\[
\Phi_{\beta\gamma}
 \;=\; (d_M,\; F_2,\; 0,\; 0,\; \ldots\,).
\]
The normally ordered product $:\!\beta\gamma\!:$ does give rise to
nontrivial iterated brackets, but these would produce a
Jacobiator in $M\check{C}^{\geq 2}$, which vanishes for the
two-element cover. There is no \v{C}ech degree in which
a nontrivial $F_3$ could live.

The non-strictness of the $\beta\gamma$ system appears on the
\emph{bar} side: the transferred $A_\infty$ operation $m_3 \neq 0$
arises from integration over the Fulton--MacPherson
compactification $\overline{\mathrm{FM}}_3(\mathbb{C})$, where
the triple collision boundary gives a nonvanishing contribution.
The comparison map $\Phi^*$ of
Conjecture~\ref{conj:cech-bar-intertwining} sends the strict
\v{C}ech data $F_3 = 0$ to the non-strict bar data $m_3 \neq 0$:
the bar operation $m_3$ is \emph{created} by the comparison, not
inherited from the \v{C}ech side.

\subsubsection{\v{C}ech--bar comparison for free fields}
\label{subsec:cech-bar-free-field-comparison}
\index{Cech--bar comparison@\v{C}ech--bar comparison!free fields}

The three examples exhibit a clean trichotomy:

\begin{center}
\renewcommand{\arraystretch}{1.3}
\begin{tabular}{lccc}
\textbf{Chiral algebra}
 & \textbf{\v{C}ech HCA}
 & \textbf{Bar $A_\infty$}
 & \textbf{$\Phi^*$ creates $m_3$?} \\
\hline
Heisenberg $\cH_k$
 & strict ($F_3 = 0$)
 & strict ($m_3 = 0$)
 & no \\
Free fermion $\mathcal{F}$
 & strict ($F_3 = 0$)
 & strict ($m_3 = 0$)
 & no \\
$\beta\gamma$
 & strict ($F_3 = 0$)
 & non-strict ($m_3 \neq 0$)
 & yes \\
\end{tabular}
\end{center}

\noindent
The same pattern holds for the affine Kac--Moody algebras
$\widehat{\mathfrak{g}}_k$ on the two-element cover: the \v{C}ech
HCA is strict, but the bar side has $m_3 \neq 0$ from the cubic
Lie bracket. In every case, $\Phi^*$ is the mechanism that
produces bar non-strictness from \v{C}ech strictness.

\begin{remark}[\v{C}ech HCA strictness depends on the cover]
\label{rem:hca-strictness-diagnostic}
\index{homotopy chiral algebra!strictness diagnostic}
The diagnostic for \v{C}ech HCA strictness is the \emph{cover},
not the algebra. For a two-element cover $\mathcal{U} =
\{U_0, U_\infty\}$ of $\mathbb{P}^1$, the Moore complex
$M\check{C}^\bullet$ is concentrated in degrees $0$ and $1$, and
$\check{C}^p = 0$ for $p \geq 2$ leaves no room for a
Jacobiator (Proposition~\textup{\ref{prop:two-element-strict}}).
Every chiral algebra on this cover has strict \v{C}ech HCA. On a
finer cover with triple intersections, a nonzero Jacobiator produces
a genuine homotopy $F_3$; strictness then depends on the algebra.
Equivalently,
\emph{non-strictness is a bar-side phenomenon on $\mathbb{P}^1$,
arising from FM compactification, not from \v{C}ech descent}.
\end{remark}

\section{Mukai--Heisenberg boundary factors}
\label{sec:ff-supplement}

\begin{remark}[Mukai-Heisenberg as free-field output of $\Phi_2$ on K3]
\label{rem:ff-mukai-heisenberg}
At $\Phi_2(D^b\mathrm{Coh}(K3))$ the free-field boundary is the
rank-$24$ Mukai-Heisenberg oscillator factor attached to
$\Lambda_{\mathrm{Muk}}=\mathrm{II}_{4,20}$. The oscillator
denominator is $\eta(q)^{-24}$, while the determinant/Euler product is
$\eta(q)^{24}$; these are inverse normalisations of the same Gaussian
factor, not a reconstruction of the Borcherds denominator. The
Hodge-supertrace invariant is
$\kappa_{\mathrm{cat}}(K3)=\chi(\mathcal{O}_{K3})=2$, while the
Theorem~C $\mathcal{B}$-row uses the Mukai positive-signature scalar
$\kappa_{\mathrm{ch}}^{\mathrm{Mukai}}=c_+(\Lambda_{\mathrm{Muk}})=4$.
The Koszul conductor on the Mukai-doubling side is
$K^{\kappa_{\mathrm{ch}}^{\mathrm{Mukai}}}
=2\,c_+(\mathrm{Mukai}(K3))=8$. The Theorem~C bucket enlarges on the
Lorentzian-lattice-parametric $\mathcal{B}$-family to
$\{0, 8, 13, 250/3, 98/3\}$.
\end{remark}

\begin{remark}[Genus-2 Siegel character factorisation]
\label{rem:ff-genus2-factorisation}
At the separating-node degeneration of the genus-$2$ Siegel character
of the K3 chiral bialgebra $\mathbf H_{\Delta_5}$ (Vol~III
Chapter~\ref{V3-ch:k3-chiral-bialgebra-platonic}), the unnormalised
Igusa square $\Phi_{10}^{\mathrm{un}}=\Delta_5^2$ satisfies
\[
\frac{1}{\Phi_{10}^{\mathrm{un}}(\mathrm{diag}(\tau,\rho))}
=
\frac{1}{\eta(\tau)^{24}}\cdot\frac{1}{\phi_{10,1}(\rho)}.
\]
The inverse eta factor $1/\eta(\tau)^{24}$ reappears on the elliptic
$\tau$-direction, multiplied by the Jacobi factor on the
$\rho$-direction. This is a boundary factorisation of the
normalised Siegel character. The DMVV product and the primitive
Gritsenko--Nikulin denominator $\Delta_5$ supply the automorphic
input; the class-$\mathsf G$ Heisenberg boundary supplies only the
Gaussian $\eta$ factor.
\end{remark}
